\section{Computational Interpretation, Enigma Analogy, and Bomba-Style Structural Search}
\begin{quote}\small
\noindent\textbf{Status.} \emph{Active compressed auxiliary layer, subordinate to the mathematics-first core.}\par
\noindent\textbf{Block function.} This chapter packages the computational, machine, and decoder reading into a short theorem-usable interface, so later arguments can cite the route/decoder layer without reopening the full exploratory genealogy.
\end{quote}
\label{sec:computational-reading}
\addcontentsline{toc}{section}{Computational Interpretation, Enigma Analogy, and Bomba-Style Structural Search}

\paragraph{Status note.}
For the current v\docversion\ release line, the computational/enigma/millennium-language material is now read through a compressed interface. The purpose of this chapter is no longer to carry a long exploratory ladder on every pass, but to preserve a short theorem-usable computational reading that remains subordinate to the mathematics-first architecture.

\subsection*{1--19J. Compressed computational interpretation surface}
The historical computational line can now be cited through a single bundled surface. Its content is unchanged in substance: the chapter still records (i) the historical genealogy from Turing, Zuse, and algorithmic state-transition thinking; (ii) the reason an Enigma-/reflector-/differencing-style comparison becomes structurally natural in the present routed architecture; (iii) the millennium-machine reading with collapse states, projector logic, HC-selection, and reversible return structure; (iv) the provisional language stack MML $\to$ MMA-0.1 $\to$ MCP together with certificate/export discipline; and (v) the Bomba/Bombe-style search idea, not as a literal cryptanalytic identity claim but as a constrained elimination picture for admissible structural menus.

The primitive computational reserve remains represented by the normalized instruction family
\[
\begin{aligned}
\{&\texttt{COLLAPSE},\texttt{PROJECT},\texttt{CLOSE},\texttt{PROBE},\texttt{TRACE},\texttt{COMPARE},\\
&\texttt{DIFF},\texttt{HCSEL},\texttt{HOPF},\texttt{REJECT},\texttt{EXPORT\_BOOL},\texttt{FORWARD}\}
\end{aligned}
\]
and the intended proof-compression reading still runs through the compilation chain
\[
\begin{aligned}
&\text{MML source}
\Longrightarrow \text{MMA-0.1 instruction stream}\\
&\Longrightarrow \mathcal P_{\mathrm{MCP}}
\Longrightarrow \text{Boolean-compatible host state}.
\end{aligned}
\]
What changes in the active line is only the packaging: the long exploratory ladder is replaced by a short citation surface.

\begin{definition}[Compressed computational interpretation surface]
Write
\[
\mathfrak C^{\mathrm{comp}}
=
\langle
\mathcal H_{\mathrm{hist}},
\mathcal E_{\mathrm{route}},
\mathcal B_{\mathrm{elim}},
\mathfrak R^{\mathrm{comp}}
\rangle
\]
for the compressed computational interpretation surface, where:
\begin{itemize}[leftmargin=2em]
    \item $\mathcal H_{\mathrm{hist}}$ denotes the historical genealogy and machine-reading motivation;
    \item $\mathcal E_{\mathrm{route}}$ denotes the reflector-/projector-/routing-style structural comparison reading;
    \item $\mathcal B_{\mathrm{elim}}$ denotes the constrained Bomba/Bombe-style elimination viewpoint on admissible local menus;
    \item $\mathfrak R^{\mathrm{comp}}$ denotes the compressed computational reserve surface already carrying the instruction, macro, proof-template, and certificate workflow layers.
\end{itemize}
\end{definition}

\begin{proposition}[Exact citation reading of the compressed computational interpretation surface]
The compressed computational interpretation surface $\mathfrak C^{\mathrm{comp}}$ is citation-equivalent to the expanded ladder previously distributed across the historical computational line, the Enigma/reflector comparison, the Boolean/projector layer, the millennium-machine viewpoint, the MML/MMA/MCP language discussion, the computational reserve examples, and the Bomba/Bombe search outlook. Consequently later arguments may cite $\mathfrak C^{\mathrm{comp}}$ whenever they use only the existence of that computational reading and do not require reopening the internal historical examples or the micro-expansion of the reserve stack.
\end{proposition}

\begin{corollary}[What the compressed computational interpretation surface already buys]
The active line keeps a usable computational reading without carrying the full exploratory staircase. Later proof-compression, host-interface, or search-orientation arguments may therefore appeal to a single computational interpretation surface instead of reopening the whole Turing/Zuse/Enigma/MML/MCP ladder each time.
\end{corollary}


\subsection*{19K. Local operational semantics and conditional route to universality}
The compressed computational interpretation surface can be sharpened one step further without overclaiming. The correct present reading is not that universality has already been proved, but that the current release line now supports a \emph{conditional operational semantics} together with a \emph{conditional route} by which a later full universality theorem could be closed if the remaining component-identification obligations are discharged.

\begin{definition}[Local operational semantics for the compressed computational surface]
Fix a finite local machine state
\[
\Sigma=(\tau,h,q,\rho,\chi),
\]
where $\tau$ denotes a finite tape snapshot, $h$ a head position, $q$ a finite control state, $\rho$ a routed/readout state, and $\chi$ the current shell-/closure certificate. Interpret the primitive instruction family
\[
\begin{aligned}
\{&\texttt{COLLAPSE},\texttt{PROJECT},\texttt{CLOSE},\texttt{PROBE},\texttt{TRACE},\texttt{COMPARE},\\
&\texttt{DIFF},\texttt{HCSEL},\texttt{HOPF},\texttt{REJECT},\texttt{EXPORT\_BOOL},\texttt{FORWARD}\}
\end{aligned}
\]
as generating a partial local transition relation
\[
\Sigma \xrightarrow{\,\iota\,} \Sigma'
\]
on shell-admissible states, with the following intended roles:
\begin{itemize}[leftmargin=2em]
  \item \texttt{PROJECT}, \texttt{COMPARE}, \texttt{DIFF}, and \texttt{EXPORT\_BOOL} produce the locally readable branch predicates;
  \item \texttt{FORWARD} and \texttt{HOPF} move the effective read/write location along the admissible routed support;
  \item \texttt{COLLAPSE}, \texttt{CLOSE}, and \texttt{HCSEL} update the local closure/control state while preserving admissibility;
  \item \texttt{PROBE} and \texttt{TRACE} realize the controlled inspection/readout layer;
  \item \texttt{REJECT} records a halted or discarded local branch.
\end{itemize}
The present Companion only uses this as a local operational semantics on admissible windows; it is not yet a global universality theorem.
\end{definition}

\begin{definition}[Spinorial tape hypothesis and conditional Z3-route hypothesis]
For the current release line, write $\mathbf{STH}$ for the \emph{spinorial tape hypothesis}: the optimal line carried by the admissible spinorial transport acts as an effectively extendable read/write support for the local state component $\tau$, so that finite tape snapshots may be prolonged along that support without breaking shell-admissibility.

Write $\mathbf{Z3H}$ for the accompanying \emph{conditional Z3-route hypothesis}: once conditional branching, locally persistent read/write support, and effective looping/return are available on that spinorial support, the present MML$\to$MMA-0.1$\to$MCP line can be organized so as to simulate a standard sequential control core in the same sense in which a Z3-style line becomes computationally universal only after the relevant strengthening assumptions are supplied.
\end{definition}

\begin{proposition}[Conditional route from the compressed computational surface to universality]
Assume $\mathbf{STH}$ and $\mathbf{Z3H}$, and assume in addition the \emph{component identification hypothesis}: the current structural components really do instantiate the required computational roles, namely
\begin{itemize}[leftmargin=2em]
  \item a persistent effective tape support carried by the spinorial optimal line,
  \item a finite control/update layer carried by the coded MML/MMA-0.1/MCP stack,
  \item a locally decidable branch/readout layer carried by the projector/Boolean comparison interface,
  \item and a return/loop discipline carried by the routed HC-selected closure structure.
\end{itemize}
Then the compressed computational interpretation surface $\mathfrak C^{\mathrm{comp}}$ supports a \emph{conditional route to universality}: to prove a true universality theorem it is sufficient to prove that these identified components realize a standard reference machine core (for example a register-machine or tape-machine core) under the local transition semantics above.
\end{proposition}

\begin{corollary}[What the conditional universality route already buys]
The current v\docversion\ line may now state the computational claim in the correct sharpened form: the Companion does \emph{not} yet prove Turing-completeness, but it does isolate a specific route by which such a result would follow. The remaining hard step is no longer ``invent a machine reading from scratch,'' but ``prove that the present structural components are indeed the tape, control, branch, and return components required by the conditional route.''
\end{corollary}


\subsection*{19L. Component role-hardening and a conditional Zuse-style primitive core}
The historical calibration from Z1 through Z3 suggests a useful discipline for the present line. The architectural idea already existed early, but reliability improved when storage, control, and arithmetic roles became cleaner and more stable across the machine family. In the same spirit, the present Companion should not jump immediately from a compressed computational reading to a full universality claim. The next correct step is to harden the four component roles separately and only then close stronger machine-theoretic statements.

\begin{definition}[Local role certificate package]
Write
\[
\mathbf{Role}=(\mathbf T,\mathbf Q,\mathbf B,\mathbf L)
\]
for a local role certificate package on an admissible computational window, where:
\begin{itemize}[leftmargin=2em]
  \item $\mathbf T$ certifies an effectively extendable sequential read/write support;
  \item $\mathbf Q$ certifies a finite control/update layer with stable local instruction semantics;
  \item $\mathbf B$ certifies a locally decidable branch predicate layer;
  \item $\mathbf L$ certifies a loop/return discipline capable of resuming or repeating the local control flow without leaving admissibility.
\end{itemize}
A local machine reading is called \emph{Zuse-primitive} on the given window if it carries such a package, even if no full universality theorem has yet been proved.
\end{definition}

\begin{proposition}[Spinorial optimal line as conditional tape role]
Assume $\mathbf{STH}$. Then the admissible spinorial optimal line carries the conditional tape certificate $\mathbf T$ in the following weak but machine-relevant sense: finite local snapshots may be written, read, and prolonged along that line by admissible transport, and the transition primitives \texttt{FORWARD} and \texttt{HOPF} provide the local sequential displacement needed for an effective head movement on admissible windows.
\end{proposition}

\begin{proposition}[Coded MML/MMA/MCP stack as conditional control role]
Assume the component identification hypothesis for the coded MML$\to$MMA-0.1$\to$MCP line. Then the coded stack carries the conditional control certificate $\mathbf Q$: it provides a finite control/update discipline which selects, routes, and stabilizes the local instruction semantics on admissible computational windows.
\end{proposition}

\begin{proposition}[Projector/Boolean layer as conditional branch role]
Assume the locally decidable projector/readout and Boolean export layer already isolated in the compressed computational surface. Then the current release line carries the conditional branch certificate $\mathbf B$: branch predicates can be extracted locally in a form sufficient to distinguish continuation, rejection, and comparison-sensitive update on admissible windows.
\end{proposition}

\begin{proposition}[Routed HC-selected closure as conditional loop/return role]
Assume the routed HC-selected closure structure persists under admissible local updates. Then the current release line carries the conditional loop/return certificate $\mathbf L$: local computational segments can be closed, resumed, or re-entered through the routed closure discipline without immediately destroying admissibility.
\end{proposition}

\begin{theorem}[Conditional Zuse-core closure theorem]
Assume $\mathbf{STH}$, the local operational semantics on admissible computational windows, the component identification hypothesis, and the conditional role certificates $\mathbf{Role}=(\mathbf T,\mathbf Q,\mathbf B,\mathbf L)$ above. Then the compressed computational interpretation surface $\mathfrak C^{\mathrm{comp}}$ closes to a \emph{conditional local Zuse-core} on every admissible certified window: there is a well-defined primitive sequential machine reading with
\begin{itemize}[leftmargin=2em]
  \item a persistent sequential storage/support component carried by the spinorial optimal line,
  \item a finite control/update component carried by the coded MML$\to$MMA-0.1$\to$MCP stack,
  \item a locally decidable branch/readout component carried by the projector/Boolean comparison interface,
  \item and a routed loop/return component carried by the HC-selected closure structure.
\end{itemize}
Moreover this local Zuse-core is closed under admissible one-step execution and under finite concatenation of locally certified segments, so the present line already supports a genuine \emph{primitive Zuse-style computational apparatus} in the conditional sense of this chapter. What remains open is not the existence of that primitive core, but the proof that these components are globally stable, extendable, and strong enough to simulate a standard reference machine.
\end{theorem}

\begin{corollary}[Conditional Zuse-style primitive computational core]
Assume $\mathbf{STH}$, the component identification hypothesis, and the conditional role certificates $\mathbf{Role}=(\mathbf T,\mathbf Q,\mathbf B,\mathbf L)$ above. Then the compressed computational interpretation surface $\mathfrak C^{\mathrm{comp}}$ already supports a \emph{conditional Zuse-style primitive computational core}: the current release line carries a genuine sequential storage/support reading, a finite control/update layer, a local branch/readout layer, and a routed loop/return discipline. What is still missing for stronger claims is not the existence of a primitive machine reading, but the final proof that these roles are globally stable and strong enough for an explicit reference-machine simulation.
\end{corollary}

\begin{corollary}[What the conditional Zuse-core closure theorem already buys]
Later arguments may now cite a single bundled closure statement for the machine-reading side of the Companion: under the certified local role package, the present line is already closed as a primitive Zuse-style computational core on admissible windows. Thus one no longer needs to reopen the tape/control/branch/loop roles separately unless the proof requires strengthening one of the four certificates themselves.
\end{corollary}


\subsection*{19M. Reduced proof-obligation package for the conditional Zuse-core}
Once the conditional Zuse-style primitive core has been isolated, the remaining machine-theoretic burden should be stated as a small explicit obligation package rather than left as a diffuse list of promises. The right next reading is therefore not ``prove universality all at once,'' but ``reduce the open universality route to a finite role-upgrade package.''

\begin{definition}[Reduced universality-obligation package for the conditional Zuse-core]
Write
\[
\mathbf U=(\mathbf U_T,\mathbf U_Q,\mathbf U_B,\mathbf U_L,\mathbf U_S)
\]
for the reduced universality-obligation package attached to the conditional Zuse-core, where:
\begin{itemize}[leftmargin=2em]
  \item $\mathbf U_T$ is the obligation that the spinorial optimal line is not merely locally readable, but globally prolongable as an effective read/write support;
  \item $\mathbf U_Q$ is the obligation that the coded MML$\to$MMA-0.1$\to$MCP control layer composes stably across admissible local windows;
  \item $\mathbf U_B$ is the obligation that the projector/Boolean layer yields executable and persistently decidable branch predicates under admissible concatenation;
  \item $\mathbf U_L$ is the obligation that the routed HC-selected closure discipline provides genuine resumable loop/return control under repeated local execution;
  \item $\mathbf U_S$ is the obligation that these upgraded components explicitly simulate a standard sequential reference-machine core on admissible windows.
\end{itemize}
A conditional Zuse-core is called \emph{universality-ready} if all five obligations in $\mathbf U$ are discharged.
\end{definition}

\begin{proposition}[Proof-obligation reduction for the conditional Zuse-core]
Assume the conditional Zuse-core closure theorem. Then the remaining gap between the present release line and a genuine universality theorem is reduced to the package $\mathbf U=(\mathbf U_T,\mathbf U_Q,\mathbf U_B,\mathbf U_L,\mathbf U_S)$ above. In particular, one no longer needs to reopen the existence of a primitive machine reading itself; the open task is only to upgrade the four certified roles from local conditional availability to globally stable executable roles and then prove the explicit reference-machine simulation.
\end{proposition}

\begin{corollary}[What the reduced universality-obligation package already buys]
The machine-theoretic remainder of the Companion is now sharply localized. Later revisions may therefore proceed componentwise: tape hardening, control hardening, branch hardening, loop/return hardening, and finally explicit simulation. This is a much narrower and more testable program than the earlier diffuse slogan ``show that the whole architecture is universal.''
\end{corollary}

\begin{remark}[What the Z1$\to$Z2$\to$Z3 improvement path suggests for the present model]
The historical machine line suggests three concrete lessons for the current architecture.
\begin{itemize}[leftmargin=2em]
  \item First, one should keep the candidate storage support as stable as possible while improving the control/branch implementation; this favors treating the spinorial optimal line as the persistent support and hardening the projector/control stack around it.
  \item Second, the decisive upgrade path may come from improving reliability of the control and branching layers before demanding full universality; this mirrors the move from more fragile early realization to a more robust relay-style execution discipline.
  \item Third, the right immediate target is not yet ``prove Turing-completeness,'' but ``prove that the present line already forms a Zuse-primitive machine in the precise conditional sense above.''
\end{itemize}
\end{remark}




\subsection*{19N. First hardening priority for the reduced universality-obligation package}
Once the remaining universality route has been reduced to a finite package, the next question is no longer merely ``what is missing,'' but ``in which order should the missing pieces be hardened so that later upgrades do not have to be reopened unnecessarily.'' The Zuse-guided reading suggests that the first useful order is not to attack full simulation immediately, but to stabilize the executable control layers around the candidate storage support before demanding the final global machine theorem.

\begin{definition}[First Zuse-guided hardening priority order for the reduced universality-obligation package]
Write
\[
\mathbf H_1=(\mathbf U_Q,\mathbf U_B,\mathbf U_L,\mathbf U_T,\mathbf U_S)
\]
for the first Zuse-guided hardening priority order attached to the reduced universality-obligation package. Its intended reading is the following:
\begin{itemize}[leftmargin=2em]
  \item first harden the finite control/update stack so that local instruction semantics compose stably across admissible windows;
  \item then harden the projector/Boolean branch layer so that executable branching survives admissible concatenation;
  \item then harden the routed HC-selected loop/return discipline so that repeated local execution remains resumable;
  \item only afterwards demand that the spinorial optimal line be upgraded from local readable support to globally prolongable effective tape;
  \item and only at the end require the explicit reference-machine simulation.
\end{itemize}
\end{definition}

\begin{proposition}[Zuse-guided priority reduction for the reduced universality-obligation package]
Assume the conditional Zuse-core closure theorem and the reduced universality-obligation package $\mathbf U=(\mathbf U_T,\mathbf U_Q,\mathbf U_B,\mathbf U_L,\mathbf U_S)$. Then the first useful hardening route may be read through the priority order $\mathbf H_1=(\mathbf U_Q,\mathbf U_B,\mathbf U_L,\mathbf U_T,\mathbf U_S)$ above: the present line already has a persistent storage candidate, but its most leverage-sensitive uncertified parts lie first in the execution stack surrounding that candidate. In particular, upgrading control, then branch, then loop/return yields the cleanest path to a later tape hardening and finally to explicit simulation, because it stabilizes the executable local dynamics before asking for the strongest global machine claim.
\end{proposition}

\begin{corollary}[What the first Zuse-guided hardening priority order already buys]
The open universality route is now not only finite but ordered. Later revisions therefore need not ask merely whether the current release line is already universal; they may instead proceed along a narrower program: first make the control layer compose, then make branching executable, then make local return resumable, then globalize the tape, and only afterwards prove explicit simulation. This turns the remaining machine-theoretic work from one undifferentiated challenge into a staged hardening schedule.
\end{corollary}



\subsection*{19O. First local hardening of the control obligation $\mathbf U_Q$}
The priority order above says that the first leverage-sensitive remaining task is no longer the tape itself but the stability of the executable control stack around that tape candidate. The natural first hardening step is therefore to isolate a local criterion under which the instruction/update layer is not merely available pointwise, but composes across adjacent admissible windows without reopening its semantics at every cut.

\begin{definition}[First local control-hardening certificate for $\mathbf U_Q$]
Write
\[
\mathbf Q_{\mathrm{loc}}(I)=1
\]
for the first local control-hardening certificate on a finite admissible window $I$ if the following hold simultaneously on $I$:
\begin{itemize}[leftmargin=2em]
  \item the local operational semantics of the compressed computational surface assigns a well-defined successor state to every admissible instruction step on $I$;
  \item adjacent admissible subwindows of $I$ compose without changing the decoded control meaning of the instruction stack;
  \item the induced controller action remains compatible with the later-depth cycle-control layer whenever that layer is active on the same window.
\end{itemize}
In words, $\mathbf Q_{\mathrm{loc}}(I)=1$ means that the present line already carries a locally executable and composition-stable control role on the window $I$.
\end{definition}

\begin{proposition}[First reduction of the control obligation $\mathbf U_Q$]
Assume the conditional Zuse-core closure theorem, the reduced universality-obligation package \(\mathbf U=(\mathbf U_T,\mathbf U_Q,\mathbf U_B,\mathbf U_L,\mathbf U_S)\), and the first Zuse-guided hardening priority order. Suppose moreover that there exists an admissible cover by finite windows on which the local control-hardening certificate \(\mathbf Q_{\mathrm{loc}}=1\) holds and is stable under admissible concatenation. Then the obligation \(\mathbf U_Q\) is discharged at the first hardening stage, and the remaining route reduces to
\[
(\mathbf U_B,\mathbf U_L,\mathbf U_T,\mathbf U_S).
\]
Equivalently, once the instruction/update layer composes locally and survives admissible concatenation without semantic drift, the first unresolved machine-theoretic gap no longer lies in control but in branch, loop/return, tape globalization, and explicit simulation.
\end{proposition}

\begin{corollary}[What the first local hardening of $\mathbf U_Q$ already buys]
The first hardening stage is now no longer just an ordering slogan. It has been turned into a concrete target: prove local control compositionality on admissible windows and its persistence under admissible concatenation. If that certificate is obtained, later revisions need not reopen the instruction stack itself and may move directly to the branch obligation as the next active gap.
\end{corollary}


\subsection*{19P. First local hardening of the branch obligation $\mathbf U_B$}
Once the control layer composes locally, the next leverage-sensitive open task is to certify that the projector/Boolean layer does not merely look branch-like, but already induces an executable local branch role on admissible windows. The natural next hardening step is therefore to isolate a local criterion under which branch alternatives are decoded, selected, and propagated without collapsing the admissible control meaning established at the previous stage.

\begin{definition}[First local branch-hardening certificate for $\mathbf U_B$]
Write
\[
\mathbf B_{\mathrm{loc}}(I)=1
\]
for the first local branch-hardening certificate on a finite admissible window $I$ if the following hold simultaneously on $I$:
\begin{itemize}[leftmargin=2em]
  \item the projector/Boolean layer assigns a well-defined admissible branch menu to every locally admissible branch site on $I$;
  \item whenever a branch discriminator is active, exactly one admissible local branch successor is selected at the level of the compressed computational surface;
  \item the selected local branch successor remains compatible with the already hardened local control meaning and with the later-depth cycle-control layer whenever that layer is active on the same window;
  \item admissible concatenation of adjacent certified subwindows does not change the decoded local branch choice except at explicitly reopened branch sites.
\end{itemize}
In words, $\mathbf B_{\mathrm{loc}}(I)=1$ means that the present line already carries a locally executable and concatenation-stable branch role on the window $I$.
\end{definition}

\begin{proposition}[First reduction of the branch obligation $\mathbf U_B$]
Assume the conditional Zuse-core closure theorem, the reduced universality-obligation package \(\mathbf U=(\mathbf U_T,\mathbf U_Q,\mathbf U_B,\mathbf U_L,\mathbf U_S)\), the first Zuse-guided hardening priority order, and that the local control-hardening certificate has already discharged \(\mathbf U_Q\). Suppose moreover that there exists an admissible cover by finite windows on which the local branch-hardening certificate \(\mathbf B_{\mathrm{loc}}=1\) holds and is stable under admissible concatenation away from explicitly reopened branch sites. Then the obligation \(\mathbf U_B\) is discharged at the second hardening stage, and the remaining route reduces to
\[
(\mathbf U_L,\mathbf U_T,\mathbf U_S).
\]
Equivalently, once the projector/Boolean layer yields an executable locally unique branch choice which survives admissible concatenation and remains compatible with the control stack, the next unresolved machine-theoretic gap no longer lies in branch but in loop/return, tape globalization, and explicit simulation.
\end{proposition}

\begin{corollary}[What the first local hardening of $\mathbf U_B$ already buys]
The second hardening stage is now also no longer only a priority promise. It has become a concrete target: certify local executable branch selection on admissible windows and its persistence under admissible concatenation, except where the theory explicitly reopens the branch site. If that certificate is obtained, later revisions may treat branch as locally settled and move directly to the loop/return obligation as the next active gap.
\end{corollary}


\subsection*{19Q. First local hardening of the loop/return obligation $\mathbf U_L$}
Once local control and local branch have been hardened, the next leverage-sensitive open task is to certify that routed HC-selected closure is not merely closure-like, but already induces a locally executable loop/return role on admissible windows. The natural third hardening step is therefore to isolate a local criterion under which admissible return sites are reopened, routed, and closed again without losing the local control-branch meaning established at the previous stages.

\begin{definition}[First local loop/return-hardening certificate for $\mathbf U_L$]
Write
\[
\mathbf L_{\mathrm{loc}}(I)=1
\]
for the first local loop/return-hardening certificate on a finite admissible window $I$ if the following hold simultaneously on $I$:
\begin{itemize}[leftmargin=2em]
  \item every admissible reopened return site on $I$ carries a well-defined routed HC-selected closure menu;
  \item whenever a local return event is triggered, exactly one admissible routed return/closure successor is selected at the level of the compressed computational surface;
  \item the selected routed return/closure successor remains compatible with the already hardened local control and local branch meanings on the same window;
  \item admissible concatenation of adjacent certified subwindows preserves the decoded routed return/closure choice except at explicitly reopened return sites.
\end{itemize}
In words, $\mathbf L_{\mathrm{loc}}(I)=1$ means that the present line already carries a locally executable and concatenation-stable loop/return role on the window $I$.
\end{definition}

\begin{proposition}[First reduction of the loop/return obligation $\mathbf U_L$]
Assume the conditional Zuse-core closure theorem, the reduced universality-obligation package \(\mathbf U=(\mathbf U_T,\mathbf U_Q,\mathbf U_B,\mathbf U_L,\mathbf U_S)\), the first Zuse-guided hardening priority order, and that the local control-hardening and local branch-hardening certificates have already discharged \(\mathbf U_Q\) and \(\mathbf U_B\). Suppose moreover that there exists an admissible cover by finite windows on which the local loop/return-hardening certificate \(\mathbf L_{\mathrm{loc}}=1\) holds and is stable under admissible concatenation away from explicitly reopened return sites. Then the obligation \(\mathbf U_L\) is discharged at the third hardening stage, and the remaining route reduces to
\[
(\mathbf U_T,\mathbf U_S).
\]
Equivalently, once routed HC-selected closure yields an executable locally unique return/closure choice which survives admissible concatenation and remains compatible with the hardened control and branch layers, the next unresolved machine-theoretic gap no longer lies in loop/return but in tape globalization and explicit reference-machine simulation.
\end{proposition}

\begin{corollary}[What the first local hardening of $\mathbf U_L$ already buys]
The third hardening stage is now also no longer only a priority promise. It has become a concrete target: certify local executable routed return/closure selection on admissible windows and its persistence under admissible concatenation, except where the theory explicitly reopens the return site. If that certificate is obtained, later revisions may treat loop/return as locally settled and move directly to the tape-globalization obligation as the next active gap.
\end{corollary}


\subsection*{19R. First globalization of the tape obligation $\mathbf U_T$}
Once local control, local branch, and local loop/return have been hardened, the remaining machine-theoretic gap before explicit reference-machine simulation is no longer a local choice problem but a persistence problem: the spinorial optimal line must be shown to support a globally extensible read/write track across admissibly concatenated windows. The natural fourth hardening step is therefore to isolate a globalization criterion under which the local tape role ceases to be merely windowwise heuristic and becomes a coherent tape support along an admissible chain.

\begin{definition}[First tape-globalization certificate for $\mathbf U_T$]
Write
\[
\mathbf T_{\mathrm{glob}}(\Gamma)=1
\]
for the first tape-globalization certificate on an admissible chain $\Gamma=(I_1,\dots,I_N)$ of finite windows if the following hold simultaneously along $\Gamma$:
\begin{itemize}[leftmargin=2em]
  \item each window $I_k$ in the chain carries the local control-hardening, local branch-hardening, and local loop/return-hardening certificates;
  \item the spinorial optimal line induces on each $I_k$ a locally readable and writable tape segment whose head-position update is compatible with the compressed computational surface;
  \item on every admissible overlap or concatenation interface $I_k\cap I_{k+1}$, the local tape segments glue to a persistent track without ambiguity in symbol transport, head continuation, or routed return-to-track placement;
  \item admissible resurface/reentry events may reopen the local track, but after each certified reentry the tape continuation rejoins the same globalized track class on the remaining suffix of $\Gamma$.
\end{itemize}
In words, $\mathbf T_{\mathrm{glob}}(\Gamma)=1$ means that the spinorial optimal line is no longer only a local tape candidate on isolated windows, but already behaves as a globally extensible effective tape support along the admissible chain $\Gamma$.
\end{definition}

\begin{proposition}[First reduction of the tape obligation $\mathbf U_T$]
Assume the conditional Zuse-core closure theorem, the reduced universality-obligation package \(\mathbf U=(\mathbf U_T,\mathbf U_Q,\mathbf U_B,\mathbf U_L,\mathbf U_S)\), the first Zuse-guided hardening priority order, and that the local control-hardening, local branch-hardening, and local loop/return-hardening certificates have already discharged \(\mathbf U_Q\), \(\mathbf U_B\), and \(\mathbf U_L\). Suppose moreover that there exists an admissible exhaustion by finite chains \(\Gamma\) on which the first tape-globalization certificate \(\mathbf T_{\mathrm{glob}}(\Gamma)=1\) holds and is stable under admissible extension of the chain. Then the obligation \(\mathbf U_T\) is discharged at the fourth hardening stage, and the remaining route reduces to
\[
(\mathbf U_S).
\]
Equivalently, once the spinorial optimal line supports a globally extensible read/write track compatible with the already hardened control, branch, and loop/return layers, the only unresolved machine-theoretic gap no longer lies in tape but solely in the explicit simulation of a standard reference machine.
\end{proposition}

\begin{corollary}[What the first globalization of $\mathbf U_T$ already buys]
The fourth hardening stage is now also no longer only a priority promise. It has become a concrete target: certify a globally extensible spinorial tape support along admissible chains and its stability under admissible extension and reentry. If that certificate is obtained, later revisions may treat tape as globally settled and focus entirely on the last remaining obligation, namely explicit reference-machine simulation.
\end{corollary}


\subsection*{19S. First reduction of the reference-machine simulation obligation $\mathbf U_S$}
Once control, branch, loop/return, and tape have been hardened, the final remaining machine-theoretic gap is no longer internal role identification but explicit comparison with a standard sequential reference core. The natural fifth step is therefore to isolate a simulation certificate which shows that the compressed computational surface can reproduce the transition semantics of a chosen admissible reference machine on the already globalized spinorial tape support.

\begin{definition}[First reference-machine simulation certificate for $\mathbf U_S$]
Write
\[
\mathbf S_{\mathrm{ref}}(\mathcal R,\Gamma)=1
\]
for the first reference-machine simulation certificate on an admissible chain $\Gamma$ for a chosen sequential reference core $\mathcal R$ if the following hold simultaneously:
\begin{itemize}[leftmargin=2em]
  \item there exists an explicit encoding of the machine state of $\mathcal R$ into the compressed computational surface, including control state, branch markers, loop/return status, head position, and tape symbols on the already globalized spinorial tape support;
  \item every one-step transition of $\mathcal R$ on the encoded configuration is reproduced by a finite admissible move sequence of the present line without loss of the hardened control, branch, loop/return, or tape certificates;
  \item halting, branch selection, and return-to-track events of $\mathcal R$ are reflected by the corresponding local cycle-control symbols and their decoder-driven actions on the encoded side;
  \item admissible extension of $\Gamma$ preserves the simulation relation on the longer prefix/suffix decomposition whenever the encoded reference computation itself extends.
\end{itemize}
In words, $\mathbf S_{\mathrm{ref}}(\mathcal R,\Gamma)=1$ means that the present line does not merely admit a primitive computational reading, but explicitly simulates the chosen sequential reference core $\mathcal R$ along the admissible chain $\Gamma$.
\end{definition}

\begin{proposition}[First reduction of the reference-machine simulation obligation $\mathbf U_S$]
Assume the conditional Zuse-core closure theorem, the reduced universality-obligation package \(\mathbf U=(\mathbf U_T,\mathbf U_Q,\mathbf U_B,\mathbf U_L,\mathbf U_S)\), the first Zuse-guided hardening priority order, and that the local control-hardening, local branch-hardening, local loop/return-hardening, and tape-globalization certificates have already discharged \(\mathbf U_Q\), \(\mathbf U_B\), \(\mathbf U_L\), and \(\mathbf U_T\). Suppose moreover that there exists a chosen admissible sequential reference core \(\mathcal R\) and an admissible exhaustion by finite chains \(\Gamma\) on which the first reference-machine simulation certificate \(\mathbf S_{\mathrm{ref}}(\mathcal R,\Gamma)=1\) holds and is stable under admissible extension. Then the final remaining obligation \(\mathbf U_S\) is discharged, and the reduced universality-obligation package collapses to the empty list.
\[
(\,).
\]
Equivalently, once the already hardened line explicitly simulates a chosen sequential reference core on the globalized spinorial tape support, the conditional Zuse-style primitive computational route is closed at the level of explicit reference-machine simulation.
\end{proposition}

\begin{corollary}[What the first reduction of $\mathbf U_S$ already buys]
The fifth hardening stage is now also no longer only a distant final promise. It has become a concrete target: produce an explicit admissible simulation of a chosen sequential reference core on the already hardened control/branch/loop/tape stack. If that certificate is obtained, later revisions may treat the conditional Zuse-core route as closed up to the chosen reference-machine level, while any stronger universality claim would then depend only on the strength of the reference core and the separate unboundedness assumptions built into that choice.
\end{corollary}


\subsection*{19T. Conditional reference-core closure surface}
Once the reduced universality-obligation package has collapsed to the empty list for a chosen admissible sequential reference core, the computational reading should no longer be cited only as a chain of discharged obligations. The right next compression is therefore a single bundled closure surface recording that the hardened line now simulates the chosen reference core on the globalized spinorial tape support.

\begin{definition}[Conditional reference-core closure surface]
Let $\mathcal R$ be a chosen admissible sequential reference core and let $(\Gamma_n)_{n\ge 1}$ be an admissible exhaustion by finite chains. Write
\[
\mathfrak Z^{\mathrm{ref}}(\mathcal R)=1
\]
if the following hold simultaneously:
\begin{itemize}[leftmargin=2em]
  \item the conditional Zuse-core closure theorem holds;
  \item the local control-hardening, local branch-hardening, local loop/return-hardening, tape-globalization, and reference-machine simulation certificates discharge $\mathbf U_Q$, $\mathbf U_B$, $\mathbf U_L$, $\mathbf U_T$, and $\mathbf U_S$ on the exhaustion $(\Gamma_n)_{n\ge 1}$;
  \item the explicit simulation relation for $\mathcal R$ is stable under admissible extension of the finite chains in the exhaustion.
\end{itemize}
In words, $\mathfrak Z^{\mathrm{ref}}(\mathcal R)=1$ means that the present line is already closed, in the conditional sense of this chapter, up to the explicit simulation of the chosen reference core $\mathcal R$.
\end{definition}

\begin{theorem}[Conditional reference-core closure theorem]
Assume the conditional Zuse-core closure theorem and the hypotheses of the first reduction of the reference-machine simulation obligation $\mathbf U_S$. Suppose moreover that the discharged hardening certificates remain stable under admissible extension along an exhaustion of finite chains for a chosen admissible sequential reference core $\mathcal R$. Then the compressed computational interpretation surface $\mathfrak C^{\mathrm{comp}}$ closes to a \emph{conditional reference-core closure surface} for $\mathcal R$:
\[
\mathfrak Z^{\mathrm{ref}}(\mathcal R)=1.
\]
Equivalently, the present line is no longer merely a conditional primitive Zuse-style computational core, but a conditionally closed simulator of the chosen reference core on the globalized spinorial tape support. Any stronger universality claim therefore depends only on the strength of the chosen reference core and on whatever additional unboundedness assumptions one wishes to attach to it.
\end{theorem}

\begin{corollary}[What the conditional reference-core closure surface already buys]
Later arguments may now cite a single bundled machine-theoretic endpoint: for a chosen admissible sequential reference core $\mathcal R$, the present line is conditionally closed up to explicit reference-core simulation. Thus one no longer needs to reopen the intermediate hardening ladder unless the proof specifically aims to strengthen the chosen reference core or to pass from reference-core closure to a stronger universality theorem.
\end{corollary}


\begin{remark}[Deferred LLM-superstructure note]
The present line does not claim that it already constitutes a large language model in the strict modern sense. What is now reasonable to record for later work is a weaker architectural possibility: the combination of graph-structured objects, routed node logic, history-sensitive compression, cache-like admissible reuse along an optimal line, and the conditional Zuse-style primitive computational core may eventually support an \emph{LLM-superstructure reading}. In that later reading, the current release line would function not as the trained language-model core itself, but as a host/control/memory architecture into which a genuine next-symbol predictor or other trainable linguistic layer could be inserted. This is therefore deferred as a future-work hypothesis rather than used as an active theorematic ingredient. It may, however, be rechecked opportunistically in later revisions whenever the token/readout, predictive, or trainability side becomes structurally clearer.
\end{remark}


\subsection*{19U. Strength hierarchy above conditional reference-core closure}

\begin{definition}[First machine-theoretic strength hierarchy on the compressed computational surface]
For the compressed computational interpretation surface $\mathfrak C^{\mathrm{comp}}$, define the following five nested strength levels:
\[
\begin{aligned}
\mathbf S_1&:=\text{disciplined computational reading},\\
\mathbf S_2&:=\text{conditional Zuse-style primitive core closure},\\
\mathbf S_3(\mathcal R)&:=\text{conditional reference-core closure for a chosen sequential reference core }\mathcal R,\\
\mathbf S_4(\mathcal R)&:=\parbox[t]{0.68\linewidth}{conditional universality relative to a chosen reference core $\mathcal R$ strong enough for the intended claim},\\
\mathbf S_5&:=\text{strict unbounded universality / Turing-style completeness}.
\end{aligned}
\]
\end{definition}

\begin{proposition}[Monotonicity of the first machine-theoretic strength hierarchy]
On the admissible computational window one has the one-way implication chain
\[
\mathbf S_5\Longrightarrow \mathbf S_4(\mathcal R)\Longrightarrow \mathbf S_3(\mathcal R)\Longrightarrow \mathbf S_2\Longrightarrow \mathbf S_1.
\]
In particular, every stronger machine-theoretic endpoint automatically contains all weaker lower-level readings, while the converse implications need not hold without additional unboundedness or simulation arguments.
\end{proposition}

\begin{corollary}[What the first machine-theoretic strength hierarchy already buys]
The present line may now be located unambiguously in the hierarchy: by the conditional reference-core closure theorem it reaches $\mathbf S_3(\mathcal R)$ for a chosen admissible sequential reference core $\mathcal R$. What remains open is therefore no longer the primitive computational reading itself, but the step from $\mathbf S_3(\mathcal R)$ to stronger levels $\mathbf S_4(\mathcal R)$ and $\mathbf S_5$. Later arguments can thus distinguish cleanly between primitive closure, reference-core closure, conditional universality claims, and strict Turing-style completeness claims.
\end{corollary}


\subsection*{19V. First upgrade criterion from reference-core closure to conditional universality}
Once the present line has reached \(\mathbf S_3(\mathcal R)\) for a chosen admissible sequential reference core, the remaining machine-theoretic question is no longer whether the line simulates that reference core, but whether the chosen reference core is itself strong enough to transport an intended universality claim back through the already established simulation relation.

\begin{definition}[First relative universality-upgrade certificate above \(\mathbf S_3(\mathcal R)\)]
Let \(\mathcal R\) be a chosen admissible sequential reference core and let \(\mathcal K\) be a target class of sequential computations one wishes to cover. Write
\[
\mathfrak U^{\mathrm{rel}}(\mathcal R,\mathcal K)=1
\]
if the following hold simultaneously:
\begin{itemize}[leftmargin=2em]
  \item the conditional reference-core closure surface satisfies \(\mathfrak Z^{\mathrm{ref}}(\mathcal R)=1\);
  \item the reference core \(\mathcal R\) carries a universality theorem for the target class \(\mathcal K\) under a specified admissible encoding/decoding discipline;
  \item the simulation relation used in \(\mathfrak Z^{\mathrm{ref}}(\mathcal R)=1\) transports that admissible encoding/decoding discipline without breaking the globalized tape support, the hardened control/branch/loop stack, or the local readout semantics.
\end{itemize}
In words, \(\mathfrak U^{\mathrm{rel}}(\mathcal R,\mathcal K)=1\) means that the present line is not only closed up to the chosen reference core, but conditionally inherits the universality of \(\mathcal R\) for the intended target class \(\mathcal K\).
\end{definition}

\begin{proposition}[First conditional upgrade from \(\mathbf S_3(\mathcal R)\) to \(\mathbf S_4(\mathcal R)\)]
Assume \(\mathfrak Z^{\mathrm{ref}}(\mathcal R)=1\) for a chosen admissible sequential reference core \(\mathcal R\). Suppose moreover that \(\mathcal R\) is already known to be universal for a target class \(\mathcal K\) under an admissible encoding/decoding discipline, and that this discipline is preserved by the conditional reference-core closure surface. Then the present line upgrades from \(\mathbf S_3(\mathcal R)\) to the relative universality level \(\mathbf S_4(\mathcal R)\) for the intended class \(\mathcal K\).
\[
\mathfrak U^{\mathrm{rel}}(\mathcal R,\mathcal K)=1
\Longrightarrow
\mathbf S_4(\mathcal R).
\]
Equivalently, no additional local hardening layer is required beyond \(\mathbf S_3(\mathcal R)\); what remains is the external strength theorem for the chosen reference core and the proof that its admissible coding discipline is faithfully transported by the already hardened simulation route.
\end{proposition}

\begin{corollary}[What the first upgrade criterion above \(\mathbf S_3(\mathcal R)\) already buys]
The gap between the present line and conditional universality is now sharply isolated. To pass from the current reference-core closure level to \(\mathbf S_4(\mathcal R)\), later revisions no longer need to reopen the primitive computational hardening ladder. They need only choose a reference core \(\mathcal R\) strong enough for the intended target class \(\mathcal K\), certify the corresponding universality theorem for \(\mathcal R\), and verify that the admissible encoding/decoding discipline survives transport through the already established reference-core closure surface.
\end{corollary}


\subsection*{19W. First admissible reference-core selection principle above \(\mathbf S_3(\mathcal R)\)}
Once the upgrade from \(\mathbf S_3(\mathcal R)\) to \(\mathbf S_4(\mathcal R)\) has been isolated, the next question is no longer whether one can choose some strong reference core at all, but which admissible reference core should be chosen first. A structure-first route should avoid importing unnecessary machine-theoretic strength too early; it should instead select the weakest admissible reference core already sufficient for the intended target claim.

\begin{definition}[First conservative reference-core selection certificate above \(\mathbf S_3(\mathcal R)\)]
Fix a target class \(\mathcal K\) and a chosen comparison family \(\mathfrak F_{\mathrm{ref}}\) of admissible sequential reference cores together with a strength preorder \(\preceq_{\mathrm{ref}}\). Write
\[
\mathfrak S^{\mathrm{ref}}(\mathcal R,\mathcal K)=1
\]
if the following hold:
\begin{itemize}[leftmargin=2em]
  \item \(\mathcal R\in\mathfrak F_{\mathrm{ref}}\) is admissible and satisfies \(\mathfrak Z^{\mathrm{ref}}(\mathcal R)=1\);
  \item \(\mathcal R\) is strong enough to support the intended target claim for \(\mathcal K\) under an admissible encoding/decoding discipline;
  \item for every \(\mathcal R'\in\mathfrak F_{\mathrm{ref}}\) with \(\mathcal R'\prec_{\mathrm{ref}}\mathcal R\), the same target claim for \(\mathcal K\) is not yet certified under the same admissible coding discipline.
\end{itemize}
Thus \(\mathfrak S^{\mathrm{ref}}(\mathcal R,\mathcal K)=1\) means that \(\mathcal R\) is not merely an admissible reference core, but the first conservatively sufficient one inside the chosen comparison family for the intended upgrade problem.
\end{definition}

\begin{proposition}[First conservative selection reduction for the universality-upgrade route]
Assume \(\mathfrak U^{\mathrm{rel}}(\mathcal R,\mathcal K)=1\) and \(\mathfrak S^{\mathrm{ref}}(\mathcal R,\mathcal K)=1\). Then the passage from \(\mathbf S_3(\mathcal R)\) to \(\mathbf S_4(\mathcal R)\) may be formulated with no stronger external reference-core assumption than the one already needed for the chosen target class \(\mathcal K\) inside the comparison family \(\mathfrak F_{\mathrm{ref}}\). In particular, later revisions do not need to import a stronger machine model than necessary in order to state the first conditional universality claim above the present reference-core closure level.
\end{proposition}

\begin{corollary}[What the first conservative reference-core selection principle already buys]
The relative universality route is now disciplined in two distinct ways: first by the already established reference-core closure surface, and second by a conservative selection rule on the external endpoint. Later upgrades above \(\mathbf S_3(\mathcal R)\) can therefore proceed by choosing the weakest admissible reference core sufficient for the intended claim, rather than by jumping prematurely to an unnecessarily strong machine-theoretic endpoint.
\end{corollary}


\subsection*{19X. First admissible tape-based reference-core family above \(\mathbf S_3(\mathcal R)\)}
The conservative selection principle above becomes materially useful only once one fixes a first concrete comparison family of admissible reference cores. The present line already carries explicit tape, control, branch, and loop/return roles, so the most conservative first family should preserve exactly that architecture instead of importing richer machine features prematurely.

\begin{definition}[First admissible tape-based reference-core family]
Let \(\mathfrak F_{\mathrm{tape}}\) denote the class of admissible sequential reference cores \(\mathcal R\) satisfying all of the following:
\begin{itemize}[leftmargin=2em]
  \item \(\mathcal R\) has finite control state, finite local symbol alphabet, and a one-track sequential update discipline;
  \item every local step of \(\mathcal R\) is determined by the currently scanned symbol together with the current control state and returns one next control state, one write/update instruction, and one move instruction in \(\{-1,0,+1\}\);
  \item branch resolution is explicit and local, in the sense that each step chooses one of finitely many admissible successor clauses by a finite Boolean/projector-type test on the current scanned data;
  \item loop/return structure is explicit and local, in the sense that every admissible return event is encoded by a finite control-side return symbol together with a deterministic reentry instruction on the same one-track support;
  \item the tape support itself is one-dimensional and sequential, so that the already globalized spinorial optimal line may serve as the direct comparison support without introducing any extra addressing geometry.
\end{itemize}
Write
\[
\mathfrak A^{\mathrm{tape}}(\mathcal R)=1
\]
if \(\mathcal R\in\mathfrak F_{\mathrm{tape}}\).
\end{definition}

\begin{proposition}[First conservative tape-family reduction for the upgrade route]
Assume the conditional reference-core closure surface \(\mathfrak Z^{\mathrm{ref}}(\mathcal R)=1\) for some chosen admissible sequential reference core \(\mathcal R\), together with the conservative selection discipline above. Suppose moreover that the intended universality-upgrade problem concerns only target classes whose admissible coding discipline is already one-track, sequential, and local in the same sense as the hardened tape/control/branch/loop architecture of the present line. Then the first conservative search for a reference core may be restricted to the admissible tape-based family \(\mathfrak F_{\mathrm{tape}}\) without loss for that upgrade problem.

Equivalently, whenever the intended target class does not yet require richer addressing, parallel branching geometry, or nonlocal memory moves, the weakest natural external comparison family is the tape-based family already matched to the present hardened architecture.
\end{proposition}

\begin{corollary}[What the first admissible tape-based reference-core family already buys]
The external universality-upgrade route is now no longer open-ended. Later revisions can first search for a minimally sufficient \emph{tape-based} reference core before importing any richer machine model. This keeps the passage above \(\mathbf S_3(\mathcal R)\) aligned with the already hardened internal architecture and reduces the chance of overstating the needed external strength.
\end{corollary}

\subsection*{19Y. First internal-to-tape transport certificate above \(\mathbf S_3(\mathcal R)\)}
Choosing a first admissible tape-based comparison family becomes materially useful only once one can express the existing reference-core closure surface in a way that is transport-compatible with that family. The next conservative step is therefore not a stronger universality claim, but a certificate saying that the already hardened internal read/write/control/branch/return package can be read inside a chosen tape-based reference core without changing the local computational discipline.

\begin{definition}[First internal-to-tape transport certificate above \(\mathbf S_3(\mathcal R)\)]
Let \(\mathcal R\in\mathfrak F_{\mathrm{tape}}\) satisfy \(\mathfrak Z^{\mathrm{ref}}(\mathcal R)=1\). Write
\[
\mathfrak T^{\mathrm{tape}}(\mathcal R)=1
\]
if the following hold on every admissible later-depth window used by the present compressed computational surface:
\begin{itemize}[leftmargin=2em]
  \item the globalized spinorial optimal line induces a one-dimensional sequential support compatible with the tape support of \(\mathcal R\);
  \item the already hardened local control state of the present line is encoded into the finite control state of \(\mathcal R\) without introducing extra addressing geometry;
  \item the already hardened projector/Boolean branch layer is transported into the branch clauses of \(\mathcal R\) by a local admissible encoding/decoding discipline;
  \item the already hardened routed HC-selected return/closure step is transported into the local return/reentry discipline of \(\mathcal R\) without changing the sequential one-track support;
  \item one local step of the present compressed computational surface is transported to one admissible local step of \(\mathcal R\) up to the chosen local coding discipline.
\end{itemize}
Thus \(\mathfrak T^{\mathrm{tape}}(\mathcal R)=1\) means that the internal reference-core closure already lives on a support that can be read directly inside the first admissible tape-based comparison family.
\end{definition}

\begin{proposition}[First transport reduction from reference-core closure to the tape-based family]
Assume \(\mathfrak Z^{\mathrm{ref}}(\mathcal R)=1\), \(\mathfrak A^{\mathrm{tape}}(\mathcal R)=1\), and \(\mathfrak T^{\mathrm{tape}}(\mathcal R)=1\). Then the external upgrade route above \(\mathbf S_3(\mathcal R)\) may be reformulated entirely inside the admissible tape-based comparison family without reopening the already hardened local obligations for control, branch, loop/return, or tape support. In particular, the passage to the first conditional universality-upgrade problem may be stated as an external strength question about the chosen tape-based reference core itself, rather than as a renewed local hardening problem for the present line.
\end{proposition}

\begin{corollary}[What the first internal-to-tape transport certificate already buys]
The tape-based upgrade route is now disciplined in three layers: by the conditional reference-core closure surface, by the conservative tape-family selection principle, and by an explicit transport certificate from the present compressed computational surface into that family. Later work above \(\mathbf S_3(\mathcal R)\) can therefore focus directly on the external strength of the chosen tape-based reference core, instead of rechecking the already stabilized local computational roles.
\end{corollary}


\subsection*{19Z. First tape-family relative universality witness above \(\mathbf S_3(\mathcal R)\)}
Once the admissible tape-based family has been selected and the internal reference-core closure has been transported into it, the remaining external question is no longer whether the present line can be read inside that family, but whether a chosen member of that family is already strong enough for the intended target class. The next conservative step is therefore to isolate the first explicit witness certifying that a transported tape-based reference core upgrades the present line from conditional reference-core closure to conditional universality relative to that core.

\begin{definition}[First tape-family relative universality witness above \(\mathbf S_3(\mathcal R)\)]
Let \(\mathcal R\in\mathfrak F_{\mathrm{tape}}\) satisfy
\[
\mathfrak Z^{\mathrm{ref}}(\mathcal R)=1,
\qquad
\mathfrak A^{\mathrm{tape}}(\mathcal R)=1,
\qquad
\mathfrak T^{\mathrm{tape}}(\mathcal R)=1.
\]
For a target class \(\mathcal K\), write
\[
\mathfrak U^{\mathrm{tape}}(\mathcal R,\mathcal K)=1
\]
if the following hold:
\begin{itemize}[leftmargin=2em]
  \item the chosen tape-based reference core \(\mathcal R\) is already known to simulate every machine in the intended target class \(\mathcal K\) under a one-track sequential local coding discipline compatible with \(\mathfrak F_{\mathrm{tape}}\);
  \item the admissible encoding/decoding discipline used in that simulation is stable under the internal-to-tape transport certificate \(\mathfrak T^{\mathrm{tape}}(\mathcal R)=1\);
  \item no extra addressing geometry, nonlocal branching discipline, or richer memory move is required beyond the tape-based family already selected.
\end{itemize}
Thus \(\mathfrak U^{\mathrm{tape}}(\mathcal R,\mathcal K)=1\) means that the chosen transported tape-based reference core is externally strong enough to witness the first conditional universality claim relevant to \(\mathcal K\).
\end{definition}

\begin{proposition}[First tape-family upgrade from \(\mathbf S_3(\mathcal R)\) to \(\mathbf S_4(\mathcal R)\)]
Assume \(\mathfrak Z^{\mathrm{ref}}(\mathcal R)=1\), \(\mathfrak A^{\mathrm{tape}}(\mathcal R)=1\), \(\mathfrak T^{\mathrm{tape}}(\mathcal R)=1\), and \(\mathfrak U^{\mathrm{tape}}(\mathcal R,\mathcal K)=1\). Then the present compressed computational surface upgrades from the reference-core closure level \(\mathbf S_3(\mathcal R)\) to the relative universality level \(\mathbf S_4(\mathcal R)\) for the intended target class \(\mathcal K\).

Equivalently, once the internal closure surface has been transported into a sufficiently strong admissible tape-based reference core, the remaining strength step above \(\mathbf S_3(\mathcal R)\) is discharged entirely by the external universality witness for that transported tape-based core.
\end{proposition}

\begin{corollary}[What the first tape-family relative universality witness already buys]
The universality-upgrade route above the present compressed computational surface is now disciplined in four layers: by the conditional reference-core closure surface, by the conservative tape-family selection rule, by the internal-to-tape transport certificate, and by an explicit external universality witness for a chosen transported tape-based reference core. Later revisions can therefore state the first conditional universality claim without reopening the already discharged primitive machine-theoretic hardening ladder.
\end{corollary}




\subsection*{19AA. First strict universality-obligation package above \(\mathbf S_4(\mathcal R)\)}
Once the present line has reached the relative universality level \(\mathbf S_4(\mathcal R)\) for a chosen transported tape-based reference core and target class, the remaining question is no longer whether the primitive computational roles are available, nor whether a conservative tape-based reference family can witness relative universality. What remains is the stricter global question whether the transported tape-based route supports the genuinely unbounded extension needed for strict Turing-style completeness. The next conservative step is therefore to isolate that last gap as a separate package instead of mixing it back into the already discharged local hardening ladder.

\begin{definition}[First strict universality-obligation package above \(\mathbf S_4(\mathcal R)\)]
Assume that \(\mathbf S_4(\mathcal R)\) has been reached for an intended target class \(\mathcal K\). Write
\[
\mathbf V=(\mathbf V_\infty,\mathbf V_{\mathrm{ext}},\mathbf V_{\mathrm{dec}})
\]
for the first strict universality-obligation package above \(\mathbf S_4(\mathcal R)\), where:
\begin{itemize}[leftmargin=2em]
  \item \(\mathbf V_\infty\) is the obligation that the chosen transported tape-based reference route admits arbitrarily long admissible extensions rather than only bounded local windows;
  \item \(\mathbf V_{\mathrm{ext}}\) is the obligation that the already established tape-globalization and concatenation discipline extends coherently along such arbitrarily long admissible chains, without hidden global collapse or cutoff;
  \item \(\mathbf V_{\mathrm{dec}}\) is the obligation that the admissible encoding/decoding discipline witnessing \(\mathbf S_4(\mathcal R)\) remains stable under those unbounded extensions and does not depend on a fixed finite window budget.
\end{itemize}
Thus \(\mathbf V\) isolates the first strict machine-theoretic gap between relative universality inside a chosen tape-based reference-core route and full unbounded Turing-style completeness.
\end{definition}

\begin{proposition}[First isolation of the strict universality gap above \(\mathbf S_4(\mathcal R)\)]
Assume that \(\mathbf S_4(\mathcal R)\) has been reached through the conditional reference-core closure surface, the conservative tape-family selection rule, the internal-to-tape transport certificate, and an explicit tape-family relative universality witness. Then the only remaining machine-theoretic step toward the strict level \(\mathbf S_5\) is discharged by the package \(\mathbf V\) above.

Equivalently, once the present line has been upgraded to \(\mathbf S_4(\mathcal R)\), no additional local hardening of control, branch, loop/return, tape support, or reference-core simulation is needed. What remains is exclusively the global unboundedness, extension, and decoding-stability problem recorded by \(\mathbf V\).
\end{proposition}

\begin{corollary}[What the first strict universality-obligation package already buys]
The machine-theoretic hierarchy is now disciplined at both of its remaining upper gaps. The passage from \(\mathbf S_3(\mathcal R)\) to \(\mathbf S_4(\mathcal R)\) is controlled by the conservative tape-family upgrade route, while the passage from \(\mathbf S_4(\mathcal R)\) to \(\mathbf S_5\) is now isolated as the strict universality-obligation package \(\mathbf V\). Later revisions can therefore state very clearly whether they are addressing relative universality inside a chosen transported tape-based reference family or the separate problem of genuine unbounded Turing-style completeness.
\end{corollary}


\subsection*{19AB. Hardening priority for the strict universality obligation above \(\mathbf S_4(\mathcal R)\)}
Once the strict universality gap above \(\mathbf S_4(\mathcal R)\) has been isolated as the package \(\mathbf V=(\mathbf V_\infty,\mathbf V_{\mathrm{ext}},\mathbf V_{\mathrm{dec}})\), the next conservative step is not to attack all three remaining obligations at once. As on the earlier Zuse-guided route, the cleaner strategy is to impose a first hardening order, so that later revisions can reduce the strict universality gap in the same staged way as the primitive computational obligations.

\begin{definition}[First hardening priority order for the strict universality-obligation package above \(\mathbf S_4(\mathcal R)\)]
Assume that the present compressed computational surface has reached \(\mathbf S_4(\mathcal R)\) for an intended target class \(\mathcal K\), and let
\[
\mathbf V=(\mathbf V_\infty,\mathbf V_{\mathrm{ext}},\mathbf V_{\mathrm{dec}})
\]
be the first strict universality-obligation package above \(\mathbf S_4(\mathcal R)\).
Define the first hardening priority order on \(\mathbf V\) by
\[
\mathbf H_2=(\mathbf V_{\mathrm{ext}},\mathbf V_{\mathrm{dec}},\mathbf V_\infty).
\]
Thus the conservative strengthening strategy is:
\begin{itemize}[leftmargin=2em]
  \item first harden the global extension obligation \(\mathbf V_{\mathrm{ext}}\),
  \item then harden decoding stability under extension \(\mathbf V_{\mathrm{dec}}\),
  \item and only at the end attack genuine unboundedness \(\mathbf V_\infty\).
\end{itemize}
\end{definition}

\begin{proposition}[First priority reduction for the strict universality-obligation package above \(\mathbf S_4(\mathcal R)\)]
Assume that \(\mathbf S_4(\mathcal R)\) has been reached and that the hardening order \(\mathbf H_2\) above is adopted. Then later revisions can reduce the strict universality gap monotonically: any successful discharge of the first two obligations \(\mathbf V_{\mathrm{ext}}\) and \(\mathbf V_{\mathrm{dec}}\) shrinks the remaining open strict gap to the single final problem \(\mathbf V_\infty\).

Equivalently, once the extension discipline and its decoding stability have been hardened, the only remaining obstruction to \(\mathbf S_5\) is the genuinely unbounded continuation problem rather than the local or semiglobal coherence of the transported tape-based route.
\end{proposition}

\begin{corollary}[What the first hardening priority order above \(\mathbf S_4(\mathcal R)\) already buys]
The strict universality gap is now not only isolated but also staged. Later revisions can therefore work above \(\mathbf S_4(\mathcal R)\) without mixing extension, decoding, and unboundedness into a single undifferentiated final problem. In the conservative reading, the current next target is no longer \(\mathbf S_5\) as a whole but specifically the first extension-hardening step \(\mathbf V_{\mathrm{ext}}\).
\end{corollary}


\subsection*{19AC. First local hardening of the extension obligation \(\mathbf V_{\mathrm{ext}}\) above \(\mathbf S_4(\mathcal R)\)}
With the strict universality-obligation package now ordered by \(\mathbf H_2\), the next conservative step is to attack the first remaining upper-gap component \(\mathbf V_{\mathrm{ext}}\). The question is no longer whether the present line supports a finite transported tape-based computation relative to a chosen reference core, but whether those admissible windows extend coherently across longer chains without hidden collapse, cutoff, or incompatible seam behavior.

\begin{definition}[First extension-hardening certificate for \(\mathbf V_{\mathrm{ext}}\) above \(\mathbf S_4(\mathcal R)\)]
Assume \(\mathbf S_4(\mathcal R)\) for an intended target class \(\mathcal K\).
Write
\[
\mathfrak E^{\mathrm{glob}}(\mathcal R)=1
\]
if the following hold:
\begin{itemize}[leftmargin=2em]
  \item every admissible finite transported tape-window in the chosen family extends to a longer admissible tape-window by a compatible local continuation step;
  \item the already established concatenation/cut and patch-normal-form discipline remains coherent under those longer continuations;
  \item no hidden global seam defect, cutoff phenomenon, or incompatible local reindexing appears when adjoining admissible windows along the transported tape-line.
\end{itemize}
Thus \(\mathfrak E^{\mathrm{glob}}(\mathcal R)=1\) certifies the first coherent extension-hardening of the tape-based route above \(\mathbf S_4(\mathcal R)\).
\end{definition}

\begin{proposition}[First reduction of the extension obligation \(\mathbf V_{\mathrm{ext}}\) above \(\mathbf S_4(\mathcal R)\)]
Assume \(\mathbf S_4(\mathcal R)\) and \(\mathfrak E^{\mathrm{glob}}(\mathcal R)=1\). Then the first strict universality-obligation package reduces from
\[
\mathbf V=(\mathbf V_\infty,\mathbf V_{\mathrm{ext}},\mathbf V_{\mathrm{dec}})
\]
to
\[
\mathbf V'=(\mathbf V_{\mathrm{dec}},\mathbf V_\infty).
\]
Equivalently, once the transported tape-based route is globally extensible in the admissible-window sense, the strict universality gap above \(\mathbf S_4(\mathcal R)\) no longer contains a separate extension obstruction; what remains are the decoding-stability problem under such extension and the final genuinely unbounded continuation problem.
\end{proposition}

\begin{corollary}[What the first local hardening of \(\mathbf V_{\mathrm{ext}}\) above \(\mathbf S_4(\mathcal R)\) already buys]
The upper strict gap above relative universality is now cleaner than before. If the extension-hardening certificate holds, later revisions do not need to reopen the global seam/coherence problem of transported admissible tape-windows. The next active upper-gap target becomes decoding stability under extension, with genuine unboundedness kept as the final separate endpoint problem.
\end{corollary}


\subsection*{19AD. First local hardening of the decoding-stability obligation \(\mathbf V_{\mathrm{dec}}\) above \(\mathbf S_4(\mathcal R)\)}
Once the extension obligation has been isolated and locally hardened, the next conservative strict-universality target is no longer the seam/coherence of admissible transported tape-windows themselves, but the stability of the admissible encoding/decoding discipline under those longer extensions. The question is now whether the same transported tape-based route that already witnesses relative universality continues to admit a coherent decoding discipline when admissible windows are enlarged, concatenated, and reread along longer chains.

\begin{definition}[First decoding-stability certificate for \(\mathbf V_{\mathrm{dec}}\) above \(\mathbf S_4(\mathcal R)\)]
Assume \(\mathbf S_4(\mathcal R)\) for an intended target class \(\mathcal K\), and assume moreover that the global extension certificate \(\mathfrak E^{\mathrm{glob}}(\mathcal R)=1\) has already been established.
Write
\[
\mathfrak D^{\mathrm{ext}}(\mathcal R)=1
\]
if the following hold:
\begin{itemize}[leftmargin=2em]
  \item the admissible decoding discipline used to witness \(\mathbf S_4(\mathcal R)\) on finite transported tape-windows extends coherently to every admissible longer transported tape-window produced by the extension certificate;
  \item restriction from a longer admissible transported tape-window to any admissible subwindow commutes with the chosen decoding/readout rule;
  \item admissible concatenation of transported tape-windows does not introduce a new decoding ambiguity, hidden branch aliasing, or context-dependent reinterpretation of already decoded local machine states.
\end{itemize}
Thus \(\mathfrak D^{\mathrm{ext}}(\mathcal R)=1\) certifies the first stable decoding discipline under admissible transported tape-extension above \(\mathbf S_4(\mathcal R)\).
\end{definition}

\begin{proposition}[First reduction of the decoding-stability obligation \(\mathbf V_{\mathrm{dec}}\) above \(\mathbf S_4(\mathcal R)\)]
Assume \(\mathbf S_4(\mathcal R)\), \(\mathfrak E^{\mathrm{glob}}(\mathcal R)=1\), and \(\mathfrak D^{\mathrm{ext}}(\mathcal R)=1\). Then the strict universality-obligation package reduces from
\[
\mathbf V=(\mathbf V_\infty,\mathbf V_{\mathrm{ext}},\mathbf V_{\mathrm{dec}})
\]
to
\[
\mathbf V''=(\mathbf V_\infty).
\]
Equivalently, once the transported tape-based route is globally extensible and its admissible decoding discipline remains stable under those extensions, the strict universality gap above \(\mathbf S_4(\mathcal R)\) no longer contains a separate decoding obstruction; what remains is only the final genuinely unbounded continuation problem.
\end{proposition}

\begin{corollary}[What the first local hardening of \(\mathbf V_{\mathrm{dec}}\) above \(\mathbf S_4(\mathcal R)\) already buys]
The upper strict gap has now been reduced as far as it can be reduced without addressing genuine unboundedness itself. If the extension-hardening and decoding-stability certificates both hold, later revisions no longer need to reopen finite-window extension or decoding coherence. The only remaining active upper-gap target is then \(\mathbf V_\infty\), namely true unbounded continuation.
\end{corollary}


\subsection*{19AE. Local hardening of genuine unboundedness \(\mathbf V_\infty\) above \(\mathbf S_4(\mathcal R)\)}
Once global extension and decoding stability under extension have both been isolated and locally hardened, the only remaining strict-universality target is no longer any finite-window transport or rereading problem. What remains is the genuinely unbounded continuation problem itself: whether the transported tape-based route above \(\mathbf S_4(\mathcal R)\) admits arbitrarily extendable admissible continuations without forced collapse, periodic truncation, or hidden resource cutoff.

\begin{definition}[First genuine-unboundedness certificate for \(\mathbf V_\infty\) above \(\mathbf S_4(\mathcal R)\)]
Assume \(\mathbf S_4(\mathcal R)\), \(\mathfrak E^{\mathrm{glob}}(\mathcal R)=1\), and \(\mathfrak D^{\mathrm{ext}}(\mathcal R)=1\).
Write
\[
\mathfrak U^{\infty}(\mathcal R)=1
\]
if the following hold:
\begin{itemize}[leftmargin=2em]
  \item for every admissible transported tape-window in the chosen tape-based family there exists a strictly longer admissible transported tape-window extending it;
  \item admissible extension can be iterated indefinitely without forcing a terminal collapse, bounded cycling artifact, or hidden global cutoff that would prevent arbitrarily long transported computations;
  \item the already established extension and decoding certificates remain valid along every finite stage of such iterated admissible continuation.
\end{itemize}
Thus \(\mathfrak U^{\infty}(\mathcal R)=1\) certifies the first genuinely unbounded admissible continuation regime above \(\mathbf S_4(\mathcal R)\).
\end{definition}

\begin{proposition}[First reduction of the genuine unboundedness obligation \(\mathbf V_\infty\) above \(\mathbf S_4(\mathcal R)\)]
Assume \(\mathbf S_4(\mathcal R)\), \(\mathfrak E^{\mathrm{glob}}(\mathcal R)=1\), \(\mathfrak D^{\mathrm{ext}}(\mathcal R)=1\), and \(\mathfrak U^{\infty}(\mathcal R)=1\). Then the strict universality-obligation package reduces from
\[
\mathbf V=(\mathbf V_\infty,\mathbf V_{\mathrm{ext}},\mathbf V_{\mathrm{dec}})
\]
to the empty list.
Equivalently, once the transported tape-based route is globally extensible, decoding-stable under extension, and genuinely unbounded in the admissible continuation sense, the strict universality gap above \(\mathbf S_4(\mathcal R)\) contains no further separate obstruction. In that case the compressed computational surface reaches the top level \(\mathbf S_5\), namely strict unbounded Turing-style completeness relative to the certified route.
\end{proposition}

\begin{corollary}[What the first local hardening of \(\mathbf V_\infty\) above \(\mathbf S_4(\mathcal R)\) already buys]
The upper gap is now completely structured. If the global extension, decoding-stability, and genuine-unboundedness certificates all hold, later revisions do not need to reopen any remaining strict-universality obstruction above \(\mathbf S_4(\mathcal R)\). The only remaining caution is interpretive: one must still keep explicit which route, tape-family, and reference-core class are being certified when speaking of strict completeness.
\end{corollary}


\subsection*{19AF. Full machine-theoretic closure surface on the compressed computational line}
After the conditional reference-core closure surface, the relative universality-upgrade witness, and the strict upper certificate package have all been isolated, the machine-theoretic reading should no longer appear as a long ladder of locally discharged obligations. It can now be bundled into a single closure surface that records exactly which level of the hierarchy has been reached under which certificate discipline.

\begin{definition}[Full machine-theoretic closure surface on the compressed computational line]
Fix a chosen admissible sequential reference core \(\mathcal R\), a chosen admissible tape-based reference family, and the associated conservative transport discipline. Write
\[
\mathfrak M^{\mathrm{full}}(\mathcal R)=1
\]
if the following package holds simultaneously:
\begin{itemize}[leftmargin=2em]
  \item the compressed computational interpretation surface reaches the conditional reference-core closure level \(\mathbf S_3(\mathcal R)\), i.e. \(\mathfrak Z^{\mathrm{ref}}(\mathcal R)=1\);
  \item the tape-family upgrade route is certified, i.e. the conservative tape-family selection rule, the internal-to-tape transport certificate, and the tape-family relative universality witness all hold, so that the line reaches \(\mathbf S_4(\mathcal R)\);
  \item the strict upper certificate package is certified, i.e.
  \[
  \mathfrak E^{\mathrm{glob}}(\mathcal R)=1,\qquad
  \mathfrak D^{\mathrm{ext}}(\mathcal R)=1,\qquad
  \mathfrak U^{\infty}(\mathcal R)=1,
  \]
  so that the remaining strict gap above \(\mathbf S_4(\mathcal R)\) is empty.
\end{itemize}
Thus \(\mathfrak M^{\mathrm{full}}(\mathcal R)=1\) denotes the first fully bundled machine-theoretic closure surface attached to the present compressed computational line relative to the chosen certified route.
\end{definition}

\begin{theorem}[First bundled machine-theoretic closure theorem on the compressed computational line]
Assume \(\mathfrak M^{\mathrm{full}}(\mathcal R)=1\). Then the compressed computational line reaches the top machine-theoretic level \(\mathbf S_5\) along the chosen certified route. Equivalently, all intermediate gaps
\[
\mathbf S_1\longrightarrow \mathbf S_2\longrightarrow \mathbf S_3(\mathcal R)\longrightarrow \mathbf S_4(\mathcal R)\longrightarrow \mathbf S_5
\]
are discharged under an explicitly named certificate package, so that no separate machine-theoretic obstruction remains inside the adopted route.
\end{theorem}

\begin{corollary}[What the full machine-theoretic closure surface already buys]
Later revisions no longer need to cite the whole local hardening ladder when the only point is that the present computational reading has been conditionally closed all the way up the hierarchy along a chosen certified route. They may instead cite \(\mathfrak M^{\mathrm{full}}(\mathcal R)=1\), while still keeping explicit which reference core, which tape-family route, and which upper certificates are meant. This compresses the machine-theoretic reading to a single surface without pretending that the route has become unconditional or model-independent.
\end{corollary}




\subsection*{19AG. First machine-theoretic status decoder below and above the full closure surface}
After the full machine-theoretic closure surface has been bundled, the text should not force later revisions to reopen the entire hierarchy merely to say which machine-theoretic level is currently justified under a given certificate package. A small status decoder is therefore useful: it reads off the strongest currently justified strength level from the available closure surfaces and upgrade certificates without changing the underlying mathematics.

\begin{definition}[First machine-theoretic status decoder on the compressed computational line]
Fix a chosen admissible sequential reference core \(\mathcal R\). Define the machine-theoretic status decoder
\[
\mathfrak S^{\mathrm{mach}}(\mathcal R)\in\{\mathbf S_1,\mathbf S_2,\mathbf S_3(\mathcal R),\mathbf S_4(\mathcal R),\mathbf S_5\}
\]
as the strongest level among the hierarchy
\[
\mathbf S_1\longrightarrow \mathbf S_2\longrightarrow \mathbf S_3(\mathcal R)\longrightarrow \mathbf S_4(\mathcal R)\longrightarrow \mathbf S_5
\]
whose required certificate package has been verified. Concretely:
\begin{itemize}[leftmargin=2em]
  \item \(\mathfrak S^{\mathrm{mach}}(\mathcal R)=\mathbf S_5\) if \(\mathfrak M^{\mathrm{full}}(\mathcal R)=1\);
  \item \(\mathfrak S^{\mathrm{mach}}(\mathcal R)=\mathbf S_4(\mathcal R)\) if the tape-family upgrade route is certified but the strict upper certificate package is not fully certified;
  \item \(\mathfrak S^{\mathrm{mach}}(\mathcal R)=\mathbf S_3(\mathcal R)\) if the conditional reference-core closure surface is certified but the tape-family upgrade route is not fully certified;
  \item \(\mathfrak S^{\mathrm{mach}}(\mathcal R)=\mathbf S_2\) if the conditional primitive Zuse-style core is certified but reference-core closure is not;
  \item \(\mathfrak S^{\mathrm{mach}}(\mathcal R)=\mathbf S_1\) otherwise.
\end{itemize}
\end{definition}

\begin{proposition}[Monotone exactness of the first machine-theoretic status decoder]
For every chosen admissible sequential reference core \(\mathcal R\), the decoder \(\mathfrak S^{\mathrm{mach}}(\mathcal R)\) is monotone with respect to certificate accumulation and exact with respect to the present hierarchy: it never outputs a higher level than has actually been certified, and whenever a full certificate package for some level in the hierarchy has been established, the decoder outputs at least that level. In particular,
\[
\mathfrak M^{\mathrm{full}}(\mathcal R)=1
\quad\Longrightarrow\quad
\mathfrak S^{\mathrm{mach}}(\mathcal R)=\mathbf S_5.
\]
Likewise, if only the reference-core closure surface is available, the decoder halts at \(\mathbf S_3(\mathcal R)\), and if the tape-family upgrade route has been certified but the strict upper package has not, it halts at \(\mathbf S_4(\mathcal R)\).
\end{proposition}

\begin{corollary}[What the first machine-theoretic status decoder already buys]
Later revisions can now state the strongest presently justified machine-theoretic claim without reopening the whole ladder of lower certificates. They may cite the single decoded status \(\mathfrak S^{\mathrm{mach}}(\mathcal R)\) and then only unfold the route-specific certificate package when a stronger or more unconditional claim is intended. This keeps the computational reading disciplined: the text can speak precisely about the current level already reached, while preventing accidental overstatement about higher levels that are not yet certified.

\end{corollary}


\subsection*{19AH. Conditional/unconditional status surface on the compressed computational line}
After the machine-theoretic status decoder has been introduced, the current release line should not remain described only in abstract terms. A small present-line status surface is therefore useful: it records, in one place, the strongest route-dependent status already certified on the present v\docversion\ release line and separates it from the weaker certificate-free status that can be asserted without invoking any chosen reference-core route.

\begin{definition}[First present-line conditional/unconditional status surface on the compressed computational line]
Fix a chosen admissible sequential reference core \(\mathcal R\). Define the present-line machine-theoretic status pair
\[
\mathfrak P^{\mathrm{mach}}(\mathcal R)=\bigl(\mathfrak S^{\mathrm{abs}},\mathfrak S^{\mathrm{cond}}(\mathcal R)\bigr)
\]
by the following rule:
\begin{itemize}[leftmargin=2em]
  \item \(\mathfrak S^{\mathrm{abs}}\) is the strongest machine-theoretic level justified without invoking any chosen reference core, tape-family route, or upper certificate package;
  \item \(\mathfrak S^{\mathrm{cond}}(\mathcal R)\) is the strongest machine-theoretic level justified under the explicit chosen route \(\mathcal R\) together with the corresponding lower and upper certificate package.
\end{itemize}
The pair \(\mathfrak P^{\mathrm{mach}}(\mathcal R)\) will be called the present-line conditional/unconditional status surface.
\end{definition}

\begin{proposition}[First explicit present-line status readout on the compressed computational line]
On the present v\docversion\ release line, and for every chosen admissible sequential reference core \(\mathcal R\) along which the bundled closure surface \(\mathfrak M^{\mathrm{full}}(\mathcal R)=1\) has been certified, the present-line status surface reads
\[
\mathfrak P^{\mathrm{mach}}(\mathcal R)=\bigl(\mathbf S_1,\mathbf S_5\bigr).
\]
Equivalently: unconditionally, the present line should still only be advertised at the baseline disciplined computational-reading level \(\mathbf S_1\); conditionally, along the explicitly certified route \(\mathcal R\), the present line reaches the top certified level \(\mathbf S_5\). In particular, the present text should avoid collapsing these two readings into a single unconditional global claim.
\end{proposition}

\begin{corollary}[What the first present-line conditional/unconditional status surface already buys]
Later revisions can now cite the present machine-theoretic state of the Companion in one disciplined sentence: the line is unconditionally at \(\mathbf S_1\), while conditionally along a chosen certified route it reaches \(\mathbf S_5\). This gives a compact status formula for future hardening passes while preventing accidental overstatement about certificate-dependent closures as if they were already unconditional.
\end{corollary}

\begin{remark}[One-paragraph synopsis of the machine-theoretic line]
The computational line of the present framework can be read in three nested stages. First, the local operational layer yields a conditional primitive machine core, in which tape-like support, control, branching, and loop/return behavior are structurally separated and locally executable. Second, this primitive core can be hardened into a reference-core closure relative to a chosen sequential reference kernel \(\mathcal R\), so that the line is no longer only machine-like in analogy, but carries a disciplined reference computation surface. Third, along a certified tape-based route above this reference-core closure, the framework admits a conditional upgrade toward stronger machine-theoretic levels, up to a Turing-style completeness claim only under explicit higher certificates; accordingly, the present line supports an unconditional computational base reading, while stronger universality claims remain route-dependent and conditional.
\end{remark}




\subsection*{19AH$'$. First bidirectional observer-side machine-shadow split below the status surface}
After the present-line conditional/unconditional status surface has been fixed, one more local question becomes unavoidable on the working continuation line: how should the currently retained public machine readout itself be decomposed at observer side? The next useful step is not yet a full new universality statement, but a disciplined split of the retained machine shadow into one symmetric total channel and one antisymmetric imbalance channel. This keeps the older no-second-carrier discipline visible while making the current machine readout auditable in a form that can be compared directly with later transport, spectral, or shell-level hardening.

\begin{definition}[Bidirectional observer-side machine-shadow coordinates]
Fix the retained overlap window \(W_{\mathrm{ret}}\) and a tested lag \(\ell\). Let \(\eta_{\mathrm{drive}}^{(\ell)}\) denote the fitted observer-side drive contribution and \(\eta_{\mathrm{carrier}}^{(\ell)}\) the fitted carrier contribution on \(W_{\mathrm{ret}}\). Define the bidirectional observer-side coordinates
\[
\eta_+^{(\ell)}=\eta_{\mathrm{drive}}^{(\ell)}+\eta_{\mathrm{carrier}}^{(\ell)},
\qquad
\eta_-^{(\ell)}=\eta_{\mathrm{drive}}^{(\ell)}-\eta_{\mathrm{carrier}}^{(\ell)}.
\]
We call \(\eta_+^{(\ell)}\) the \emph{total machine-shadow channel} and \(\eta_-^{(\ell)}\) the \emph{observer-side imbalance channel}.
\end{definition}

\begin{proposition}[Current retained-overlap checkpoint at the public optimum]
On the retained overlap window currently kept on the working line, the bidirectional observer-side coordinates reproduce the present public optimum at lag
\[
\ell_*=-14
\]
with weighted defect approximately \(0.6746\) over \(42\) overlap points.
\end{proposition}

\begin{remark}[Interpretation of the checkpoint]
The point of this checkpoint is not yet a final dynamical theorem about the origin of the machine readout. It is the weaker but already useful fact that the presently best public lag can be read consistently through one bidirectional observer-side split rather than through two unrelated fitted channels. In that sense, the present step converts one numerical optimum into a structurally readable local machine-shadow record.
\end{remark}

\begin{proposition}[The imbalance channel is substantial but not primitive]
On the same retained overlap window at the public optimum \(\ell_*=-14\), the observer-side imbalance channel satisfies
\[
\frac{\lVert \eta_-^{(\ell_*)}\rVert_1}{\lVert \eta_+^{(\ell_*)}\rVert_1}\approx 0.6827,
\qquad
\operatorname{corr}\!\bigl(\eta_+^{(\ell_*)},\eta_-^{(\ell_*)}\bigr)\approx 0.9532.
\]
Hence the imbalance channel is not negligible, but it remains tightly slaved to the same total machine-shadow profile.
\end{proposition}

\begin{corollary}[First no-second-primitive-carrier reading from the machine split]
The current bidirectional observer-side readout does not force the introduction of a second primitive carrier. What it supports is one dominant machine shadow together with a substantial but highly correlated imbalance readout. In particular, the current continuation line remains compatible with the older no-second-carrier discipline: the extra visible degree of freedom appears at observer side as a locked asymmetry channel, not as evidence for a new independent carrier species.
\end{corollary}


\begin{remark}[Why this matters for the next hardening step]
The split creates a useful next leverage point. Future hardening no longer needs to ask only whether a retained carrier exists; it can ask whether a proposed theorem, transport step, or spectral proxy controls the symmetric total channel \(\eta_+\), the locked imbalance channel \(\eta_-\), or both at once. That distinction should make it easier to separate genuine carrier obligations from observer-shell correction obligations in the next working passes.
\end{remark}

\subsection*{19AH$''$. First channel-control dictionary induced by the machine-shadow split}
Once the observer-side machine readout has been split into \(\eta_+\) and \(\eta_-\), the next useful step is a control dictionary telling later arguments what kind of burden they actually carry. This is still weaker than a new dynamical theorem. Its role is bookkeeping with mathematical consequences: it distinguishes propositions that genuinely control the retained machine shadow from propositions that merely reorganize the observer-side asymmetry readout around that shadow.

\begin{definition}[\(\eta_+\)-control, \(\eta_-\)-control, and joint channel control]
Fix the retained overlap window \(W_{\mathrm{ret}}\) at the current public optimum \(\ell_*=-14\). We say that a theorem, estimate, transport rule, or spectral proxy on the present working line has
\begin{enumerate}[leftmargin=2em]
  \item \emph{total-channel control} if it bounds, reconstructs, or rigidifies \(\eta_+^{(\ell_*)}\) up to the standing admissible equivalences without requiring an independent primitive contribution beyond the present drive/carrier split;
  \item \emph{imbalance-channel control} if, once the total channel is fixed in that admissible sense, it bounds, reconstructs, or rigidifies \(\eta_-^{(\ell_*)}\) as an observer-side asymmetry readout around the same total machine shadow;
  \item \emph{joint channel control} if it controls both channels at once and in addition preserves or explains a structural lock between them, such as the present norm ratio or the present high correlation.
\end{enumerate}
\end{definition}

\begin{proposition}[First machine-side reading discipline induced by the dictionary]
On the current working line, arguments aimed at the retained carrier or at the existence of one common machine shadow should be read as \(\eta_+\)-control problems, whereas observer-shell, asymmetry, phase-correction, or frame-flicker refinements should be read as \(\eta_-\)-control problems unless they also alter the total channel itself. Mixed claims that try to upgrade both at once should be treated as joint-channel problems and therefore carry the heaviest burden of proof.
\end{proposition}

\begin{corollary}[First reformulation of the no-second-carrier discipline at machine level]
A future refinement is compatible with the present no-second-carrier discipline whenever its extra visible freedom can be organized as \(\eta_-\)-control over a fixed or jointly controlled \(\eta_+\)-channel, rather than by introducing a second primitive contribution already at the total-channel level. In particular, the present checkpoint pushes the burden of any would-be second-carrier claim upward: such a claim would now have to break the current reading at the level of \(\eta_+\), not merely by producing additional observer-side imbalance structure.
\end{corollary}

\begin{remark}[Why the dictionary is a real compression tool]
The dictionary turns one vague question --- ``does this next theorem improve the machine line?'' --- into a sharper trichotomy. A proposed step may improve the total machine shadow, the locked imbalance readout, or the coupling data between them. This should make later proof compression cleaner because many older optical and shell-side corrections can now be classified immediately as \(\eta_-\)-side obligations, while the rarer genuinely carrier-bearing steps can be isolated as \(\eta_+\)-side obligations.
\end{remark}

\subsection*{19AH$'''$. First backward channel reading of the older optical and shell-side baselines}
The dictionary becomes substantially more useful once it reorganizes not only future arguments but also the already stabilized comparison/shell archive. The present working line therefore records a first backward reading of the older optical-comparison, shortlist, bombe/sieve, shell/menu/window, and flicker-style material through the split into \(\eta_+\) and \(\eta_-\).

\begin{proposition}[First backward channel reading of the retained public baselines]
In this inherited public-release baseline block, the older public baselines should be read according to the following burden split.
\begin{enumerate}[leftmargin=2em]
  \item The retained-carrier comparison layer, public optimum selection, no-second-carrier discipline, and every comparison/back-reading/bombe/sieve step insofar as it isolates one common machine-bearing shadow are primarily \(\eta_+\)-control statements.
  \item The shell/menu/window refinements, phase-correction re-readings, and frame-flicker-style observer corrections are primarily \(\eta_-\)-control statements whenever they refine how that common shadow appears on the observer side without changing the retained optimum itself.
  \item Any statement claiming at once that a shell-side correction preserves the retained optimum, explains the asymmetry pattern, and leaves the common machine shadow identifiable should be treated as a joint-channel statement.
\end{enumerate}
\end{proposition}

\begin{corollary}[First re-reading of the optical approximation baseline]
The earlier first retained-carrier optical approximation should presently be read as one \(\eta_+\)-bearing baseline together with a subordinate \(\eta_-\)-side correction envelope. Accordingly, visible mismatch inside that optical approximation is not yet evidence against the retained-carrier discipline unless it forces an additional primitive contribution already at the total-channel level.
\end{corollary}

\begin{remark}[Compression payoff of the backward reading]
The backward reading turns a long mixed archive into a more ordered burden ledger. Carrier isolation, common-window retention, and no-second-carrier arguments can now be cited as total-channel obligations, while shell refinements, asymmetry diagnostics, and flicker-style observer corrections can be cited as imbalance-side obligations. Only the rarer steps that claim both retention and corrected observer-side readout need to reopen the full joint burden.
\end{remark}


\subsection*{19AH$^{(4)}$. First theorematic burden split for the remaining continuation knots}
Once the older baselines have been re-read through \(\eta_+\) and \(\eta_-\), the next practical question is which unresolved continuation knots still carry genuine carrier burden and which ones are only observer-side refinement burdens. The present working line therefore records one first theorematic burden split: unresolved shell-, phase-, or flicker-side refinement is not automatically a retained-carrier obstruction.

\begin{proposition}[First theorematic burden split for the remaining continuation knots]
On the present continuation line, the remaining theorematic knots should be read according to the following burden discipline.
\begin{enumerate}[leftmargin=2em]
  \item Any future statement whose conclusion asserts preservation of the retained public optimum, one-common-machine-shadow identification, or the no-second-primitive-carrier discipline must control the total channel \(\eta_+\). Purely imbalance-side refinements on \(\eta_-\) can neither prove nor refute such a carrier-level statement by themselves.
  \item The frame-flicker ansatz, shell/menu/window refinements, phase-correction re-readings, and analogous observer-side asymmetry refinements are to be treated as \(\eta_-\)-side knots as long as they only change how the retained machine shadow is seen and do not force a new total-channel profile.
  \item Any theorem claiming simultaneously that an observer-side correction explains the asymmetry pattern, improves the retained fit, and preserves one common carrier must be treated as a joint-channel knot: it has to control both \(\eta_+\) and \(\eta_-\).
  \item Consequently, a visible residual defect that can be confined to \(\eta_-\) is not yet a retained-carrier failure of the present line. A genuine carrier-level obstruction appears only if the retained optimum, the one-shadow reading, or the no-second-carrier discipline is already forced to fail on the total-channel side.
\end{enumerate}
\end{proposition}

\begin{corollary}[Default downgrade rule for unresolved observer-side mismatch]
On the present continuation line, an unresolved optical, shell-side, or flicker-style mismatch should by default be downgraded to an \(\eta_-\)-side obligation unless it demonstrably changes the retained optimum, destroys the one-shadow reading at the total-channel level, or forces an additional primitive contribution already in \(\eta_+\). This keeps the current no-second-carrier discipline from being overburdened by observer-side correction tasks that have not yet crossed into carrier space.
\end{corollary}

\begin{remark}[Why this sharpened burden split matters]
The practical gain is a much cleaner proof queue. Future hardening can now ask first whether an open point is really about the common retained shadow or only about its observer-side distortion. If it is only the latter, then it belongs to the asymmetry shell and should not delay carrier-level statements. Only the comparatively rare steps that touch both channels at once need the full coupled burden again.
\end{remark}

\subsection*{19AH$^{(5)}$. First formal demotion criterion from joint-channel burden to pure imbalance-side burden}
The burden split of 19AH$^{(4)}$ becomes operational only once one can say when a candidate continuation step no longer deserves to be treated as a genuinely joint-channel knot. The next local tool is therefore a formal demotion criterion: a pending observer-side refinement may be demoted from joint-channel burden to a pure \(\eta_-\)-side burden as soon as its unresolved remainder provably leaves the retained total profile untouched.

\begin{definition}[Observer-side demotion certificate]
Fix a pending continuation statement \(\mathcal C\) whose unresolved part currently concerns shell, phase, flicker, menu, or analogous observer-side refinement. We say that \(\mathcal C\) admits an \emph{observer-side demotion certificate} if all three of the following conditions hold on the retained admissible window under discussion.
\begin{enumerate}[leftmargin=2em]
  \item \textbf{Total-profile invariance.} The retained public optimum, the one-machine-shadow identification, and the no-second-primitive-carrier reading remain unchanged on the \(\eta_+\)-side under the unresolved remainder of \(\mathcal C\).
  \item \textbf{No forced new primitive contribution.} The unresolved remainder of \(\mathcal C\) does not require an additional primitive carrier already at the total-channel level; equivalently, the outstanding defect can still be represented without altering the retained \(\eta_+\)-carrier architecture.
  \item \textbf{Imbalance localization.} The visible mismatch produced by the unresolved remainder of \(\mathcal C\) can be localized to the observer-side imbalance channel \(\eta_-\), meaning that its effect is read as a change of asymmetry, shell, or observer-side correction profile rather than as a change of the retained common total shadow.
\end{enumerate}
\end{definition}

\begin{proposition}[First formal demotion criterion]
If a pending continuation statement \(\mathcal C\) admits an observer-side demotion certificate, then \(\mathcal C\) may be demoted from a provisional joint-channel knot to a pure \(\eta_-\)-side obligation on the present continuation line. In particular, the unresolved remainder of \(\mathcal C\) may no longer be cited as an open carrier-level obstruction against the retained public optimum, the one-shadow reading, or the no-second-primitive-carrier discipline.
\end{proposition}

\begin{proof}
The three certificate clauses remove exactly the reasons why \(\mathcal C\) would still need joint status. Clause~(1) says that the retained total profile already survives unchanged at the \(\eta_+\)-level. Clause~(2) rules out the only remaining carrier-level escape route, namely that the unresolved part secretly forces an additional primitive contribution already in the total channel. Clause~(3) places the visible remainder entirely on the observer-side imbalance ledger. Once these three facts are available simultaneously, the outstanding burden of \(\mathcal C\) no longer concerns whether the present line carries one retained common machine shadow, but only how the observer-side asymmetry shell should be refined. The knot is therefore genuinely \(\eta_-\)-side and may be demoted accordingly.
\end{proof}

\begin{corollary}[Safe demotion rule for future continuation work]
Future continuation steps should only be kept on the joint-channel ledger if at least one of the three certificate clauses fails or is not yet known. Conversely, once the retained \(\eta_+\)-profile is stable, no extra primitive carrier is forced there, and the remaining defect is localized to \(\eta_-\), the step should be queued as observer-side hardening rather than as unfinished carrier proof.
\end{corollary}

\begin{remark}[Why this tool is useful]
This criterion turns the burden split into a reusable proof tool. It prevents later drafts from repeatedly reopening the full carrier question whenever a shell-side mismatch or flicker-style correction is still under discussion. The proof queue can now be triaged by certificate: unresolved points without \(\eta_+\)-damage stay on the asymmetry ledger, while only genuine total-channel risk keeps a knot on the joint ledger.
\end{remark}


\subsection*{19AH$^{(6)}$. First provisional demotion ledger for inherited observer-side continuation knots}
\noindent\textbf{Reader cue.} Read the next two subsections as a queue-classification pair. Their purpose is to sort inherited shell/phase/flicker-style leftovers first onto the observer-side ledger by default and then to state the exact evidence that would force those same knots back into joint or total-channel burden.

The formal demotion criterion of 19AH$^{(5)}$ becomes most useful once it is applied back to the inherited queue. The point is not to declare old open points ``solved'' by fiat, but to state explicitly which of them should from now on be carried on the observer-side ledger unless and until they actually damage the retained total profile. The next local step is therefore a provisional demotion ledger for the main inherited observer-side knot families.

\begin{definition}[Provisional demotion ledger for inherited observer-side knots]
On the present continuation line, the following inherited knot families are placed on the \emph{provisional observer-side demotion ledger}, subject to re-promotion only if later work shows a genuine failure of at least one certificate clause from 19AH$^{(5)}$:
\begin{enumerate}[leftmargin=2em]
  \item shell/menu/window-refinement mismatches that alter the observer-side shell geometry without changing the retained total optimum profile,
  \item phase-correction and frame-flicker residuals that modify how the retained profile is read from the observer side while leaving the common total shadow intact,
  \item optical-approximation residuals whose visible discrepancy can still be localized to the asymmetry envelope rather than to the retained common carrier, and
  \item analogous observer-side refinements inherited from the older archive whenever they do not force a second primitive carrier and do not move the retained public optimum on the total-channel side.
\end{enumerate}
\end{definition}

\begin{proposition}[First provisional demotion ledger theorem]
The inherited knot families listed above belong, by default, on the observer-side demotion ledger rather than on the joint-channel obstruction ledger. More precisely, on the present compressed computational line they should be read as provisional \(\eta_-\)-side obligations unless later analysis shows that one of them changes the retained \(\eta_+\)-profile, breaks the one-shadow reading at total-channel level, or forces an additional primitive carrier already in \(\eta_+\).
\end{proposition}

\begin{proof}
Each listed family is inherited precisely from the shell/phase/flicker/optical side of the archive, hence from the side that 19AH$'''$ and 19AH$^{(4)}$ already identified as the natural observer-side burden domain. The formal criterion of 19AH$^{(5)}$ then supplies the missing operational rule: as long as the retained total profile remains fixed, no new primitive carrier is forced in the total channel, and the visible mismatch stays localizable to the imbalance channel, the corresponding inherited knot is not entitled to remain on the joint ledger. The listed families therefore enter the queue under default demotion. Re-promotion is reserved only for the exceptional case that later work actually detects \(\eta_+\)-damage or a genuine carrier-level escape route.
\end{proof}

\begin{corollary}[How to read the old queue from now on]
From this point onward, older shell/menu/window, phase-correction, flicker-style, and observer-side optical residuals should no longer be cited as open carrier objections by default. They remain real work items, but they are to be read first as asymmetry-hardening tasks. Only after a concrete failure of the demotion certificate has been exhibited should one move such a knot back to the coupled or carrier-level queue.
\end{corollary}

\begin{remark}[What ``demotion'' means here in plain language]
This is only a bookkeeping and proof-burden move. ``Demotion'' does \emph{not} mean that the point is unimportant, false, or removed. It means: this point still matters, but at the current stage it counts as a problem of observer-side correction rather than as evidence that the common retained carrier picture has broken. In ordinary language, we stop treating every leftover wobble in the readout shell as if it were a collapse of the whole machine-bearing line.
\end{remark}

\subsection*{19AH$^{(7)}$. First formal re-promotion criterion from observer-side demotion back to joint or total burden}
Once a provisional demotion ledger exists, the opposite discipline becomes equally important. Demotion is safe only if later work can also say when a demoted knot must be upgraded again. The next local tool is therefore a formal re-promotion criterion: if a supposedly observer-side remainder turns out to damage the retained total machine-shadow profile or to force a new primitive carrier already at total-channel level, then the knot must leave the pure \(\eta_-\)-ledger and return to a higher burden class.

\begin{definition}[Observer-side re-promotion trigger]
Let a continuation knot family \(K\) currently lie on the provisional observer-side demotion ledger. We say that \(K\) carries an \emph{observer-side re-promotion trigger} if later analysis exhibits at least one of the following three failures:
\begin{enumerate}[leftmargin=2em]
  \item the retained total machine-shadow profile \(\eta_+\) is no longer fixed up to the standing admissible equivalences;
  \item the one-shadow reading at total-channel level breaks, in the sense that the visible fit can no longer be organized without introducing an additional primitive carrier already in \(\eta_+\);
  \item the observer-side remainder can no longer be localized purely to \(\eta_-\), because its correction necessarily changes the total-channel optimum, the total-channel rigidity data, or the total-channel carrier count.
\end{enumerate}
\end{definition}

\begin{proposition}[First re-promotion criterion]
Let \(K\) be a knot family on the provisional observer-side demotion ledger. If \(K\) satisfies an observer-side re-promotion trigger, then \(K\) is no longer admissible as a pure \(\eta_-\)-side obligation. More precisely:
\begin{enumerate}[leftmargin=2em]
  \item if the trigger only shows that \(K\) cannot be separated from the total-channel data but does not yet force a new primitive carrier, then \(K\) must be re-promoted from pure \(\eta_-\)-burden to \emph{joint-channel burden};
  \item if the trigger forces a change in the retained total machine-shadow profile or compels an additional primitive carrier already in \(\eta_+\), then \(K\) must be re-promoted all the way to \emph{total-channel burden}.
\end{enumerate}
\end{proposition}

\begin{proof}
This is the exact converse of the demotion logic from 19AH$^{(5)}$ and the ledger rule from 19AH$^{(6)}$. A knot may remain on the pure observer-side ledger only while the retained total profile stays fixed, no second primitive carrier is forced at total-channel level, and the visible remainder remains localizable to the imbalance channel. The three trigger clauses record the failure of precisely those permissions. Once a clause fails, pure \(\eta_-\)-treatment is no longer justified. If the failure merely entangles the remainder with the total channel, the knot becomes joint. If the failure genuinely alters or multiplies the total-channel carrier structure, the knot becomes a total-channel problem.
\end{proof}

\begin{corollary}[Plain-language queue rule]
A leftover shell, phase, flicker, or optical residual may stay in the ``observer-side corrections'' queue only while it leaves the main common machine shadow intact. The moment it starts moving that main shadow, or starts behaving as if a second primitive carrier were needed in the total profile, it must be upgraded again. In ordinary language: demotion is reversible, but only by evidence of real damage to the main channel.
\end{corollary}

\begin{remark}[Why the demotion/re-promotion pair is useful]
Together, 19AH$^{(5)}$--19AH$^{(8)}$ now give a two-way burden discipline. Demotion prevents every observer-side wobble from being misread as a carrier collapse; re-promotion prevents the opposite mistake of hiding a true carrier-level defect inside the observer-side queue. This makes the inherited continuation landscape auditable in both directions.
\end{remark}

\subsection*{19AH$^{(8)}$. First channel traffic-light ledger induced by the demotion/re-promotion discipline}
\noindent\textbf{Reader cue.} Read the next ledger as a dashboard. Its purpose is to summarize the already stated demotion and re-promotion rules in green/yellow/red form, so that the inherited queue can be scanned quickly before later live burdens are isolated.

The next local consolidation step is to turn the two-way burden discipline into a directly readable queue picture. Instead of carrying every inherited continuation knot under one undifferentiated status word, one may sort the inherited families by a traffic-light rule. Green means that the family presently remains admissible as a pure observer-side burden; yellow means that the family is already entangled enough with the total channel that it must be treated jointly; red is reserved for the genuinely exceptional case that the family damages the retained total profile itself or forces a second primitive carrier already in \(\eta_+\).

\begin{definition}[Channel traffic-light ledger]
A \emph{channel traffic-light ledger} for inherited continuation knots is an assignment of each inherited knot family \(K\) to exactly one of the following three burden classes:
\begin{enumerate}[leftmargin=2em]
  \item \textbf{green / observer-side only}: \(K\) satisfies the demotion certificate from 19AH$^{(5)}$ and no re-promotion trigger from 19AH$^{(7)}$ is presently visible;
  \item \textbf{yellow / joint-channel}: \(K\) no longer qualifies for pure observer-side treatment, but the available evidence still stops short of forcing a genuine total-channel carrier change;
  \item \textbf{red / total-channel}: \(K\) damages the retained total profile \(\eta_+\) itself or forces an additional primitive carrier already at total-channel level.
\end{enumerate}
\end{definition}

\begin{proposition}[First provisional traffic-light placement for the inherited queue]
On the present compressed computational line, the inherited queue may provisionally be sorted as follows.
\begin{enumerate}[leftmargin=2em]
  \item \textbf{Green:} shell/menu/window refinements, phase-correction refinements, frame-flicker compensation refinements, and observer-side optical residual refinements remain on the green ledger, because they presently leave the retained total machine-shadow profile fixed, force no second primitive carrier in \(\eta_+\), and stay localizable to imbalance-side repair.
  \item \textbf{Yellow:} any inherited knot whose proposed correction already shifts the retained optimum location, modifies the rigidity data needed to keep the one-shadow total reading stable, or otherwise couples the observer-side mismatch back into the total-channel fit must enter the yellow ledger and be treated as a joint-channel obligation.
  \item \textbf{Red:} no presently inherited knot family is forced onto the red ledger on the current line; red is reserved only for the case that later work exhibits an actual breakdown of the retained total profile or a genuine second primitive carrier already in \(\eta_+\).
\end{enumerate}
\end{proposition}

\begin{proof}
The green placement is exactly the provisional demotion ledger from 19AH$^{(6)}$, now read through the demotion certificate of 19AH$^{(5)}$. The yellow placement is the first nontrivial output of the re-promotion rule from 19AH$^{(7)}$: once a knot cannot be repaired purely on the imbalance side, yet still does not force a total-channel carrier split, it must be carried jointly. The red placement is currently empty because the present line has not yet produced evidence that the retained total machine shadow itself fails or that a second primitive carrier is already required in the total channel. Hence the inherited queue is presently green-dominant, with yellow reserved for future entanglement cases and red reserved for genuine carrier-level failure.
\end{proof}

\begin{corollary}[Plain-language traffic-light reading]
At the moment, most inherited leftovers are green: they are still real work, but only on the observer side. Yellow means: this leftover is beginning to touch the main common shadow and therefore needs joint treatment. Red would mean: the main common shadow itself is no longer enough. On the current line, we are not at red.
\end{corollary}

\begin{remark}[Why this corrects possible drift]
The point of the traffic-light ledger is precisely to stop the continuation block from drifting into an ever longer meta-language. Instead of adding more abstract burden labels, the present line now has a single concrete reading device: green for pure \(\eta_-\)-repair, yellow for genuine coupling, red for actual \(\eta_+\)-damage. This keeps the new diagnostic layer operational and tied to concrete queue management.
\end{remark}

\subsection*{19AH$^{(9)}$. First synchronized OP ledger after the channel-discipline consolidation}
\noindent\textbf{Reader cue.} Read the next synchronized OP ledger as a theorem-summary device. Its purpose is to collapse historical OP labels into \emph{closed}, \emph{absorbed}, and \emph{residual-core} status before the roadmap is read, not to reopen already discharged fronts.

Once the inherited continuation queue has been sorted by total/imbalance burden and by the traffic-light rule, the older named OP fronts should no longer be carried unchanged merely for historical convenience. The next consolidation step is therefore to synchronize the inherited OP labels against the presently certified structure itself. The point is not to rename difficulty away, but to state explicitly which earlier OP fronts are already structurally discharged, which no longer survive as independent fronts, and which remain as the genuine residual core for the next version.

\begin{definition}[Synchronized OP ledger]
A \emph{synchronized OP ledger} on the present compressed computational line assigns every inherited named OP front \(P\) to exactly one of the following three statuses:
\begin{enumerate}[leftmargin=2em]
  \item \textbf{closed}: the relevant structural burden of \(P\) is already discharged on the present line and should no longer be listed as a primary open problem;
  \item \textbf{absorbed}: \(P\) is no longer independent, because its only remaining unresolved content is already contained in another residual core;
  \item \textbf{residual core}: \(P\) still carries genuinely new proof burden that is not yet discharged by the present line.
\end{enumerate}
\end{definition}

\begin{proposition}[First synchronized placement of the inherited OP queue]
On the present compressed computational line, the inherited named OP queue may be synchronized as follows.
\begin{enumerate}[leftmargin=2em]
  \item \textbf{Closed:} OP~3 is closed as an independent front by the shell-certificate/residual-interface line already exported in the present text, and OP~5 Step~(A) is closed by the source-return, branchwise, and routed carrier-first closure package.
  \item \textbf{Absorbed:} OP~2 is absorbed into OP~1, because the remaining frame-offset/continuum-correction burden no longer survives as an independent Dirac-side front once the present line is synchronized; OP~4 is absorbed into OP~1, because its only still-open quantitative content is precisely the same continuum-limit correction \(\delta_{\mathrm{CL}}\).
  \item \textbf{Residual core:} the live queue is reduced to OP~1 and OP~5 Step~(B). OP~1 remains the continuum-limit/total-scaling residual centered on \(\delta_{\mathrm{CL}}\). OP~5 Step~(B) remains the exact \(\Theta\)--\(R\) balance problem on the routed photon-frame side.
\end{enumerate}
\end{proposition}

\begin{proof}
The earlier OP~3 burden has already been converted into a shell-certified interface theorem together with reusable residual/shell reading rules, so it no longer behaves as an unreduced primary front. Likewise, the source-return closure theorem, its branchwise refinement, and the routed carrier-first OP5 interface already discharge the source/return and admissible-routing part of OP~5, leaving only the exact total balance problem as live burden. Earlier in the text, OP~4 is explicitly reduced to the same continuum-limit correction \(\delta_{\mathrm{CL}}\) as OP~1, and the remaining OP~2 frame-offset burden has already collapsed onto that same continuum-limit core rather than surviving as an independent front. Hence the present synchronized ledger has only two genuine residual cores left: OP~1 and OP~5 Step~(B).
\end{proof}

\begin{corollary}[Plain-language OP synchronization]
The inherited OP queue is no longer a five-front queue. On the current line, OP~3 is closed, OP~2 and OP~4 are absorbed into OP~1, OP~5 survives only through Step~(B), and the live queue is therefore a two-core queue: OP~1 and OP~5 Step~(B).
\end{corollary}

\begin{remark}[Why this synchronization matters]
This synchronization keeps the continuation line honest in both directions. It prevents already discharged or absorbed fronts from being carried forever as if nothing had changed, but it also refuses to pretend that the two genuine residual cores are already gone. The next real version target is therefore sharply focused rather than historically bloated.
\end{remark}

\subsection*{19AH$^{(10)}$. Two-core OP-closing roadmap before the later freeze deferral}
\noindent\textbf{Reader cue.} Read the next roadmap as an earlier compression snapshot of the line. Its purpose is to orient burden genealogy by recording how the queue had once been narrowed to two cores before the later freeze deferral; it does not override the present post-freeze priority order.

After the synchronized ledger from 19AH$^{(9)}$, the line first compressed the next-version target to a narrow two-core closure program. At that stage, the intended next version was to try to close the remaining queue without reopening the green observer-side ledger except where a genuine yellow or red trigger is produced by the re-promotion rule. The purpose is therefore not to add another diagnostic family, but to turn the present compression into two direct closure attempts: one on the continuum-limit residual and one on the exact routed photon-frame balance.

\begin{definition}[Residual-core closing roadmap]
A \emph{residual-core closing roadmap} in this earlier two-core reading is an ordered list of closure tasks that leaves all green observer-side knot families parked by default and reopens them only if a later step actually produces a re-promotion trigger.
\end{definition}

\begin{proposition}[First two-core OP-closing roadmap before the later freeze deferral]
In that earlier two-core reading, the next version would have proceeded through the following tasks.
\begin{enumerate}[leftmargin=2em]
  \item \textbf{OP~1 isolation.} Reduce OP~1 to one final explicit continuum-limit residual identity for \(\delta_{\mathrm{CL}}\), so that no auxiliary historical formulation remains in the live queue.
  \item \textbf{OP~1 closure attempt.} Attack the vanishing or exact identification of that residual directly from the retained total-channel data and the lattice-to-continuum transport structure already built into the present line.
  \item \textbf{OP~5(B) isolation.} Rewrite OP~5 Step~(B) as one exact routed balance statement for the remaining \(\Theta\)--\(R\) defect term, with all source-return, branchwise, and carrier-first routing pieces treated as fixed infrastructure rather than reopened.
  \item \textbf{OP~5(B) closure attempt.} Prove or refute the vanishing of the remaining routed balance defect without importing a second primitive total-channel carrier.
  \item \textbf{Final synchronization pass.} If the two closure attempts succeed, update the synchronized OP ledger so that no live residual core remains; if one attempt fails, record precisely which single residual core survives and keep the green observer-side queue parked.
\end{enumerate}
\end{proposition}

\begin{corollary}[Roadmap discipline in the earlier two-core reading]
In that earlier two-core reading, shell/menu/window, phase-correction, frame-flicker, and other green observer-side families are not first targets. They may be revisited only if closing OP~1 or OP~5 Step~(B) actually forces a yellow or red re-promotion.
\end{corollary}

\begin{remark}[Why the roadmap is intentionally narrow]
The synchronized ledger shows that the continuation line is now past the stage where every historical knot needs to be touched again. The next mathematically serious step in that earlier reading was to spend the new leverage exactly where the residual burden still lives: on the continuum-limit source residual of OP~1 and on the exact routed balance residual of OP~5 Step~(B). A narrower roadmap is therefore not a loss of ambition, but a sign that the line has finally compressed enough to attack the remaining core directly.
\end{remark}




\subsection*{19AH$^{(11)}$. Two-core residual dictionary opening the v\docversion\ closure line}
After the synchronized OP ledger from 19AH$^{(9)}$ and the next-version roadmap from 19AH$^{(10)}$, the correct opening move for v\docversion\ is to stop carrying OP~1 and OP~5 Step~(B) as diffuse historical labels and to rewrite them as two explicit residual objects. Every later closure argument can then be judged by one simple rule: it either changes one of these residuals directly, or it is auxiliary infrastructure only.

\begin{definition}[Continuum-limit residual for OP~1]
\label{def:delta-cl-residual}
Define the \emph{continuum-limit residual}
\[
  \Delta_{\mathrm{CL}}
  :=
  \log \delta_{\mathrm{CL}} + \frac{\gamma_1}{4\pi\gamma_L^5}.
\]
Equivalently, \(\Delta_{\mathrm{CL}}=0\) if and only if the corrected continuum-limit formula from Remark~\ref{conj:cl-correction},
\[
  \delta_{\mathrm{CL}} = \exp\!\left(-\frac{\gamma_1}{4\pi\gamma_L^5}\right),
\]
holds exactly.
\end{definition}

\begin{definition}[Routed photon-frame balance defect for OP~5 Step~(B)]
\label{def:delta-theta-r-residual}
Let \(T:=\gamma_L^{-3/2}\) and let \(\mathcal T(T)\) be the photon-frame Theta sum from Definition~\ref{def:theta-sum-B}. Define the \emph{routed photon-frame balance defect}
\[
  \Delta_{\Theta R}(T)
  :=
  \mathcal T(T)
  - \left(\frac{3T}{2} - \log T - \frac{T^2}{24} - \gamma_E - \bsym^2 T^2\right).
\]
Equivalently, \(\Delta_{\Theta R}(T)=0\) if and only if the exact Theta-sum balance conjecture of Proposition~\ref{conj:theta-sum-B} holds at the routed photon-frame cutoff.
\end{definition}

\begin{proposition}[Two-core reduction of the live OP queue]
\label{prop:two-core-residual-dictionary}
On the synchronized ledger of 19AH$^{(9)}$, the live OP queue is reduced exactly as follows.
\begin{enumerate}[leftmargin=2em]
  \item OP~1 is equivalent to the vanishing or exact first-principles identification of the single scalar residual \(\Delta_{\mathrm{CL}}\).
  \item OP~5 Step~(B) is equivalent to the vanishing of the single routed balance defect \(\Delta_{\Theta R}(T)\).
  \item Any inherited continuation task that leaves both \(\Delta_{\mathrm{CL}}\) and \(\Delta_{\Theta R}(T)\) unchanged remains auxiliary unless it triggers a yellow or red re-promotion in the sense of 19AH$^{(7)}$--19AH$^{(8)}$.
\end{enumerate}
\end{proposition}

\begin{proof}
The OP~1 side was already compressed in theorem form by Theorem~\ref{thm:op1op4-unified}, which reduces the open quantitative content of OP~1 and OP~4 to the single continuum-limit correction \(\delta_{\mathrm{CL}}\). Writing that remaining burden as \(\Delta_{\mathrm{CL}}\) only turns the same statement into one scalar residual identity. On the OP~5 side, Step~(A) is already closed by Proposition~\ref{prop:theta-remainder-trivial}, while Step~(B) is explicitly reduced in Proposition~\ref{conj:theta-sum-B} to one exact routed photon-frame Theta-sum balance statement. Writing that remaining burden as \(\Delta_{\Theta R}(T)\) therefore changes no mathematical content; it only compresses the live queue into one explicit defect variable. The third clause is then the direct reading of the synchronized ledger 19AH$^{(9)}$, the roadmap 19AH$^{(10)}$, and the re-promotion discipline 19AH$^{(7)}$--19AH$^{(8)}$.
\end{proof}

\begin{remark}[Why the new version should start here]
The point of this reduction is not rhetorical neatness. Once the queue is written as \(\Delta_{\mathrm{CL}}\) and \(\Delta_{\Theta R}(T)\), every later step can be audited immediately: either it shrinks one of these two residuals, identifies one of them exactly, or it belongs to infrastructure/support rather than to the remaining core closure burden. This prevents the next version from drifting back into a diffuse many-front reading of the inherited archive.
\end{remark}


\subsection*{19AH$^{(12)}$. First OP~1 slope certificate at the continuum-limit residual}
After the two-core reduction from 19AH$^{(11)}$, the next honest frontal attack on OP~1 is to compress the residual once more: not into a new conjectural formula, but into one normalized attenuation index. This separates the ``what has to be derived'' question from the ``how much surrounding continuum-limit story one tells'' question. The live OP~1 burden is then no longer an undifferentiated derivation of \(\delta_{\mathrm{CL}}\), but the identification of one single scalar exponent/slope quantity.

\begin{definition}[Normalized continuum-limit attenuation index]
\label{def:xi-cl}
Define the \emph{normalized continuum-limit attenuation index}
\[
  \Xi_{\mathrm{CL}}
  :=
  -\frac{4\pi\gamma_L^5}{\gamma_1}\log \delta_{\mathrm{CL}}.
\]
Using Theorem~\ref{thm:op1op4-unified}, equivalently
\[
  \Xi_{\mathrm{CL}}
  =
  -\frac{4\pi\gamma_L^5}{\gamma_1}\log\!\bigl(\sqrt{3}\,\beta_{\mathrm{sym}}\bigr).
\]
Thus \(\Xi_{\mathrm{CL}}\) measures exactly how far the still-open continuum-limit correction deviates from the candidate exponent in Remark~\ref{conj:cl-correction}, but in normalized unit form.
\end{definition}

\begin{proposition}[Equivalent slope certificate for OP~1]
\label{prop:op1-slope-certificate}
The following statements are equivalent.
\begin{enumerate}[leftmargin=2em]
  \item The OP~1 continuum-limit residual vanishes, i.e. \(\Delta_{\mathrm{CL}}=0\).
  \item The normalized attenuation index satisfies \(\Xi_{\mathrm{CL}}=1\).
  \item Whenever one can write the remaining continuum-limit correction at the first live scale in exponential form
  \[
    \delta_{\mathrm{CL}} = e^{-c_{\mathrm{CL}}\gamma_1},
  \]
  for some dimensionless attenuation slope \(c_{\mathrm{CL}}\), that slope is exactly
  \[
    c_{\mathrm{CL}} = \frac{1}{4\pi\gamma_L^5}.
  \]
\end{enumerate}
\end{proposition}

\begin{proof}
By Definition~\ref{def:delta-cl-residual},
\[
  \Delta_{\mathrm{CL}}=0
  \quad\Longleftrightarrow\quad
  \log \delta_{\mathrm{CL}} = -\frac{\gamma_1}{4\pi\gamma_L^5}.
\]
Multiplying by \(-4\pi\gamma_L^5/\gamma_1\) gives
\[
  \Delta_{\mathrm{CL}}=0
  \quad\Longleftrightarrow\quad
  -\frac{4\pi\gamma_L^5}{\gamma_1}\log \delta_{\mathrm{CL}} = 1,
\]
which is exactly \(\Xi_{\mathrm{CL}}=1\), proving the equivalence of (1) and (2). If in addition \(\delta_{\mathrm{CL}} = e^{-c_{\mathrm{CL}}\gamma_1}\), then
\[
  \log \delta_{\mathrm{CL}} = -c_{\mathrm{CL}}\gamma_1,
\]
so the same identity becomes
\[
  c_{\mathrm{CL}}\gamma_1 = \frac{\gamma_1}{4\pi\gamma_L^5}
  \quad\Longleftrightarrow\quad
  c_{\mathrm{CL}} = \frac{1}{4\pi\gamma_L^5}.
\]
Conversely, this slope value immediately reconstructs the vanishing residual formula from Definition~\ref{def:delta-cl-residual}. Hence all three statements are equivalent.
\end{proof}

\begin{corollary}[What the r02 OP~1 attack reduces to]
\label{cor:op1-r02-target}
On the present compressed line, a successful OP~1 closure attempt no longer needs to present itself as a broad continuum-limit narrative. It is enough to derive, certify, or uniquely force the single normalized attenuation statement
\[
  \Xi_{\mathrm{CL}}=1.
\]
Equivalently, under any admissible exponential first-scale attenuation ansatz, OP~1 is closed once the slope \(c_{\mathrm{CL}}\) is shown to equal \(1/(4\pi\gamma_L^5)\).
\end{corollary}

\begin{remark}[Why this is a real reduction and not just notation]
This step still does \emph{not} prove OP~1. What it does remove is one layer of ambiguity about what the next version must actually derive. The live target is no longer a vague \emph{prime-lattice / continuum-limit correction story}; it is the single normalized exponent encoded by \(\Xi_{\mathrm{CL}}\). Any later lemma about lattice counting, prime spacing, frame conversion, or first-zero transport is now auditable by one question only: does it fix \(\Xi_{\mathrm{CL}}\), leave it unchanged, or force it away from~1? In that sense r02 turns OP~1 into the narrowest honest one-scalar closure gate currently available.
\end{remark}


\subsection*{19AH$^{(13)}$. First frame-transport split certificate for OP~1 in photon-first-zero normalization}
After the slope certificate from 19AH$^{(12)}$, the next useful compression is not to invent a new exponent but to split the already isolated OP~1 load into its frame-conversion part and its photon-side transport part. Remark~\ref{conj:cl-correction} already records that the candidate residual formula may be written in two equivalent ways,
\[
  \delta_{\mathrm{CL}}
  = \exp\!\left(-\frac{\gamma_1}{4\pi\gamma_L^5}\right)
  = \exp\!\left(-\frac{\gamma_1^{(\mathrm{photon})}}{4\pi\gamma_L^4}\right),
\]
with \(\gamma_1^{(\mathrm{photon})}=\gamma_1/\gamma_L\). The fifth power of \(\gamma_L\) in the matter-frame formula is therefore already visibly split into one frame-change factor and four photon-side transport factors. This means that OP~1 can be attacked in a sharper way: not as a mysterious global \(\gamma_L^5\)-law, but as a photon-first-zero transport law with four residual shell factors after the frame conversion has been peeled off.

\begin{definition}[Photon-normalized continuum-transport index]
\label{def:psi-cl}
Define the \emph{photon-normalized continuum-transport index}
\[
  \Psi_{\mathrm{CL}}
  :=
  -\frac{4\pi\gamma_L^4}{\gamma_1^{(\mathrm{photon})}}\log \delta_{\mathrm{CL}}.
\]
Equivalently, since \(\gamma_1^{(\mathrm{photon})}=\gamma_1/\gamma_L\),
\[
  \Psi_{\mathrm{CL}}
  =
  -\frac{4\pi\gamma_L^5}{\gamma_1}\log \delta_{\mathrm{CL}}
  = \Xi_{\mathrm{CL}}.
\]
Thus \(\Psi_{\mathrm{CL}}\) carries the same open numerical burden as \(\Xi_{\mathrm{CL}}\), but written after the matter-to-photon frame passage has been made explicit.
\end{definition}

\begin{proposition}[Frame-transport split certificate for OP~1]
\label{prop:op1-frame-transport-split}
The following statements are equivalent.
\begin{enumerate}[leftmargin=2em]
  \item The OP~1 residual vanishes, i.e. \(\Delta_{\mathrm{CL}}=0\).
  \item The slope certificate of 19AH$^{(12)}$ holds, i.e. \(\Xi_{\mathrm{CL}}=1\).
  \item The photon-normalized transport index satisfies \(\Psi_{\mathrm{CL}}=1\).
  \item The continuum-limit correction has the photon-first-zero transport form
  \[
    \delta_{\mathrm{CL}}
    = \exp\!\left(-\frac{\gamma_1^{(\mathrm{photon})}}{4\pi\gamma_L^4}\right).
  \]
\end{enumerate}
In particular, OP~1 is equivalent to proving that the entire residual attenuation beyond the frame change from \(\gamma_1\) to \(\gamma_1^{(\mathrm{photon})}\) is exactly one unit of photon-normalized four-layer transport load.
\end{proposition}

\begin{proof}
The equivalence of (1) and (2) is Proposition~\ref{prop:op1-slope-certificate}. By Definition~\ref{def:psi-cl} and the identity \(\gamma_1^{(\mathrm{photon})}=\gamma_1/\gamma_L\), one has
\[
  \Psi_{\mathrm{CL}}
  =
  -\frac{4\pi\gamma_L^4}{\gamma_1^{(\mathrm{photon})}}\log \delta_{\mathrm{CL}}
  =
  -\frac{4\pi\gamma_L^5}{\gamma_1}\log \delta_{\mathrm{CL}}
  = \Xi_{\mathrm{CL}},
\]
so (2) and (3) are equivalent. Solving \(\Psi_{\mathrm{CL}}=1\) for \(\delta_{\mathrm{CL}}\) gives
\[
  \log \delta_{\mathrm{CL}} = -\frac{\gamma_1^{(\mathrm{photon})}}{4\pi\gamma_L^4},
\]
and exponentiating yields (4). Conversely, (4) inserted into Definition~\ref{def:psi-cl} gives \(\Psi_{\mathrm{CL}}=1\). Hence all four statements are equivalent.
\end{proof}

\begin{corollary}[What the r03 OP~1 attack really changes]
\label{cor:op1-r03-target}
The remaining OP~1 burden may now be read in a more structured way. One no longer has to ``explain a \(\gamma_L^5\) exponent all at once''. It is enough to justify the following split picture:
\begin{enumerate}[leftmargin=2em]
  \item one power of \(\gamma_L\) comes from the explicit passage from the matter-frame first zero \(\gamma_1\) to the photon-frame first zero \(\gamma_1^{(\mathrm{photon})}\);
  \item the remaining four powers form the actual photon-side transport attenuation burden;
  \item that photon-side burden has normalized weight exactly one.
\end{enumerate}
Therefore any future counting, lattice, shell, or first-zero transport lemma may now be audited by the sharper question: does it explain the unit value of \(\Psi_{\mathrm{CL}}\), rather than merely restating the matter-frame exponent?
\end{corollary}

\begin{remark}[Why this is progress without pretending closure]
This still does not prove OP~1. The real gain is conceptual localization. In r02 the target was the one-scalar index \(\Xi_{\mathrm{CL}}\). In r03 the same scalar is re-read in a frame-split form that matches the standing photon/matter distinction of the Companion. That removes one source of merge-style ambiguity: the fifth power of \(\gamma_L\) no longer needs to be read as an indivisible miracle. The present line may now search specifically for a \emph{four-layer photon transport law after one explicit frame conversion}. That is narrower, structurally cleaner, and closer to the actual lattice/transport heuristics already used elsewhere in the document.
\end{remark}

\subsection*{19AH$^{(14)}$. First finite four-shell balance certificate for OP~1 after the frame split}
After the frame-transport split of 19AH$^{(13)}$, the next honest compression is to stop carrying the photon-side transport burden as one undifferentiated four-power block and to write it as a finite shell balance. This does not yet explain \emph{why} the four-layer photon transport should have normalized weight one, but it does isolate the exact finite statement that any such explanation would have to prove. In particular, OP~1 can now be re-read as a question about the total normalized load carried by four post-frame transport layers, rather than as a raw continuum-limit estimate.

\begin{definition}[Finite four-shell transport loads for the OP~1 residual]
\label{def:op1-four-shell-loads}
Assume the photon-side continuum-limit attenuation admits a four-factor decomposition
\[
  \delta_{\mathrm{CL}} = \prod_{j=1}^4 \tau_j,
\]
with positive shell factors \(\tau_j>0\). Define the corresponding \emph{normalized shell loads}
\[
  \lambda_j := -\frac{4\pi\gamma_L^4}{\gamma_1^{(\mathrm{photon})}}\log \tau_j
  \qquad (j=1,2,3,4).
\]
Whenever the above product formula holds, the photon-normalized transport index of Definition~\ref{def:psi-cl} satisfies
\[
  \Psi_{\mathrm{CL}} = \sum_{j=1}^4 \lambda_j.
\]
\end{definition}

\begin{proposition}[Finite four-shell closure certificate for OP~1]
\label{prop:op1-four-shell-certificate}
Suppose \(\delta_{\mathrm{CL}}=\prod_{j=1}^4 \tau_j\) is a positive four-shell decomposition and let \(\lambda_j\) be the normalized shell loads from Definition~\ref{def:op1-four-shell-loads}. Then the following are equivalent.
\begin{enumerate}[leftmargin=2em]
  \item The OP~1 residual vanishes, i.e. \(\Delta_{\mathrm{CL}}=0\).
  \item The photon-normalized transport index satisfies \(\Psi_{\mathrm{CL}}=1\).
  \item The four normalized shell loads satisfy the finite balance law
  \[
    \sum_{j=1}^4 \lambda_j = 1.
  \]
\end{enumerate}
In particular, after the frame split from 19AH$^{(13)}$, OP~1 is equivalent to proving that the post-frame four-shell transport carries total normalized load exactly one.
\end{proposition}

\begin{proof}
By Proposition~\ref{prop:op1-frame-transport-split}, statement (1) is equivalent to \(\Psi_{\mathrm{CL}}=1\), so (1) and (2) are equivalent. If \(\delta_{\mathrm{CL}}=\prod_{j=1}^4 \tau_j\), then
\[
  \log \delta_{\mathrm{CL}} = \sum_{j=1}^4 \log \tau_j.
\]
Multiplying by \(-4\pi\gamma_L^4/\gamma_1^{(\mathrm{photon})}\) gives
\[
  \Psi_{\mathrm{CL}}
  = -\frac{4\pi\gamma_L^4}{\gamma_1^{(\mathrm{photon})}}\log \delta_{\mathrm{CL}}
  = \sum_{j=1}^4 \left(-\frac{4\pi\gamma_L^4}{\gamma_1^{(\mathrm{photon})}}\log \tau_j\right)
  = \sum_{j=1}^4 \lambda_j,
\]
which is exactly statement (3). Hence (2) and (3) are equivalent, and all three assertions are equivalent.
\end{proof}

\begin{corollary}[Equal-shell partition is sufficient but not necessary]
\label{cor:op1-quarter-shells}
Under the hypotheses of Proposition~\ref{prop:op1-four-shell-certificate}, any of the following is sufficient to close OP~1.
\begin{enumerate}[leftmargin=2em]
  \item \(\lambda_j=\tfrac14\) for each \(j=1,2,3,4\);
  \item more generally, any asymmetric four-shell profile with
  \[
    \lambda_1+\lambda_2+\lambda_3+\lambda_4 = 1.
  \]
\end{enumerate}
Thus the present line does \emph{not} require equal quarter-load shells; equal partition is merely the simplest admissible closure pattern.
\end{corollary}



\begin{remark}[Why r04 is a genuine narrowing of OP~1]
This still does not prove OP~1, and it should not be oversold as such. The genuine gain is that OP~1 is now finite twice over. First, r02 reduced it to one scalar index; second, r03 split that index into one frame factor plus a four-layer photon burden; now r04 turns the surviving photon burden into a finite additive balance law. The live question is no longer \emph{``what is the whole continuum-limit correction story?''} but rather: \emph{which four post-frame shell loads are present, and why do they sum to one?} Any future lattice, shell, counting, or transport lemma can therefore be audited even more sharply: does it produce one of the \(\lambda_j\), constrain their sum, or leave the finite balance untouched?
\end{remark}

\subsection*{19AH$^{(15)}$. First two-pass half-burden certificate for the OP~1 four-shell balance}
The finite four-shell balance of 19AH$^{(14)}$ still leaves open where the four loads \(\lambda_j\) should come from. The first structurally natural candidate is a symmetrized two-pass reading of the post-frame photon transport: two shell loads belong to the outward/forward pass and two belong to the return/backward pass. This does not yet determine the individual shell factors, but it reduces the origin question from four unrelated loads to two pass totals. That is already enough to expose the first genuinely bidirectional closure pattern for OP~1.

\begin{definition}[Two-pass shell-pair loads for the OP~1 residual]
\label{def:op1-two-pass-loads}
Under the hypotheses of Definition~\ref{def:op1-four-shell-loads}, define the forward- and backward-pass totals
\[
  \Lambda_{\rightarrow} := \lambda_1+\lambda_2,
  \qquad
  \Lambda_{\leftarrow} := \lambda_3+\lambda_4.
\]
We call the four-shell transport profile \emph{two-pass grouped} if the post-frame shell loads are read as two forward-side loads and two backward-side loads in this way. We call it \emph{pass-symmetric} if in addition
\[
  \Lambda_{\rightarrow}=\Lambda_{\leftarrow}.
\]
\end{definition}

\begin{proposition}[Two-pass half-burden certificate for OP~1]
\label{prop:op1-two-pass-half-burden}
Suppose the OP~1 residual admits a two-pass grouped four-shell decomposition in the sense of Definition~\ref{def:op1-two-pass-loads}. Then the following are equivalent.
\begin{enumerate}[leftmargin=2em]
  \item The OP~1 residual vanishes, i.e. \(\Delta_{\mathrm{CL}}=0\).
  \item The photon-side transport index satisfies \(\Psi_{\mathrm{CL}}=1\).
  \item The two pass totals satisfy
  \[
    \Lambda_{\rightarrow}+\Lambda_{\leftarrow}=1.
  \]
\end{enumerate}
If the profile is moreover pass-symmetric, then these are also equivalent to
\[
  \Lambda_{\rightarrow}=\Lambda_{\leftarrow}=\frac12.
\]
Thus, under pass symmetry, OP~1 is equivalent to saying that the outward and return photon-side passes each carry normalized half-burden.
\end{proposition}

\begin{proof}
By Proposition~\ref{prop:op1-four-shell-certificate}, \(\Delta_{\mathrm{CL}}=0\) is equivalent to
\[
  \lambda_1+\lambda_2+\lambda_3+\lambda_4=1.
\]
Using Definition~\ref{def:op1-two-pass-loads}, this is exactly
\[
  \Lambda_{\rightarrow}+\Lambda_{\leftarrow}=1,
\]
so statements (1)--(3) are equivalent. If, in addition, the profile is pass-symmetric, then \(\Lambda_{\rightarrow}=\Lambda_{\leftarrow}\); together with their sum being one, this gives
\[
  \Lambda_{\rightarrow}=\Lambda_{\leftarrow}=\tfrac12.
\]
Conversely, if both pass totals equal \(\tfrac12\), then their sum is one, and OP~1 closes by the already established equivalences.
\end{proof}

\begin{corollary}[What r05 really reduces]
\label{cor:op1-two-pass-reduction}
Under a pass-symmetric two-pass reading, OP~1 no longer needs a direct identification of all four individual shell loads. It is enough to show that the forward pass contributes normalized burden \(\tfrac12\) and the backward pass contributes normalized burden \(\tfrac12\). In particular, any future lattice, first-zero, shell, or transport lemma that yields pass symmetry together with one half-burden identity closes OP~1 without resolving the internal split inside each pass.
\end{corollary}

\begin{remark}[Why this is the first honest origin candidate for the four \texorpdfstring{$\lambda_j$}{lambdaj}]
This is not yet a first-principles derivation, and it should not be advertised as one. The gain is narrower and more honest. The four shell loads now have a concrete structural reading as \emph{two forward-side loads plus two backward-side loads}. That matches the standing bidirectional transport intuition already present elsewhere in the project and turns OP~1 from a raw four-term balance into a two-pass half-burden question. The next attack can therefore be sharper: not ``guess the four \(\lambda_j\) independently,'' but derive why the post-frame photon transport should be pass-symmetric and why each pass should carry normalized load \(\tfrac12\).
\end{remark}


\subsection*{19AH$^{(16)}$. First reciprocal pass-symmetry certificate for the OP~1 two-pass balance}
After the two-pass half-burden certificate of 19AH$^{(15)}$, the next narrow question is no longer how to guess the four shell loads independently, but why the forward and backward pass totals should agree in the first place. The weakest structurally natural answer is a reciprocity principle: each forward-side shell load has a backward-side partner carrying the same normalized attenuation burden. This does not yet derive the common value of the pairs, but it is enough to force pass symmetry and therefore collapse OP~1 to a one-pass half-balance statement.

\begin{definition}[Reciprocal shell-pairing certificate for the OP~1 transport]
\label{def:op1-reciprocal-shell-pairing}
Under the hypotheses of Definition~\ref{def:op1-two-pass-loads}, we say that the two-pass grouped four-shell profile satisfies the \emph{reciprocal shell-pairing certificate} if the forward/backward shell labels are paired by
\[
  \rho(1)=3,\qquad \rho(2)=4,
\]
and the paired normalized loads agree,
\[
  \lambda_{\rho(j)}=\lambda_j\qquad (j=1,2).
\]
Equivalently,
\[
  \lambda_1=\lambda_3,
  \qquad
  \lambda_2=\lambda_4.
\]
We call this \emph{pairwise reciprocal} because each outward shell burden is mirrored by one return-shell burden of the same normalized weight.
\end{definition}

\begin{proposition}[Reciprocal pairing forces pass symmetry for OP~1]
\label{prop:op1-reciprocal-pass-symmetry}
Suppose the OP~1 residual admits a two-pass grouped four-shell decomposition in the sense of Definition~\ref{def:op1-two-pass-loads}. If the reciprocal shell-pairing certificate of Definition~\ref{def:op1-reciprocal-shell-pairing} holds, then the profile is automatically pass-symmetric:
\[
  \Lambda_{\rightarrow}=\Lambda_{\leftarrow}.
\]
Consequently,
\[
  \Delta_{\mathrm{CL}}=0
  \iff
  \Psi_{\mathrm{CL}}=1
  \iff
  \Lambda_{\rightarrow}=\Lambda_{\leftarrow}=\frac12.
\]
Equivalently, under reciprocal shell pairing, OP~1 is reduced to the single one-pass half-burden identity
\[
  \lambda_1+\lambda_2=\frac12,
\]
or, identically, to the return-pass statement
\[
  \lambda_3+\lambda_4=\frac12.
\]
\end{proposition}

\begin{proof}
If reciprocal shell pairing holds, then
\[
  \lambda_1=\lambda_3,
  \qquad
  \lambda_2=\lambda_4.
\]
Therefore
\[
  \Lambda_{\rightarrow}=\lambda_1+\lambda_2
  =\lambda_3+\lambda_4=\Lambda_{\leftarrow},
\]
so the profile is pass-symmetric. The remaining equivalences are then exactly the pass-symmetric conclusion of Proposition~\ref{prop:op1-two-pass-half-burden}. In particular, once reciprocal pairing is known, OP~1 closes precisely when one of the two pass totals equals \(\tfrac12\), because the other is then forced to equal the same value.
\end{proof}

\begin{corollary}[The live OP~1 burden after r06]
\label{cor:op1-live-burden-after-r06}
If a future shell, lattice, counting, or first-zero transport lemma yields reciprocal shell pairing, then the separate proof obligation for pass symmetry disappears. The live OP~1 burden is then only the one-pass half-balance statement
\[
  \lambda_1+\lambda_2=\frac12,
\]
with the backward pass following automatically by reciprocity.
\end{corollary}

\begin{remark}[Why this is a real narrowing without pretending OP~1 is closed]
This is still not a first-principles derivation of the continuum-limit correction, and it should not be sold as such. The gain is more disciplined. After r05, OP~1 asked for two ingredients: pass symmetry and half-burden. After r06, any concrete reciprocity mechanism collapses the first ingredient into a taut consequence of pairwise equal shell loads. The next attack can therefore become even sharper: not ``why should the two passes be equal in aggregate?'' but ``what concrete structural mechanism enforces reciprocal shell pairing, and why does one pass carry normalized load \(\tfrac12\)?''
\end{remark}


\subsection*{19AH$^{(17)}$. First same-corridor reversal certificate for reciprocal shell pairing in OP~1}
After the reciprocal pass-symmetry certificate of 19AH$^{(16)}$, the next honest question is what concrete structural mechanism could force the pairwise equalities
\[
  \lambda_1=\lambda_3,
  \qquad
  \lambda_2=\lambda_4.
\]
The narrowest bidirectional candidate is the standing two-pass reading itself: the post-frame photon transport uses the \emph{same shell corridor} once in the forward direction and once in the reverse direction. If the normalized attenuation burden is attached to the corridor shells rather than to the direction label, then reciprocal shell pairing follows automatically. This still does not determine the common shell values, but it supplies the first concrete origin mechanism for the r06 reciprocity pattern.

\begin{definition}[Same-corridor reversal certificate for the OP~1 shell transport]
\label{def:op1-same-corridor-reversal}
Assume the grouped four-shell decomposition of Definition~\ref{def:op1-two-pass-loads}. We say that the profile satisfies the \emph{same-corridor reversal certificate} if there exist two nonnegative corridor loads
\[
  \mu_1,\mu_2\ge 0
\]
and a single two-shell corridor such that the forward pass reads the corridor in order
\[
  (\lambda_1,\lambda_2)=(\mu_1,\mu_2),
\]
while the backward pass reads the same corridor in reverse order but with the same shell-attached loads,
\[
  (\lambda_3,\lambda_4)=(\mu_1,\mu_2).
\]
Equivalently, the direction of traversal changes, but the shell-burden assignment is carried by the corridor itself and therefore survives reversal unchanged.
\end{definition}

\begin{proposition}[Same-corridor reversal implies reciprocal shell pairing]
\label{prop:op1-same-corridor-implies-reciprocity}
If the OP~1 grouped four-shell profile satisfies the same-corridor reversal certificate of Definition~\ref{def:op1-same-corridor-reversal}, then the reciprocal shell-pairing certificate of Definition~\ref{def:op1-reciprocal-shell-pairing} holds. Hence
\[
  \lambda_1=\lambda_3=\mu_1,
  \qquad
  \lambda_2=\lambda_4=\mu_2,
\]
and therefore
\[
  \Lambda_{\rightarrow}=\Lambda_{\leftarrow}=\mu_1+\mu_2.
\]
Consequently,
\[
  \Delta_{\mathrm{CL}}=0
  \iff
  \Psi_{\mathrm{CL}}=1
  \iff
  \mu_1+\mu_2=\frac12.
\]
In particular, under same-corridor reversal the full OP~1 burden is reduced to one two-shell half-balance on a single direction-free corridor.
\end{proposition}

\begin{proof}
By definition of same-corridor reversal we have
\[
  \lambda_1=\mu_1=\lambda_3,
  \qquad
  \lambda_2=\mu_2=\lambda_4.
\]
Thus the reciprocal shell-pairing equalities of Definition~\ref{def:op1-reciprocal-shell-pairing} hold automatically. Proposition~\ref{prop:op1-reciprocal-pass-symmetry} then yields pass symmetry, and the pass totals are
\[
  \Lambda_{\rightarrow}=\lambda_1+\lambda_2=\mu_1+\mu_2,
  \qquad
  \Lambda_{\leftarrow}=\lambda_3+\lambda_4=\mu_1+\mu_2.
\]
Since OP~1 is equivalent to \(\Lambda_{\rightarrow}=\Lambda_{\leftarrow}=\tfrac12\) under reciprocal pairing, the final equivalence follows.
\end{proof}

\begin{corollary}[The live OP~1 burden after r07]
\label{cor:op1-live-burden-after-r07}
If a future prime-lattice, shell, or first-zero transport lemma identifies the post-frame photon transport with one reversal-invariant two-shell corridor, then OP~1 no longer asks for a separate reciprocity proof and no longer asks for two independent pass totals. The live burden is only the scalar two-shell corridor identity
\[
  \mu_1+\mu_2=\frac12.
\]
\end{corollary}

\begin{remark}[Why this is the first genuinely concrete reciprocity mechanism]
This is still not the closure of OP~1. The unresolved content has merely been pushed to an even narrower place: the one-corridor half-balance \(\mu_1+\mu_2=\tfrac12\). But unlike the earlier certificates, r07 now gives a concrete structural reason for the reciprocity equalities themselves. The forward and backward passes are no longer treated as two independently weighted routes; they are two traversals of the same post-frame corridor. That is the first point in the present line where the project's standing bidirectional transport picture becomes an explicit OP~1 origin mechanism rather than only a heuristic slogan.
\end{remark}



\subsection*{19AH$^{(18)}$. First unit-loop half-burden certificate for the one-corridor OP~1 transport}
After the same-corridor reversal certificate of 19AH$^{(17)}$, the unresolved OP~1 content has been pushed onto one direction-free two-shell corridor with total pass load \(\mu_1+\mu_2\). The next honest question is therefore no longer where reciprocity comes from, but why this one corridor should carry exactly the half-burden \(\tfrac12\). The narrowest bidirectional candidate is the standing two-pass transport picture itself: the post-frame photon corridor is traversed once forward and once backward, and the normalized burden is attached to the full closed two-pass loop rather than to one preferred orientation. If that closed loop is unit-normalized and the corridor loads are reversal-invariant, then each traversal must carry exactly half of the total loop burden.

\begin{definition}[Unit-loop half-burden certificate for the OP~1 one-corridor transport]
\label{def:op1-unit-loop-half-burden}
Assume the same-corridor reversal certificate of Definition~\ref{def:op1-same-corridor-reversal}, so that a single direction-free two-shell corridor with loads \(\mu_1,\mu_2\ge 0\) is traversed once forward and once backward. We say that the profile satisfies the \emph{unit-loop half-burden certificate} if the full closed two-pass corridor loop is normalized to total burden one, i.e.
\[
  (\mu_1+\mu_2)+(\mu_1+\mu_2)=1.
\]
Equivalently, the forward and backward traversals are not assigned independent total weights; they are the two orientations of one unit-normalized closed transport loop.
\end{definition}

\begin{proposition}[Unit-loop normalization forces the OP~1 half-burden]
\label{prop:op1-unit-loop-half-burden}
If the OP~1 transport satisfies the same-corridor reversal certificate of Definition~\ref{def:op1-same-corridor-reversal} and the unit-loop half-burden certificate of Definition~\ref{def:op1-unit-loop-half-burden}, then
\[
  \mu_1+\mu_2=\frac12.
\]
Consequently,
\[
  \Lambda_{\rightarrow}=\Lambda_{\leftarrow}=\frac12,
\qquad
  \Psi_{\mathrm{CL}}=1,
\qquad
  \Delta_{\mathrm{CL}}=0.
\]
In particular, under same-corridor reversal plus unit-loop normalization, OP~1 is closed.
\end{proposition}

\begin{proof}
By Definition~\ref{def:op1-unit-loop-half-burden},
\[
  2(\mu_1+\mu_2)=1.
\]
Hence \(\mu_1+\mu_2=\tfrac12\). Proposition~\ref{prop:op1-same-corridor-implies-reciprocity} then gives
\[
  \Lambda_{\rightarrow}=\Lambda_{\leftarrow}=\mu_1+\mu_2=\frac12.
\]
By Proposition~\ref{prop:op1-same-corridor-implies-reciprocity}, this is equivalent to \(\Psi_{\mathrm{CL}}=1\), and by Proposition~\ref{prop:op1-slope-certificate} it is equivalent to \(\Delta_{\mathrm{CL}}=0\). Thus OP~1 is closed under the stated hypotheses.
\end{proof}

\begin{corollary}[The live OP~1 burden after r08]
\label{cor:op1-live-burden-after-r08}
After r08, the remaining live content of OP~1 is no longer a free four-shell, two-pass, or reciprocity problem. It is the identification of the post-frame photon corridor with a \emph{unit-normalized closed two-pass loop}. Once that normalization is obtained from a future prime-lattice, shell, or first-zero transport lemma, the one-corridor half-balance and therefore OP~1 follow immediately.
\end{corollary}

\begin{remark}[Why r08 is a real advance without pretending the normalization is already proved]
This does not yet prove that the post-frame photon corridor really is unit-normalized. That point remains the live burden. But the nature of the burden has changed sharply. After r07, OP~1 asked for one scalar half-balance on a direction-free corridor. After r08, even that half-balance is no longer independent: it is forced as soon as the corridor is shown to be the correct unit-normalized closed two-pass object. So the next attack does not need to juggle shell weights, pass symmetry, or reciprocity any more. It only needs to justify one normalization principle for one already isolated bidirectional corridor.
\end{remark}


\subsection*{19AH$^{(19)}$. First one-first-zero loop-budget certificate for OP~1 after the unit-loop reduction}
After the unit-loop half-burden certificate of 19AH$^{(18)}$, the remaining burden is no longer to guess a half-load on the corridor. What remains is to explain why the isolated closed two-pass corridor should carry \emph{one} normalized transport unit in the first place. The frame split from 19AH$^{(13)}$ already shows that the surviving OP~1 normalization is measured in units of the photon-normalized first zero \(\gamma_1^{(\mathrm{photon})}\). The sharpest next compression is therefore to re-read the unit-loop hypothesis not as an external normalization convention, but as a multiplicity-one statement: the post-frame corridor is exactly one closed first-zero transport loop, not a fraction of one loop and not a multiple copy of several such loops.

\begin{definition}[One-first-zero loop-budget certificate for the OP~1 corridor]
\label{def:op1-one-first-zero-loop-budget}
Assume the same-corridor reversal certificate of Definition~\ref{def:op1-same-corridor-reversal}, so that the post-frame OP~1 transport is carried by one direction-free two-shell corridor with loads \(\mu_1,\mu_2\). Define its \emph{closed-loop budget index} by
\[
  \Omega_{\mathrm{loop}} := 2(\mu_1+\mu_2).
\]
We say that the profile satisfies the \emph{one-first-zero loop-budget certificate} if the isolated closed two-pass corridor represents exactly one photon-first-zero transport unit, i.e.
\[
  \Omega_{\mathrm{loop}} = 1.
\]
Equivalently, the corridor is read as a multiplicity-one closed loop in the photon-first-zero normalization of 19AH$^{(13)}$.
\end{definition}

\begin{proposition}[One-first-zero loop budget identifies the r08 unit-loop normalization]
\label{prop:op1-one-first-zero-loop-budget}
Assume the same-corridor reversal certificate of Definition~\ref{def:op1-same-corridor-reversal}. Then
\[
  \Omega_{\mathrm{loop}} = \Psi_{\mathrm{CL}}.
\]
Consequently, if the one-first-zero loop-budget certificate of Definition~\ref{def:op1-one-first-zero-loop-budget} holds, then the unit-loop half-burden certificate of Definition~\ref{def:op1-unit-loop-half-burden} holds as well, and therefore
\[
  \Delta_{\mathrm{CL}} = 0.
\]
In particular, under same-corridor reversal, OP~1 is closed as soon as the isolated corridor is shown to be a multiplicity-one photon-first-zero loop.
\end{proposition}

\begin{proof}
By Proposition~\ref{prop:op1-same-corridor-implies-reciprocity},
\[
  \Lambda_{\rightarrow}=\Lambda_{\leftarrow}=\mu_1+\mu_2.
\]
Hence Proposition~\ref{prop:op1-two-pass-half-burden} gives
\[
  \Psi_{\mathrm{CL}} = \Lambda_{\rightarrow}+\Lambda_{\leftarrow}
  = (\mu_1+\mu_2)+(\mu_1+\mu_2)
  = 2(\mu_1+\mu_2)
  = \Omega_{\mathrm{loop}}.
\]
If now \(\Omega_{\mathrm{loop}}=1\), then by definition
\[
  (\mu_1+\mu_2)+(\mu_1+\mu_2)=1,
\]
which is exactly the unit-loop half-burden certificate of Definition~\ref{def:op1-unit-loop-half-burden}. Proposition~\ref{prop:op1-unit-loop-half-burden} then yields \(\Delta_{\mathrm{CL}}=0\).
\end{proof}

\begin{corollary}[The live OP~1 burden after r09]
\label{cor:op1-live-burden-after-r09}
After r09, the remaining live burden of OP~1 is no longer to supply an abstract unit-loop normalization from outside the frame split. It is to prove a \emph{multiplicity-one statement}: the isolated post-frame corridor is exactly one photon-first-zero closed loop. Equivalently, one must rule out both undercounting and overcounting, i.e. show that the corridor does not represent a fractional share of the first-zero budget and does not split into several loop copies carrying total index \(\Omega_{\mathrm{loop}}\neq 1\).
\end{corollary}

\begin{remark}[Why r09 is a real tightening of the OP~1 queue]
This is still not the final closure proof, because the multiplicity-one statement has not yet been derived from the prime-lattice/first-zero transport itself. But r09 removes one remaining piece of arbitrariness. In r08, the decisive hypothesis was a unit-normalized closed loop. After r09, that hypothesis is no longer a free normalization slogan: it is identified with the claim that the already isolated corridor carries exactly one photon-first-zero transport unit. So the next attack on OP~1 is even narrower than before. It should no longer ask for ``why does some loop have weight one?'' in the abstract, but for a concrete first-zero transport reason why the admissible corridor is a single normalized loop rather than a fractional or multiple-count loop.
\end{remark}



\subsection*{19AH$^{(20)}$. First admissible-window uniqueness certificate for the multiplicity-one OP~1 corridor}
After the one-first-zero loop-budget certificate of 19AH$^{(19)}$, the live OP~1 burden is no longer a free normalization question. It is a uniqueness question on the retained admissible corridor: why should the isolated post-frame same-corridor loop represent \emph{exactly one} admissible photon-first-zero transport unit rather than several parallel admissible copies? The sharpest next compression is therefore to move from multiplicity language to admissible-window uniqueness language. On the present line, this is the first point at which the older no-second-carrier discipline and the newer total/imbalance burden split directly feed back into OP~1: extra loop multiplicity can only matter if it either changes the retained total-channel profile or introduces a genuinely new primitive carrier contribution on the same admissible corridor.

\begin{definition}[Admissible-window uniqueness certificate for the OP~1 corridor]
\label{def:op1-admissible-window-uniqueness}
Assume the same-corridor reversal certificate of Definition~\ref{def:op1-same-corridor-reversal}, so that the post-frame OP~1 transport is carried by one direction-free two-shell corridor. We say that the profile satisfies the \emph{admissible-window uniqueness certificate} if, on the intended admissible machine-bearing windows of the present line, that isolated corridor admits exactly one primitive photon-first-zero closed-loop realization compatible with the retained total-channel profile and with the no-second-carrier discipline. Equivalently, every admissible realization of the corridor that preserves the retained total-channel load is identified with the same single primitive first-zero loop.
\end{definition}

\begin{proposition}[Admissible-window uniqueness upgrades multiplicity-one to OP~1 closure]
\label{prop:op1-admissible-window-uniqueness}
Assume the same-corridor reversal certificate of Definition~\ref{def:op1-same-corridor-reversal}. If the admissible-window uniqueness certificate of Definition~\ref{def:op1-admissible-window-uniqueness} holds, then the one-first-zero loop-budget certificate of Definition~\ref{def:op1-one-first-zero-loop-budget} holds. Consequently,
\[
  \Delta_{\mathrm{CL}} = 0.
\]
In particular, after r10 the surviving OP~1 burden may be read as a pure admissible-window uniqueness question for the retained post-frame corridor.
\end{proposition}

\begin{proof}
By Definition~\ref{def:op1-admissible-window-uniqueness}, the isolated post-frame same-corridor loop has exactly one admissible primitive photon-first-zero realization compatible with the retained total-channel profile and the no-second-carrier discipline. But Definition~\ref{def:op1-one-first-zero-loop-budget} identifies precisely this multiplicity-one situation with the statement that the corridor carries one photon-first-zero transport unit, i.e. \(\Omega_{\mathrm{loop}}=1\). Hence the one-first-zero loop-budget certificate holds. Proposition~\ref{prop:op1-one-first-zero-loop-budget} then yields \(\Delta_{\mathrm{CL}}=0\).
\end{proof}

\begin{corollary}[The live OP~1 burden after r10]
\label{cor:op1-live-burden-after-r10}
After r10, the remaining live burden of OP~1 is no longer to choose a normalization, a pass split, or even an abstract multiplicity slogan. It is to prove a \emph{uniqueness statement on admissible windows}: the retained post-frame corridor supports one and only one primitive photon-first-zero loop compatible with the preserved total-channel load. Any alternative admissible realization would have to appear either as a total-channel distortion or as a second primitive carrier contribution, and therefore lies outside the intended admissible OP~1 corridor class.
\end{corollary}

\begin{remark}[Why r10 is a genuine narrowing and not just another renaming]
This still does not by itself derive the uniqueness statement from first principles; so OP~1 is not yet honestly closed on theorem level. But the queue is now much sharper. Before r10 one could still phrase the task as ``explain multiplicity one'' in a somewhat free way. After r10 the task is pinned to the already existing admissible-window discipline of the present line: one must show that the retained corridor does not admit several distinct primitive first-zero loop realizations once the total-channel profile and no-second-carrier constraints are held fixed. So the next attack on OP~1 should not invent new continuum-limit parameters. It should directly use retained-window rigidity/uniqueness to rule out duplicate admissible loop occupancy on the same corridor.
\end{remark}


\subsection*{19AH$^{(21)}$. Retained-window-rigidity / no-second-carrier bridge for OP~1}
After the admissible-window uniqueness certificate of 19AH$^{(20)}$, the remaining honest task is no longer to invent a new normalization slogan or a new loop-counting parameter. It is to derive uniqueness from structural disciplines that are already native to the present line. The sharpest triad is now visible. First, the retained admissible corridor should be rigid once the retained total-channel profile and optimum data are fixed. Second, the older no-second-carrier discipline should forbid an additional primitive occupancy of that same retained corridor. Third, the strict-upper-gap discipline should forbid a genuinely distinct admissible completion route above the same retained corridor that would duplicate the first-zero budget without altering the present-line status surface. Together these three pressures are exactly what the r10 uniqueness question is asking for in compressed form.

\begin{definition}[Retained-window rigidity certificate for the OP~1 corridor]
\label{def:op1-retained-window-rigidity}
Assume the same-corridor reversal certificate of Definition~\ref{def:op1-same-corridor-reversal}. We say that the profile satisfies the \emph{retained-window rigidity certificate} on the intended admissible windows if every admissible primitive photon-first-zero closed-loop realization on the retained post-frame corridor that preserves the retained total-channel profile, optimum location, and standing admissible equivalence data occupies the same retained corridor slot up to the standing admissible corridor equivalence. Equivalently, there is no nontrivial admissible primitive displacement of the retained first-zero loop inside the same preserved total-channel window.
\end{definition}

\begin{proposition}[Rigidity plus no-second-carrier plus strict upper gap imply OP~1 uniqueness]
\label{prop:op1-rigidity-gap-uniqueness-bridge}
Assume the same-corridor reversal certificate of Definition~\ref{def:op1-same-corridor-reversal}. Assume further that the retained-window rigidity certificate of Definition~\ref{def:op1-retained-window-rigidity} holds on the intended admissible windows, that the no-second-carrier discipline continues to hold at the retained total-channel level, and that the strict upper gap admits no second independent admissible completion route above the same retained corridor on those windows. Then the admissible-window uniqueness certificate of Definition~\ref{def:op1-admissible-window-uniqueness} holds. Consequently,
\[
  \Delta_{\mathrm{CL}} = 0.
\]
In particular, after r11 the live OP~1 burden may be read as the extraction of retained-window rigidity on the post-frame admissible corridor, with the no-second-carrier and strict-upper-gap disciplines supplying the exclusion of duplicate primitive occupancy and duplicate admissible completion.
\end{proposition}

\begin{proof}
Let \(C'\) be any admissible primitive photon-first-zero closed-loop realization on the retained post-frame corridor that preserves the retained total-channel profile. By the retained-window rigidity certificate, \(C'\) occupies the same retained corridor slot as the original loop up to the standing admissible corridor equivalence. If \(C'\) were still a genuinely different primitive realization, then either it would contribute an extra primitive occupancy on that same retained corridor, contradicting the no-second-carrier discipline at total-channel level, or it would represent a second independent admissible completion route above the same corridor while preserving the retained total-channel surface, contradicting the standing strict-upper-gap exclusion on the intended windows. Hence no genuinely distinct admissible primitive realization exists, and the admissible-window uniqueness certificate of Definition~\ref{def:op1-admissible-window-uniqueness} follows. Proposition~\ref{prop:op1-admissible-window-uniqueness} then yields \(\Delta_{\mathrm{CL}}=0\).
\end{proof}

\begin{corollary}[The live OP~1 burden after r11]
\label{cor:op1-live-burden-after-r11}
After r11, OP~1 is no longer honestly read as a continuum-limit mystery in general, nor even as a free uniqueness slogan. The live burden is narrower: extract a retained-window rigidity theorem for the already isolated post-frame admissible corridor. Once that rigidity is in hand, the existing no-second-carrier discipline blocks duplicate primitive corridor occupancy and the strict-upper-gap discipline blocks duplicate admissible completion above the same retained corridor, so the uniqueness required in 19AH$^{(20)}$ follows automatically.
\end{corollary}

\begin{remark}[Why r11 is a real narrowing]
This is still not a final unconditional OP~1 theorem, because the retained-window rigidity certificate itself has not yet been discharged from the active comparison/transport line. But the queue is now much sharper than in r10. The remaining question is no longer ``why uniqueness?'' in the abstract. It is ``why is the retained admissible corridor rigid once the present-line optimum and total-channel data are fixed?'' That is exactly the kind of local structural statement the existing retained-window diagnostics, no-second-carrier discipline, and strict-upper-gap language are already built to support.
\end{remark}


\subsection*{19AH$^{(22)}$. First optimum-pinning comparison-rigidity bridge for retained-window rigidity in OP~1}
After the retained-window-rigidity / no-second-carrier / strict-upper-gap bridge of 19AH$^{(21)}$, the live OP~1 task is sharper still: why should the retained admissible corridor be rigid once the present-line comparison data are fixed? The most honest next compression is to stop treating rigidity as a free geometric slogan and to pin it directly to the already retained comparison surface. On the present line, the retained corridor is not specified by corridor language alone. It is specified together with the preserved total-channel profile, the retained optimum location, and the local comparison/defect data that survive the machine-shadow compression. If those comparison data already determine the admissible slot uniquely inside the retained corridor class, then retained-window rigidity is no longer an extra burden; it is simply the local comparison line read in the right coordinates.

\begin{definition}[Optimum-pinning comparison-rigidity certificate for OP~1]
\label{def:op1-optimum-pinning-comparison-rigidity}
Assume the same-corridor reversal certificate of Definition~\ref{def:op1-same-corridor-reversal}. We say that the profile satisfies the \emph{optimum-pinning comparison-rigidity certificate} on the intended admissible windows if every admissible primitive photon-first-zero closed-loop realization on the retained post-frame corridor that preserves the retained total-channel profile, the retained optimum location, and the standing local comparison/defect data is forced to occupy the same retained corridor slot up to the standing admissible corridor equivalence. Equivalently, within the retained admissible corridor class there is no second primitive loop placement that keeps the same optimum pinning and the same local comparison readout while remaining genuinely distinct.
\end{definition}

\begin{proposition}[Optimum pinning and comparison rigidity imply retained-window rigidity for OP~1]
\label{prop:op1-optimum-pinning-implies-rigidity}
Assume the same-corridor reversal certificate of Definition~\ref{def:op1-same-corridor-reversal}. If the optimum-pinning comparison-rigidity certificate of Definition~\ref{def:op1-optimum-pinning-comparison-rigidity} holds on the intended admissible windows, then the retained-window rigidity certificate of Definition~\ref{def:op1-retained-window-rigidity} holds there as well. Consequently, together with the no-second-carrier and strict-upper-gap disciplines from Proposition~\ref{prop:op1-rigidity-gap-uniqueness-bridge}, one obtains the admissible-window uniqueness certificate of Definition~\ref{def:op1-admissible-window-uniqueness} and therefore
\[
  \Delta_{\mathrm{CL}} = 0.
\]
In particular, after r12 the live OP~1 burden may be read as a local optimum-pinning comparison-rigidity statement on the retained post-frame corridor.
\end{proposition}

\begin{proof}
Take any admissible primitive photon-first-zero closed-loop realization on the retained post-frame corridor that preserves the retained total-channel profile. On the intended admissible windows of the present line, the retained OP~1 corridor is not read independently of the standing comparison surface: the intended loop class is already tracked together with the retained optimum location and the surviving local comparison/defect data of the compressed machine-shadow line. If the optimum-pinning comparison-rigidity certificate holds, then any admissible primitive realization preserving those retained data must occupy the same retained corridor slot up to the standing admissible corridor equivalence. But this is exactly the content of the retained-window rigidity certificate. The rest follows from Proposition~\ref{prop:op1-rigidity-gap-uniqueness-bridge}, which upgrades retained-window rigidity to admissible-window uniqueness under the standing no-second-carrier and strict-upper-gap exclusions, hence yields \(\Delta_{\mathrm{CL}}=0\).
\end{proof}

\begin{corollary}[The live OP~1 burden after r12]
\label{cor:op1-live-burden-after-r12}
After r12, OP~1 is no longer honestly carried as a general rigidity problem. The live burden is narrower: show that the retained optimum and local comparison data already pin the admissible post-frame first-zero corridor slot uniquely. Once that local comparison rigidity is extracted, retained-window rigidity follows, and with the standing no-second-carrier and strict-upper-gap disciplines the OP~1 closure chain completes automatically.
\end{corollary}

\begin{remark}[Why r12 is the first direct attack from the comparison line]
This is still not a full unconditional theorem extracted entirely from the earlier comparison blocks; the optimum-pinning comparison-rigidity certificate itself is now the remaining live statement. But it is a real narrowing. r11 said, in effect, ``prove retained-window rigidity.'' r12 says something more local and more testable: ``prove that no genuinely distinct primitive loop placement preserves the retained optimum and the surviving comparison/defect readout on the same admissible corridor.'' That is the first point at which the residual OP~1 burden is read directly as a comparison-line pinning problem rather than as a free rigidity slogan.
\end{remark}

\subsection*{19AI. Bundled overlay-completion capstone for the machine-theoretic line}
After the machine-theoretic line has reached full machine-theoretic closure, present-line status readout, and window-equivalence, the remaining overlay refinements should no longer be read as an invitation to keep appending parallel adjectives indefinitely. Their real purpose is to certify, in a bundled way, that the machine layer is the right kind of organization of the already certified machine-bearing windows: it adds no foreign admissible content, forgets none of the operative distinctions relevant for the machine hierarchy, tracks the intended statuses exactly, and returns the same machine-bearing layer up to the standing admissible equivalences. The correct freeze-level conclusion is therefore not a new ladder of further micro-properties, but a single capstone statement that the machine overlay is a conservative, faithful, exact, relatively complete, adequate, reconstructive, closed, idempotent, minimal, canonical, invariant, natural, universally characterized, rigid, terminal, and equivalence-level organization of the intended certified machine-bearing windows.

\begin{definition}[Bundled machine-overlay completion certificate on the compressed computational line]
Fix a chosen sequential reference core \(\mathcal R\) and an intended admissible machine-bearing window class \(\mathcal W^{\mathrm{mach}}\). We say that the compressed computational line satisfies the \emph{bundled machine-overlay completion certificate} on \(\mathcal W^{\mathrm{mach}}\) if, once the full machine-theoretic closure surface \(\mathfrak M^{\mathrm{full}}(\mathcal R)=1\), the present-line status surface \(\mathfrak P^{\mathrm{mach}}(\mathcal R)\), and the window-equivalence/equivalence-surface layer hold on \(\mathcal W^{\mathrm{mach}}\), the resulting machine-theoretic overlay is conservative, faithful, exact, relatively complete, adequate, reconstructive, closed, idempotent, minimal, canonical, invariant, natural, universally characterized, rigid, terminal, and equivalent to the certified machine-bearing structural content, all up to the standing admissible window-equivalence and overlay-equivalence. Equivalently, the intended admissible machine-bearing windows and the canonical machine-overlay objects form two fully compatible organizations of the same certified machine-bearing layer.
\end{definition}

\begin{proposition}[Bundled overlay-completion theorem for the machine-theoretic line]
On the intended admissible machine-bearing windows of the present compressed computational line, the above bundled machine-overlay completion certificate holds. More precisely, the machine-theoretic line of Sections~19J--19AH together with the diagnostic continuation block 19AH$\prime$--19AH$^{(22)}$ yields the indispensable ladder from local operational semantics to the present-line conditional/unconditional machine status plus the first disciplined re-reading of the retained machine shadow, while the subsequent overlay-completion work can be read in bundled form as showing that this machine layer is conservative, faithful, exact, relatively complete, adequate, reconstructive, closed, idempotent, minimal, canonical, invariant, natural, universally characterized, rigid, terminal, and equivalent to the intended certified machine-bearing windows. Consequently, for the present Companion purpose the machine-theoretic overlay may be cited as a closed reversible machine-level organization of the certified machine-bearing structural content, without needing to unfold each of these overlay attributes separately unless a later revision introduces a genuinely new machine-theoretic distinction.
\end{proposition}

\begin{corollary}[What the bundled overlay-completion capstone already buys]
The machine line now reaches a true freeze-level stopping point. For subsequent work it is enough to cite the indispensable machine ladder together with the bundled overlay-completion capstone, rather than to continue lengthening the line by parallel closure adjectives. In particular, the present line can now be read through the three decisive bundled surfaces --- the full machine-theoretic closure surface, the present-line status surface, and the bundled overlay-completion capstone --- while the inserted channel-split diagnostics remain the preferred local reading tools whenever carrier and observer-side burdens need to be separated.
\end{corollary}

\subsection*{19AZ. Freeze-audit verdict on the machine-theoretic line}
In this inherited release-baseline block, the machine-theoretic line reaches the release-level freeze threshold used for the then-present public version. A local audit of Sections~19J--19AY shows that the computational route itself is already complete enough for the present Companion purpose, while the last overlay block from 19AI onward increasingly repackages nearby closure properties rather than opening substantially new machine content. For this reason that inherited line should now be read with a strict two-tier discipline. Theorems~19J--19AH remain the indispensable machine ladder: local operational semantics, conditional Zuse-style primitive core, reference-core closure, relative tape-family upgrade, the strict upper gap, the full machine-theoretic closure surface, the status decoder, and the present-line conditional/unconditional status surface. The local continuation block 19AH$\prime$--19AH$^{(22)}$ should be read only as a diagnostic refinement layer: it separates total and imbalance channels, turns that split into a control dictionary, sorts inherited continuation knots by burden type, and compresses the surviving OP~1 queue to the present optimum-pinning comparison-rigidity kernel. Section~19AI now functions as the bundled overlay-completion capstone: it certifies in compressed form that the machine line is conservative, faithful, exact, relatively complete, adequate, reconstructive, closed, idempotent, minimal, canonical, invariant, natural, universally characterized, rigid, terminal, and equivalent to the certified machine-bearing windows on the intended admissible domain. The audit verdict in that inherited release block is therefore to stop expansion there: future revisions should preferentially cite the bundled surfaces already obtained --- in particular 19AF, 19AH, and 19AI --- rather than append further parallel overlay adjectives unless a genuinely new machine-theoretic distinction is introduced.

\subsection*{19AZ$\prime$. Post-release status note after tool/QAM completion and local OP-kernel closure}
The public v3.28.0 release completed the formerly deferred tool/QAM/decoder layer and fixed the route-control infrastructure of the compressed computational line. The present post-release working line then used that infrastructure exactly as intended: it reopened the two parked residual burdens through the smallest licensed routes and followed them to local kernel closure. OP~1 was normalized first as a retained pinning-signature injectivity problem, then split into occupancy-side and completion-side duplicate-signature witnesses, and finally discharged under the standing retained no-second-carrier and strict-upper-gap exclusions. OP~5 Step~(B) was normalized as a carrier-upgraded routed-balance witness problem, split into carrier-stable versus obstructed routed returns, then into upstream-screened versus downstream carrier-grade obstruction, and finally discharged on the retained branch-shell corridor by the inherited shell-closure package.

The honest status after these steps is therefore no longer ``two parked residual burdens.'' It is: the completed tool/QAM/decoder infrastructure now stands as fixed route-control surface, and the live OP~1 and OP~5 Step~(B) kernels have both been discharged locally under the standing retained disciplines. What remains conditional is not the local kernel closure but the broader machine-theoretic reading beyond the certified retained corridors and routes. The governing public DOI in the chronology is \href{https://doi.org/10.5281/zenodo.19699827}{10.5281/zenodo.19699827} for the public v3.31.0 release.

\subsection*{19AZ$''$. Immediate normalized channel-sorted tool ledger}
\label{sec:immediate-normalized-tool-ledger}
The first concrete task of the post-freeze continuation line is not to invent a fourth proof language, but to normalize the already extracted tools into a channel-sorted ledger. On the present line this means three reusable packets: one for statements that are already total-channel stable, one for statements that compress the surviving imbalance-side burden, and one for statements that translate between both surfaces or between their QAM/matrix realizations. The point is discipline: later arguments should be able to cite the smallest packet that already matches their burden class.

\begin{definition}[Normalized channel-sorted tool packets]
\label{def:normalized-channel-sorted-tool-packets}
The immediate continuation line organizes the current extracted tools into three provisional packets.
\begin{enumerate}[leftmargin=2em]
  \item \textbf{Total-channel packet} \(\mathfrak T_{\mathrm{tot}}\). This packet contains the tools whose role is to keep the retained total-channel architecture fixed once admissible transport has already selected the active branch. Its current generators are the transport-first tool, the readout-after-transport tool, the admissible-readout uniqueness tool, and the global pass/transfer/diagnostic-subordination interface from Propositions~\ref{prop:global-pass-tool}--\ref{prop:global-diagnostic-subordination-tool}.

  \item \textbf{Imbalance packet} \(\mathfrak T_{\mathrm{imb}}\). This packet contains the tools that compress the surviving OP~1 burden after the total-channel surface has been fixed. Its current generators are the frame-transport split certificate (Proposition~\ref{prop:op1-frame-transport-split}), the finite four-shell certificate (Proposition~\ref{prop:op1-four-shell-certificate}), the two-pass half-burden certificate (Proposition~\ref{prop:op1-two-pass-half-burden}), the reciprocal pass-symmetry certificate (Proposition~\ref{prop:op1-reciprocal-pass-symmetry}), the same-corridor reversal implication (Proposition~\ref{prop:op1-same-corridor-implies-reciprocity}), the unit-loop half-burden certificate (Proposition~\ref{prop:op1-unit-loop-half-burden}), the one-first-zero loop-budget bridge (Proposition~\ref{prop:op1-one-first-zero-loop-budget}), the admissible-window uniqueness bridge (Proposition~\ref{prop:op1-admissible-window-uniqueness}), the retained-window-rigidity / no-second-carrier / strict-upper-gap bridge (Proposition~\ref{prop:op1-rigidity-gap-uniqueness-bridge}), and the optimum-pinning comparison-rigidity bridge (Proposition~\ref{prop:op1-optimum-pinning-implies-rigidity}).

  \item \textbf{Bridge packet} \(\mathfrak T_{\mathrm{br}}\). This packet contains the tools whose job is to move between total and imbalance language, or between channel language and its shell/menu/matrix realizations. Its current generators are the shared matrix-leakage tool, the cryptanalytic translation tool, the shared phase-coupling tool, and on the QAM side the shell-stable readout tool, the commuting-square tool, the obstruction-stratification tool, and the carrier-compression tool.
\end{enumerate}
\end{definition}

\begin{proposition}[Channel-sorted citation rule for the immediate continuation line]
\label{prop:channel-sorted-citation-rule}
Let a later argument remain inside the currently admitted OP1$\to$branch$\to$OP3/QAM$\to$OP5 architecture. Then the default citation discipline is as follows.
\begin{enumerate}[leftmargin=2em]
  \item If the argument only needs preservation of the retained total-channel profile or the subordination of visible diagnostics to one already fixed carrier package, it should cite the total-channel packet \(\mathfrak T_{\mathrm{tot}}\).
  \item If the argument accepts the retained total-channel surface as fixed and only compresses the surviving OP~1 burden toward the optimum-pinning kernel, it should cite the imbalance packet \(\mathfrak T_{\mathrm{imb}}\).
  \item If the argument must translate between total and imbalance language, or between channel language and a shell/menu/matrix realization, it should cite the bridge packet \(\mathfrak T_{\mathrm{br}}\).
\end{enumerate}
In particular, the default rule of the next revision is: reopen only the smallest packet that still sees the live burden.
\end{proposition}

\begin{proof}
This is exactly the burden split already extracted in the frozen line. The present proposition merely normalizes that split into a reusable citation discipline: total-channel stability is carried by the total packet, OP~1 compression by the imbalance packet, and all language-changing or layer-changing moves by the bridge packet.
\end{proof}

\begin{remark}[What the normalized ledger already gains]
\label{rem:normalized-ledger-gain}
The gain is immediate even before any new theorem is proved. First, repeated later arguments no longer need to cite a mixed list of total, imbalance, and translation tools in arbitrary order. Second, the surviving live burden is seen more honestly: OP~1 is not ``everything left,'' but specifically the active content of \(\mathfrak T_{\mathrm{imb}}\) once \(\mathfrak T_{\mathrm{tot}}\) has already been frozen. Third, the bridge packet isolates exactly where QAM, matrix, and cryptanalytic language should re-enter: only when one genuinely needs a change of surface.
\end{remark}

\subsection*{19AZ$'''$. QAM recoding of the normalized tool ledger}
\label{sec:qam-recoding-normalized-tool-ledger}
Once the channel-sorted ledger has been fixed, the right immediate QAM task is no longer to widen the public syntax blindly. It is to recode the already stabilized tool packets as a finite theorem-store/menu layer so that later proofs can query, combine, and route them without reopening the full historical derivation. This keeps QAM in its intended role: a definitional overlay and machine-readable proof interface, not a second independent mathematics.

\begin{definition}[QAM-coded tool menu]
\label{def:qam-coded-tool-menu}
Let \(\mathfrak T_{\mathrm{tot}}\), \(\mathfrak T_{\mathrm{imb}}\), and \(\mathfrak T_{\mathrm{br}}\) be the normalized packets from Definition~\ref{def:normalized-channel-sorted-tool-packets}. Their immediate QAM recoding is the finite coded menu
\[
\mathcal Q_{\mathrm{tool}}
:=
\bigl\{\langle \mathrm{tot},\tau\rangle : \tau\in\mathfrak T_{\mathrm{tot}}\bigr\}
\cup
\bigl\{\langle \mathrm{imb},\tau\rangle : \tau\in\mathfrak T_{\mathrm{imb}}\bigr\}
\cup
\bigl\{\langle \mathrm{br},\tau\rangle : \tau\in\mathfrak T_{\mathrm{br}}\bigr\}.
\]
Here the first coordinate records the burden class and the second coordinate records the already extracted uncoded tool. The purpose of \(\mathcal Q_{\mathrm{tool}}\) is indexing, routing, and proof compression; it does not create a new semantic layer beyond the underlying tool statements.
\end{definition}

\begin{definition}[Menu-to-carrier reduction]
\label{def:menu-to-carrier-reduction}
A current local proof obligation is said to admit a \emph{menu-to-carrier reduction} if, after coding the obligation into the finite menu \(\mathcal Q_{\mathrm{tool}}\), there exists a finite submenu \(M\subseteq \mathcal Q_{\mathrm{tool}}\) such that the underlying uncoded tools already imply either
\begin{enumerate}[leftmargin=2em]
  \item one bundled carrier-stable pass statement, or
  \item one strictly smaller localized obstruction class than the original uncoded presentation displayed.
\end{enumerate}
Thus the menu does not replace proof content; it compresses the route to the smallest carrier-bearing or obstruction-bearing interface.
\end{definition}

\begin{proposition}[QAM reuse principle for the immediate continuation line]
\label{prop:qam-reuse-principle-immediate-line}
If a current proof obligation admits menu-to-carrier reduction in the sense of Definition~\ref{def:menu-to-carrier-reduction}, then the next continuation line may process that obligation through the finite coded menu \(\mathcal Q_{\mathrm{tool}}\) instead of reopening the full local historical derivation. In particular, QAM now functions as the normalized theorem-store and routing overlay for the post-freeze tool layer.
\end{proposition}

\begin{proof}
By construction, every element of \(\mathcal Q_{\mathrm{tool}}\) is only a coded wrapper around an already extracted uncoded tool. If a finite submenu already implies the same bundled carrier statement or a strictly smaller obstruction class, then reopening the full derivation adds no new structural information. The coded menu therefore acts as a conservative proof-compression interface and not as a source of new independent mathematical content.
\end{proof}

\begin{remark}[Why this QAM step is mathematically honest]
\label{rem:qam-step-mathematically-honest}
This is the right scale for the immediate QAM revision. The present line does \emph{not} claim that the finite coded menu solves OP~1 or OP~5 automatically. It claims something narrower and more useful: once the main burden has already been split into total, imbalance, and bridge packets, QAM can store that split explicitly, route later obligations to the smallest relevant packet, and expose where proof compression is real rather than cosmetic.
\end{remark}


\subsection*{19AZ$''''$. First finite submenus and certified reuse routes inside the QAM tool menu}
\label{sec:first-finite-submenus-certified-reuse-routes}
The next optimization step is therefore to stop treating \(\mathcal Q_{\mathrm{tool}}\) as one undifferentiated finite store. For later proof work one wants a small number of \emph{canonical finite submenus} that already match the most common post-freeze obligations. These submenus do not add new mathematics; they only isolate the smallest coded interfaces that repeatedly suffice.

\begin{definition}[First canonical finite submenus of \(\mathcal Q_{\mathrm{tool}}\)]
\label{def:first-canonical-finite-submenus-qtool}
Inside the coded menu \(\mathcal Q_{\mathrm{tool}}\), define the following four canonical finite submenus.
\begin{enumerate}[leftmargin=2em]
  \item \textbf{Total-rigidity submenu} \(M^{\mathrm{rig}}_{\mathrm{tot}}\). This is the coded submenu generated by the transport-first tool, the readout-after-transport tool, the admissible-readout uniqueness tool, and the global pass/transfer/diagnostic-subordination interface already collected in the total-channel packet \(\mathfrak T_{\mathrm{tot}}\).

  \item \textbf{Optimum-pinning submenu} \(M^{\mathrm{pin}}_{\mathrm{imb}}\). This is the coded submenu generated by the frame-transport split certificate, the finite four-shell certificate, the two-pass half-burden certificate, the reciprocal pass-symmetry certificate, the same-corridor reversal implication, the unit-loop half-burden certificate, the one-first-zero loop-budget bridge, the admissible-window uniqueness bridge, the retained-window-rigidity / no-second-carrier / strict-upper-gap bridge, and the optimum-pinning comparison-rigidity bridge already collected in the imbalance packet \(\mathfrak T_{\mathrm{imb}}\).

  \item \textbf{Surface-change submenu} \(M^{\mathrm{surf}}_{\mathrm{br}}\). This is the coded submenu generated by the shared matrix-leakage tool, the cryptanalytic translation tool, the shared phase-coupling tool, the shell-stable readout tool, and the commuting-square tool from the bridge packet \(\mathfrak T_{\mathrm{br}}\). Its purpose is to move between channel language and shell/menu/matrix language without reopening the whole local history.

  \item \textbf{Obstruction-localization submenu} \(M^{\mathrm{loc}}_{\mathrm{br}}\). This is the coded submenu generated by the shared matrix-leakage tool, the shared phase-coupling tool, the shell-stable readout tool, the obstruction-stratification tool, and the carrier-compression tool from the same bridge packet. Its purpose is to decide where a still-live residual should be reopened first once the total-channel baseline is already frozen.
\end{enumerate}
\end{definition}

\begin{proposition}[First certified submenu reduction routes]
\label{prop:first-certified-submenu-reduction-routes}
Each of the following recurring post-freeze proof obligations admits menu-to-carrier reduction through one of the canonical finite submenus from Definition~\ref{def:first-canonical-finite-submenus-qtool}.
\begin{enumerate}[leftmargin=2em]
  \item Any obligation whose only job is to preserve the retained total-channel profile, carrier-stable pass status, or diagnostic subordination may be routed through \(M^{\mathrm{rig}}_{\mathrm{tot}}\).

  \item Any obligation that accepts the retained total-channel profile as fixed and only seeks to compress the surviving OP~1 burden toward retained-window rigidity or optimum pinning may be routed through \(M^{\mathrm{pin}}_{\mathrm{imb}}\).

  \item Any obligation whose main difficulty is a controlled change of surface --- from channel language to shell, menu, or matrix language, or back again --- may be routed through \(M^{\mathrm{surf}}_{\mathrm{br}}\).

  \item Any obligation whose main difficulty is not proof of a new carrier statement but localization of the first still-live residual interface may be routed through \(M^{\mathrm{loc}}_{\mathrm{br}}\).
\end{enumerate}
In particular, the immediate continuation line now has four explicit reuse routes rather than only one undifferentiated theorem store.
\end{proposition}

\begin{proof}
The total-rigidity submenu is by definition exactly the finite coded portion of \(\mathfrak T_{\mathrm{tot}}\) that already controls transport-first stability, readout-after-transport stability, admissible readout uniqueness, and global diagnostic subordination. Therefore every obligation of type \emph{(1)} is already reduced to the smallest currently relevant carrier-stable total-channel interface.

The optimum-pinning submenu is by definition the finite coded chain inside \(\mathfrak T_{\mathrm{imb}}\) that successively compresses the surviving OP~1 burden from frame split to shell balance, two-pass symmetry, loop budget, admissible-window uniqueness, retained-window rigidity, and finally optimum-pinning comparison-rigidity. Therefore every obligation of type \emph{(2)} is already reduced to the smallest currently relevant imbalance-side compression route.

The surface-change submenu isolates exactly those bridge tools whose role is transport across surfaces rather than new burden creation: matrix leakage, cryptanalytic translation, shared phase coupling, shell-stable readout, and commuting-square compatibility. Hence every obligation of type \emph{(3)} is reduced to the smallest currently relevant change-of-surface route.

Finally, the obstruction-localization submenu isolates the bridge tools whose role is not to certify a new global pass statement but to decide where the first still-live residual now sits: shared leakage/phase diagnostics identify the candidate bridge side, shell-stable readout preserves the common carrier-bearing content where available, obstruction stratification localizes the first failure stratum, and carrier compression collapses the readout/kernel side to one smallest missing carrier package. Hence every obligation of type \emph{(4)} admits menu-to-carrier reduction through \(M^{\mathrm{loc}}_{\mathrm{br}}\). This is exactly the claimed four-route reuse structure.
\end{proof}

\begin{corollary}[Working citation rule for the next immediate revision]
\label{cor:working-citation-rule-next-immediate-revision}
In the next immediate revision, a later proof should normally cite one of the four canonical finite submenus first and only reopen the larger packet description if that submenu fails to carry the obligation. Concretely:
\begin{enumerate}[leftmargin=2em]
  \item cite \(M^{\mathrm{rig}}_{\mathrm{tot}}\) first for total-channel rigidity or status-preservation steps;
  \item cite \(M^{\mathrm{pin}}_{\mathrm{imb}}\) first for OP~1 compression toward the optimum-pinning kernel;
  \item cite \(M^{\mathrm{surf}}_{\mathrm{br}}\) first for channel/shell/menu/matrix translation steps;
  \item cite \(M^{\mathrm{loc}}_{\mathrm{br}}\) first for localization of the first still-live bridge or readout-side obstruction.
\end{enumerate}
Thus the QAM layer has now moved from abstract finite storage to a concrete four-route routing discipline.
\end{corollary}

\begin{remark}[What has been gained before proving anything new]
\label{rem:what-gained-before-proving-anything-new-submenus}
The gain is practical and immediate. A later argument no longer has to say only ``use QAM'' or ``reopen the bridge packet.'' It can now say much more sharply: preserve the frozen total profile via \(M^{\mathrm{rig}}_{\mathrm{tot}}\), compress the surviving OP~1 route via \(M^{\mathrm{pin}}_{\mathrm{imb}}\), switch surfaces via \(M^{\mathrm{surf}}_{\mathrm{br}}\), or localize the first still-live residual via \(M^{\mathrm{loc}}_{\mathrm{br}}\). That is exactly the kind of proof-shortening interface the immediate tool/QAM revision was supposed to extract.
\end{remark}

\subsection*{19AZ$^{(3a)}$. Normalized QAM nomenclature and finite tag layer for the four canonical submenus}
\label{sec:normalized-qam-tag-layer}
The four canonical submenus should now receive a short address language. Without that extra layer, the current QAM revision would still have finite storage and routing, but not yet a stable nomenclature by which later proofs or diagnostics can point to one smallest already-normalized carrier without reopening the full sentence around it.

\begin{definition}[Primary tool-tag families]
\label{def:primary-tool-tag-families}
Read the four canonical finite submenus of Definition~\ref{def:qam-coded-tool-menu} in the order induced by their listed entries. The \emph{primary tool-tag families} are the finite code families
\[
\mathrm{TR}_1,\mathrm{TR}_2,\dots;\qquad
\mathrm{IP}_1,\mathrm{IP}_2,\dots;\qquad
\mathrm{BS}_1,\mathrm{BS}_2,\dots;\qquad
\mathrm{BL}_1,\mathrm{BL}_2,\dots
\]
indexed respectively by the ordered entries of
\[
M^{\mathrm{rig}}_{\mathrm{tot}},\qquad
M^{\mathrm{pin}}_{\mathrm{imb}},\qquad
M^{\mathrm{surf}}_{\mathrm{br}},\qquad
M^{\mathrm{loc}}_{\mathrm{br}}.
\]
Here \(\mathrm{TR}\) stands for \emph{total rigidity}, \(\mathrm{IP}\) for \emph{imbalance pinning}, \(\mathrm{BS}\) for \emph{bridge/surface change}, and \(\mathrm{BL}\) for \emph{bridge/localization}. The ordered disjoint union of these four families is the \emph{primary tool-tag alphabet} of the immediate continuation line.
\end{definition}

\begin{definition}[Primary tool-tag registry and dereferencing map]
\label{def:primary-tool-tag-registry-dereferencing-map}
The \emph{primary tool-tag registry}
\[
\mathcal R^{\mathrm{tag}}_{\mathrm{tool}}
\]
is the finite family of quadruples
\[
(\tau,M,\kappa,x)
\]
with the following meaning:
\begin{enumerate}[leftmargin=2em]
  \item \(\tau\) is one primary tag from Definition~\ref{def:primary-tool-tag-families};
  \item \(M\) is the unique canonical submenu among \(M^{\mathrm{rig}}_{\mathrm{tot}}, M^{\mathrm{pin}}_{\mathrm{imb}}, M^{\mathrm{surf}}_{\mathrm{br}}, M^{\mathrm{loc}}_{\mathrm{br}}\) in which the tagged item has its primary home;
  \item \(\kappa\in\{\mathrm{status},\mathrm{pinning},\mathrm{surface},\mathrm{localization}\}\) is its normalized burden type;
  \item \(x\in M\) is the unique carrier statement, submenu entry, or smallest normalized reusable route addressed by \(\tau\).
\end{enumerate}
The induced map
\[
\mathsf{Ref}:\tau\longmapsto x
\]
is called the \emph{dereferencing map}. If one and the same carrier statement is reusable from more than one submenu, then exactly one of those occurrences is designated as its \emph{primary home}; every other appearance is recorded only as a secondary pointer and receives no second primary tag. The registry order defines a precedence rank
\[
\pi(\tau)\in\mathbb N
\]
inside the chosen submenu, so that smaller rank means earlier reopening inside the already chosen route.
\end{definition}

\begin{proposition}[Tag soundness, uniqueness, and minimal reopening discipline]
\label{prop:tag-soundness-uniqueness-minimal-reopening-discipline}
For every entry \(x\) of the four canonical finite submenus there exists exactly one primary tag \(\tau\) such that
\[
(\tau,M,\kappa,x)\in \mathcal R^{\mathrm{tag}}_{\mathrm{tool}}.
\]
Consequently:
\begin{enumerate}[leftmargin=2em]
  \item the primary tag \(\tau\) determines the canonical submenu \(M\), the normalized burden type \(\kappa\), and the dereferenced carrier \(x\);
  \item no collision between the four nomenclature families can occur at the primary level;
  \item once a later argument has chosen its submenu, reopening should proceed in increasing precedence rank \(\pi(\tau)\) until the first sufficient carrier is reached.
\end{enumerate}
Hence the immediate QAM layer now carries a finite nomenclature with exact dereferencing and a smallest-first reopening discipline.
\end{proposition}

\begin{proof}
The four canonical submenus are finite and ordered by construction. Assigning one tag from the corresponding family to each ordered entry therefore produces exactly one registry quadruple \((\tau,M,\kappa,x)\) for each such entry. Because the tag families \(\{\mathrm{TR}_i\}, \{\mathrm{IP}_j\}, \{\mathrm{BS}_k\}, \{\mathrm{BL}_\ell\}\) are disjoint by notation, primary collisions cannot occur across different submenu types. If one carrier statement appears in more than one submenu role, the definition of primary home forces uniqueness at the registry level by declaring every other occurrence secondary only. Finally, the submenu order itself supplies the precedence rank \(\pi(\tau)\), so once a route has been chosen the smallest-first reopening order is already fixed. This is exactly the claimed tag soundness, uniqueness, and reopening discipline.
\end{proof}

\begin{corollary}[Tag-first citation syntax for the immediate continuation line]
\label{cor:tag-first-citation-syntax-immediate-continuation-line}
The next immediate revision may cite the normalized tool layer in either of the following two equivalent ways:
\begin{enumerate}[leftmargin=2em]
  \item by naming the submenu and then the primary tag, for example ``invoke \(M^{\mathrm{rig}}_{\mathrm{tot}}\) at \(\mathrm{TR}_i\)'' or ``reopen \(M^{\mathrm{loc}}_{\mathrm{br}}\) at \(\mathrm{BL}_\ell\)'';
  \item once the submenu is fixed by context, by citing the primary tag alone and dereferencing it through \(\mathsf{Ref}\).
\end{enumerate}
Accordingly, the four primary tag families should be read as the first normalized QAM nomenclature layer for the immediate continuation line:
\begin{center}
\begin{tabular}{@{}ll@{}}
\toprule
Tag family & Intended immediate use \\
\midrule
\(\mathrm{TR}_i\) & total-channel rigidity and status preservation \\
\(\mathrm{IP}_j\) & OP~1 compression toward the optimum-pinning kernel \\
\(\mathrm{BS}_k\) & channel/shell/menu/matrix surface change \\
\(\mathrm{BL}_\ell\) & localization of the first still-live obstruction \\
\bottomrule
\end{tabular}
\end{center}
Thus the QAM layer now has not only finite menus but also a short stable address language for them.
\end{corollary}

\begin{remark}[What the tag layer buys before any new theorem is added]
\label{rem:what-tag-layer-buys-before-new-theorem}
This is the first genuinely addressable QAM tool surface on the present line. A later proof can now point to a smallest normalized carrier by a short typed tag instead of reopening a long theorem sentence; a later diagnostic routine can record that a residual cluster repeatedly points toward \(\mathrm{IP}_7\) or \(\mathrm{BL}_2\) rather than toward an entire section; and a later public extraction can export the tags as the front-face nomenclature without exporting the whole internal packet architecture. In that precise sense, the tag layer is the missing bridge from finite storage to reusable proof vocabulary.
\end{remark}



\subsection*{19AZ$^{(3b)}$. First explicit primary registry assignments for the normalized tool tags}
\label{sec:first-explicit-primary-registry-assignments}
The normalized tag alphabet should now be instantiated once, so that the immediate continuation line can cite concrete dereferenced carriers instead of only a generic family pattern. At the present stage the assignments are intentionally conservative: each tag points only to a carrier statement or smallest reusable route that already exists in the frozen working text.

\paragraph{First concrete primary registry.}
Write the finite registry in four family blocks. For compactness, the ``carrier'' column records the dereferenced item together with its smallest currently stable reference surface.
\begin{quote}\small
\noindent\textbf{Reader cue.} Read the following registry tables as lookup surfaces for the normalized QAM tags. Their purpose is to pin short addresses to already stabilized statements or reusable micro-routes; they do not introduce new carriers.
\end{quote}

{\small
\renewcommand{\arraystretch}{1.08}
\begin{center}
\begin{tabular}{@{}>{\raggedright\arraybackslash}p{0.10\linewidth}>{\raggedright\arraybackslash}p{0.14\linewidth}>{\raggedright\arraybackslash}p{0.66\linewidth}@{}}
\toprule
Tag & $\kappa$ & Dereferenced carrier \\
\midrule
\multicolumn{3}{@{}l@{}}{\textbf{Total-rigidity submenu} $M^{\mathrm{rig}}_{\mathrm{tot}}$}\\[0.2em]
$\mathrm{TR}_1$ & status & the transport-first route read from the opening clause of the proof of Theorem~\ref{thm:working-op1-admissible-readout-uniqueness}, namely: branch selection is fixed before visible readout is interpreted;\\
$\mathrm{TR}_2$ & status & the readout-after-transport route read from the next clause of the same proof of Theorem~\ref{thm:working-op1-admissible-readout-uniqueness}, namely: every canonical readout remains downstream from the already fixed transport branch;\\
$\mathrm{TR}_3$ & status & the admissible-readout uniqueness carrier of Theorem~\ref{thm:working-op1-admissible-readout-uniqueness};\\
$\mathrm{TR}_4$ & status & the bundled global pass/transfer/diagnostic-subordination interface from Propositions~\ref{prop:global-pass-tool}, \ref{prop:global-transfer-tool}, and \ref{prop:global-diagnostic-subordination-tool}.\\
\bottomrule
\end{tabular}
\end{center}

\begin{center}
\begin{tabular}{@{}>{\raggedright\arraybackslash}p{0.10\linewidth}>{\raggedright\arraybackslash}p{0.14\linewidth}>{\raggedright\arraybackslash}p{0.66\linewidth}@{}}
\toprule
Tag & $\kappa$ & Dereferenced carrier \\
\midrule
\multicolumn{3}{@{}l@{}}{\textbf{Optimum-pinning submenu} $M^{\mathrm{pin}}_{\mathrm{imb}}$}\\[0.2em]
$\mathrm{IP}_1$ & pinning & frame-transport split certificate, Proposition~\ref{prop:op1-frame-transport-split};\\
$\mathrm{IP}_2$ & pinning & finite four-shell certificate, Proposition~\ref{prop:op1-four-shell-certificate};\\
$\mathrm{IP}_3$ & pinning & two-pass half-burden certificate, Proposition~\ref{prop:op1-two-pass-half-burden};\\
$\mathrm{IP}_4$ & pinning & reciprocal pass-symmetry certificate, Proposition~\ref{prop:op1-reciprocal-pass-symmetry};\\
$\mathrm{IP}_5$ & pinning & same-corridor reversal implication, Proposition~\ref{prop:op1-same-corridor-implies-reciprocity};\\
$\mathrm{IP}_6$ & pinning & unit-loop half-burden certificate, Proposition~\ref{prop:op1-unit-loop-half-burden};\\
$\mathrm{IP}_7$ & pinning & one-first-zero loop-budget bridge, Proposition~\ref{prop:op1-one-first-zero-loop-budget};\\
$\mathrm{IP}_8$ & pinning & admissible-window uniqueness bridge, Proposition~\ref{prop:op1-admissible-window-uniqueness};\\
$\mathrm{IP}_9$ & pinning & retained-window-rigidity / no-second-carrier / strict-upper-gap bridge, Proposition~\ref{prop:op1-rigidity-gap-uniqueness-bridge};\\
$\mathrm{IP}_{10}$ & pinning & optimum-pinning comparison-rigidity bridge, Proposition~\ref{prop:op1-optimum-pinning-implies-rigidity}.\\
\bottomrule
\end{tabular}
\end{center}

\begin{center}
\begin{tabular}{@{}>{\raggedright\arraybackslash}p{0.10\linewidth}>{\raggedright\arraybackslash}p{0.14\linewidth}>{\raggedright\arraybackslash}p{0.66\linewidth}@{}}
\toprule
Tag & $\kappa$ & Dereferenced carrier \\
\midrule
\multicolumn{3}{@{}l@{}}{\textbf{Surface-change submenu} $M^{\mathrm{surf}}_{\mathrm{br}}$}\\[0.2em]
$\mathrm{BS}_1$ & surface & shared matrix-leakage surface gate, read through Definition~\ref{def:working-shared-matrix-leakage} together with Proposition~\ref{prop:working-stable-channel-criterion};\\
$\mathrm{BS}_2$ & surface & the cryptanalytic translation route from the shared matrix-leakage / cryptanalytic block;\\
$\mathrm{BS}_3$ & surface & shared phase-coupling surface diagnostic from Definitions~\ref{def:working-shared-phase-traces}--\ref{def:working-common-carrier-diagnostic} together with Propositions~\ref{prop:working-one-mode-image-test} and \ref{prop:working-quadrature-criterion};\\
$\mathrm{BS}_4$ & surface & the shell-stable readout tool from the extracted OP~3/QAM tool layer;\\
$\mathrm{BS}_5$ & surface & the commuting-square tool from the extracted OP~3/QAM tool layer.\\
\bottomrule
\end{tabular}
\end{center}

\begin{center}
\begin{tabular}{@{}>{\raggedright\arraybackslash}p{0.10\linewidth}>{\raggedright\arraybackslash}p{0.14\linewidth}>{\raggedright\arraybackslash}p{0.66\linewidth}@{}}
\toprule
Tag & $\kappa$ & Dereferenced carrier \\
\midrule
\multicolumn{3}{@{}l@{}}{\textbf{Obstruction-localization submenu} $M^{\mathrm{loc}}_{\mathrm{br}}$}\\[0.2em]
$\mathrm{BL}_1$ & localization & matrix-leakage-first localization gate toward the underfixed HC/MM/J coupling, read again through Definition~\ref{def:working-shared-matrix-leakage} and Proposition~\ref{prop:working-stable-channel-criterion};\\
$\mathrm{BL}_2$ & localization & phase-coupling-first localization gate asking whether the residual pair is still one masked carrier, read through Propositions~\ref{prop:working-one-mode-image-test} and \ref{prop:working-quadrature-criterion};\\
$\mathrm{BL}_3$ & localization & the shell-stable readout localization gate from the extracted OP~3/QAM tool layer;\\
$\mathrm{BL}_4$ & localization & the obstruction-stratification tool from the extracted OP~3/QAM tool layer;\\
$\mathrm{BL}_5$ & localization & the carrier-compression tool from the extracted OP~3/QAM tool layer.\\
\bottomrule
\end{tabular}
\end{center}
}

\paragraph{Primary-home convention at the concrete level.}
The bridge families contain repeated antecedent references on purpose. This does not violate the primary-home rule from Definition~\ref{def:primary-tool-tag-registry-dereferencing-map}, because the registry items above are not identified merely by their antecedent theorem labels, but by their \emph{normalized route role}. Thus the same matrix-leakage or shell-stable-readout material may appear once as a \emph{surface-change gate} and once as a \emph{localization gate} without creating a primary collision.

\paragraph{Immediate working use.}
From this point onward the continuation line may cite concrete tags without any placeholder notation. For example, ``freeze the total channel at $\mathrm{TR}_4$'', ``compress the surviving OP~1 kernel through $\mathrm{IP}_8\to \mathrm{IP}_{10}$'', ``change surface via $\mathrm{BS}_2$'', or ``localize the first still-live residual by testing $\mathrm{BL}_1$, then $\mathrm{BL}_2$, then $\mathrm{BL}_4$'' are now literal dereferenceable instructions on the QAM side.


\subsection*{19AZ$^{(3c)}$. First tag-first rewrite of the surviving OP~1 route}
\label{sec:first-tag-first-rewrite-of-the-surviving-op1-route}
The normalized registry should now be used once on a real proof route rather than remaining only a citation promise. The cleanest first test is the currently surviving OP~1 compression chain, because it already has a sharply localized terminal burden and because its hypotheses were precisely the ones that motivated the split into total, pinning, and bridge roles.

\begin{proposition}[First tag-first reformulation of the surviving OP~1 closure route]
\label{prop:first-tag-first-op1-route}
Assume the same-corridor reversal certificate of Definition~\ref{def:op1-same-corridor-reversal}, and assume moreover that the optimum-pinning comparison-rigidity certificate of Definition~\ref{def:op1-optimum-pinning-comparison-rigidity} holds on the intended admissible windows. Then the surviving OP~1 closure route may be cited on the normalized QAM side in the compact form
\[
  \mathrm{TR}_4 \;\Longrightarrow\; \mathrm{IP}_{10} \;\Longrightarrow\; \mathrm{IP}_9 \;\Longrightarrow\; \mathrm{IP}_8,
\]
and therefore
\[
  \Delta_{\mathrm{CL}} = 0.
\]
Equivalently, once the retained total-channel baseline and its bundled transfer/diagnostic discipline are frozen through \(\mathrm{TR}_4\), the live OP~1 burden is legitimately reopened only as the short pinning chain
\[
  \mathrm{IP}_{10} \to \mathrm{IP}_9 \to \mathrm{IP}_8,
\]
rather than as a return to the full historical OP~1 queue.
\end{proposition}

\begin{proof}
By the concrete registry from Section~\ref{sec:first-explicit-primary-registry-assignments}, the tag \(\mathrm{TR}_4\) dereferences the bundled global pass/transfer/diagnostic-subordination interface from Propositions~\ref{prop:global-pass-tool}, \ref{prop:global-transfer-tool}, and \ref{prop:global-diagnostic-subordination-tool}. On the present line this is exactly the smallest normalized total-channel surface that keeps fixed the retained total profile and the comparison/diagnostic discipline under which the OP~1 corridor is being read.

Under the standing same-corridor reversal hypothesis, the additional assumption of optimum-pinning comparison rigidity is precisely the carrier tagged by \(\mathrm{IP}_{10}\), namely Proposition~\ref{prop:op1-optimum-pinning-implies-rigidity}. That proposition says that, on the intended admissible windows, preservation of the retained total-channel profile together with the retained optimum location and local comparison/defect data forces retained-window rigidity.

The next normalized reopening step is then \(\mathrm{IP}_9\), i.e. Proposition~\ref{prop:op1-rigidity-gap-uniqueness-bridge}. Once retained-window rigidity is available, and with the standing no-second-carrier and strict-upper-gap exclusions built into that bridge, admissible-window uniqueness follows on the retained corridor.

Finally one reopens \(\mathrm{IP}_8\), i.e. Proposition~\ref{prop:op1-admissible-window-uniqueness}. This upgrades admissible-window uniqueness to the one-first-zero loop-budget certificate and hence yields
\[
  \Delta_{\mathrm{CL}} = 0.
\]
No new mathematics has been added here: the content is exactly the already established OP~1 compression chain, now rewritten in normalized tag-first form. The gain is that the route can be cited as one frozen total-channel anchor plus three smallest pinning tags instead of reopening the larger sentence-level proof history.
\end{proof}

\begin{remark}[What this first practical rewrite already buys]
\label{rem:what-the-first-practical-tag-rewrite-buys}
This is the first point at which the tag layer stops being merely architectural. A later continuation can now say, in a mathematically literal way, ``freeze at \(\mathrm{TR}_4\), then test \(\mathrm{IP}_{10}\to\mathrm{IP}_9\to\mathrm{IP}_8\)'' and thereby refer to an already normalized proof route rather than to an entire subsection block. In other words, the QAM layer has started to function as a true routing vocabulary for the live burden, not just as a storage format for tools.
\end{remark}


\subsection*{19AZ$^{(3d)}$. First tag-first rewrite of the bridge/localization route}
\label{sec:first-tag-first-bridge-localization-route}
The second practical test of the normalized QAM registry should not pretend to close a new theorematic burden. Its job is narrower and more honest: once the retained total-channel baseline is frozen, it should rewrite the currently live bridge-side reopening order into one compact route that first changes surface at the smallest matrix/phase interface and only then localizes the first still-live obstruction.

\begin{proposition}[First tag-first reformulation of the bridge/localization reopening route]
\label{prop:first-tag-first-bridge-localization-route}
Assume the retained total-channel baseline is frozen through \(\mathrm{TR}_4\), and suppose a still-live bridge-side residual is to be tested without yet postulating two independent primitive carriers. Then the first normalized QAM reopening route may be cited in the compact form
\[
  \mathrm{TR}_4 \;\Longrightarrow\; \mathrm{BS}_1 \;\Longrightarrow\; \mathrm{BS}_3 \;\Longrightarrow\; \mathrm{BL}_2 \;\Longrightarrow\; \mathrm{BL}_4,
\]
with the additional compression step
\[
  \mathrm{BL}_4 \;\Longrightarrow\; \mathrm{BL}_5
\]
whenever the first localized obstruction lies in the readout/kernel stratum.
Equivalently, once the retained total profile and its bundled comparison/diagnostic discipline are frozen, the live bridge burden is reopened first as matrix leakage, then as common-carrier/phase coupling, and only then as stratified obstruction localization; a separate carrier is not introduced before this normalized route has been exhausted.
\end{proposition}

\begin{proof}
By the concrete registry from Section~\ref{sec:first-explicit-primary-registry-assignments}, the tag \(\mathrm{TR}_4\) dereferences the bundled global pass/transfer/diagnostic-subordination interface from Propositions~\ref{prop:global-pass-tool}, \ref{prop:global-transfer-tool}, and \ref{prop:global-diagnostic-subordination-tool}. On the present line this is again the smallest frozen total-channel surface from which later local re-openings are allowed.

The first bridge-side reopening step is \(\mathrm{BS}_1\), i.e. the shared matrix-leakage surface gate from Definition~\ref{def:working-shared-matrix-leakage} together with Proposition~\ref{prop:working-stable-channel-criterion}. This is the right first test because it asks whether there is still matrix-induced off-sector drift at all. In other words, before any finer localization language is used, the route first checks whether the admissible HC channel is already stable under the Millennium/Umkehrwalze interface.

If the bridge burden is still live after that first gate, then one changes surface through \(\mathrm{BS}_3\), namely the shared phase-coupling diagnostic from Definitions~\ref{def:working-shared-phase-traces}--\ref{def:working-common-carrier-diagnostic} together with Propositions~\ref{prop:working-one-mode-image-test} and \ref{prop:working-quadrature-criterion}. This does not yet decide where the first remaining obstruction sits; it rewrites the visible residual pair into the smallest matrix/phase language in which one can honestly ask whether the pair is still one masked carrier rather than two unrelated end defects.

That localization question is then reopened directly as \(\mathrm{BL}_2\), which is precisely the phase-coupling-first localization gate recorded in the same registry. If the residual pair still behaves as one masked carrier, the route stays on the one-carrier side. If it remains unresolved or fails at that stage, the next legitimate reopening is \(\mathrm{BL}_4\), i.e. the obstruction-stratification tool from the extracted OP~3/QAM layer. This is the first point at which the present line is allowed to ask whether the surviving obstruction lives at the branch, shell, or readout/kernel level.

Finally, when \(\mathrm{BL}_4\) places the first residual in the readout/kernel stratum, the carrier-compression tool \(\mathrm{BL}_5\) compresses that stratum to absence of one common shell-stable carrier. Hence the entire normalized reopening order is exactly the displayed tag route. No new closure theorem has been proved here. The content is the already extracted bridge packet and the already extracted localization packet, now rewritten as one smallest admissible tag-first diagnostic route.
\end{proof}

\begin{remark}[What this second practical rewrite already buys]
\label{rem:what-the-second-practical-tag-rewrite-buys}
The present line can now cite both kinds of live work in the same compact language. For OP~1 compression one may say ``freeze at \(\mathrm{TR}_4\), then test \(\mathrm{IP}_{10}\to\mathrm{IP}_9\to\mathrm{IP}_8\)''; for bridge/localization one may now say ``freeze at \(\mathrm{TR}_4\), then test \(\mathrm{BS}_1\to\mathrm{BS}_3\), then localize by \(\mathrm{BL}_2\to\mathrm{BL}_4\), and if necessary compress by \(\mathrm{BL}_5\).'' In that precise sense, the QAM layer is no longer only a menu of stored tools: it has become a two-sided routing vocabulary, one side for the surviving OP~1 pinning chain and one side for the still-live bridge/localization diagnostics.
\end{remark}

\subsection*{19AZ$^{(4)}$. First Monte-Carlo residual-signature route beneath the normalized QAM tool packets}
\label{sec:first-monte-carlo-residual-signature-route}
The Monte-Carlo programme should therefore enter only \emph{after} the normalized ledger and its QAM recoding are in place. On the present line, Monte Carlo is not a replacement proof engine. It is a structural probe that samples how residual deviations distribute relative to the already fixed total-channel baseline and to the still-live imbalance/bridge packets.

\begin{definition}[Residual-signature word]
\label{def:residual-signature-word}
Fix a retained admissible window and a finite family of normalized diagnostics
\[
D=(D_1,\dots,D_m)
\]
measured relative to the retained total-channel baseline. For one Monte-Carlo sample \(\omega\), let
\[
\delta(\omega)=(\delta_1(\omega),\dots,\delta_m(\omega))
\]
be the resulting normalized deviation vector. A \emph{residual-signature word} for \(\omega\) is any finite coded word
\[
\mathsf{Sig}(\omega;D)
\]
obtained from the ordered sign, bin, rank, or slot pattern of the components of \(\delta(\omega)\). Such a word is interpreted as a structural signature or hash of the sampled residual behavior, not as a theorem by itself.
\end{definition}

\begin{definition}[Kasiski-style recurrence scan]
\label{def:kasiski-style-recurrence-scan}
Given a family of residual-signature words, the \emph{Kasiski-style recurrence scan} is the first combinatorial pass that looks for repeated blocks, repeated slot offsets, and stable periodicity or near-periodicity classes among those words. Its role is diagnostic only: it asks whether the sampled residuals point repeatedly toward the same bridge or imbalance surface.
\end{definition}

\begin{proposition}[Conservative Monte-Carlo/QAM probe rule]
\label{prop:conservative-monte-carlo-qam-probe-rule}
Residual-signature words and their Kasiski-style recurrence scan may be used in the immediate continuation line only to prioritize which normalized packet from Definition~\ref{def:normalized-channel-sorted-tool-packets} should be reopened first. They may \emph{not} by themselves introduce a second primitive carrier, overturn the frozen total-channel baseline, or replace an already theorematic pass/obstruction statement.
\end{proposition}

\begin{proof}
This is exactly the methodological discipline stated in the deferred roadmap above. The retained total-channel baseline is already part of the frozen theorem layer, while the Monte-Carlo route is introduced only as a structural probe beneath that layer. Therefore sampled residual signatures may guide priority and localization, but they cannot overrule the already established carrier and status surfaces.
\end{proof}

\begin{remark}[Immediate diagnostic reading of the first residual signatures]
\label{rem:immediate-diagnostic-reading-signatures}
The intended first reading is now precise. If the residual signatures cluster while the retained total-channel profile stays rigid, then the next reopen candidate should be searched first inside \(\mathfrak T_{\mathrm{imb}}\) or \(\mathfrak T_{\mathrm{br}}\), depending on whether the recurrence is purely local/imbalance-side or genuinely translation-sensitive. If no stable recurrence class appears at all, then the theorem layer should remain untouched and the Monte-Carlo route should be treated as presently non-informative rather than overinterpreted.
\end{remark}


\subsection*{19AZ$^{(4a)}$. First decoding of residual-signature recurrences into normalized BS/BL reopening priorities}
\label{sec:first-decoding-of-signature-recurrences-into-bs-bl-priorities}
The next conservative step is now to make the bridge/localization reading operational. The residual-signature words should not merely say that ``something bridge-like remains.'' They should decode into a smallest prioritized reopening order. In particular, the present line needs one first rule saying when a stable recurrence class should push the next diagnostic pass onto the normalized bridge/localization route rather than back into the OP~1 pinning chain.

\begin{definition}[Route-priority decoder for residual signatures]
\label{def:route-priority-decoder-residual-signatures}
Let \(\mathcal C\) be a stable recurrence class produced by the Kasiski-style scan from Definition~\ref{def:kasiski-style-recurrence-scan}. A \emph{route-priority decoder} for \(\mathcal C\) is any conservative rule
\[
\mathsf{Dec}(\mathcal C)\in \bigl\{M^{\mathrm{pin}}_{\mathrm{imb}},\,M^{\mathrm{surf}}_{\mathrm{br}},\,M^{\mathrm{loc}}_{\mathrm{br}}\bigr\}
\]
that assigns \(\mathcal C\) to the smallest normalized menu whose recorded carriers already account for the repeated slot/rank pattern of the sampled residuals.
The decoder is called \emph{translation-sensitive} if the repeated pattern is not stable under purely local imbalance-side rereading, but stabilizes once the residual pair is reread through the matrix/phase common-carrier surface. In that case the first prioritized reopening route attached to \(\mathcal C\) is
\[
\mathrm{BS}_1 \to \mathrm{BS}_3 \to \mathrm{BL}_2,
\]
with \(\mathrm{BL}_4\) and, when needed, \(\mathrm{BL}_5\) appended only after the phase/common-carrier gate has been exhausted.
\end{definition}

\begin{proposition}[Translation-sensitive signature classes prioritize the BS/BL route]
\label{prop:translation-sensitive-signature-classes-prioritize-bs-bl-route}
Assume the retained total-channel baseline remains frozen through \(\mathrm{TR}_4\), and let \(\mathcal C\) be a stable recurrence class of residual-signature words. Suppose moreover that \(\mathcal C\) is translation-sensitive in the sense of Definition~\ref{def:route-priority-decoder-residual-signatures}: its repeated slot/rank pattern does not settle under a purely local imbalance-side rereading, but does settle once the same sampled residuals are reread through the matrix/phase common-carrier surface. Then the first normalized reopening priority attached to \(\mathcal C\) is the bridge/localization route
\[
\mathrm{TR}_4 \;\Longrightarrow\; \mathrm{BS}_1 \;\Longrightarrow\; \mathrm{BS}_3 \;\Longrightarrow\; \mathrm{BL}_2,
\]
followed, only if a still-live obstruction remains, by
\[
\mathrm{BL}_2 \;\Longrightarrow\; \mathrm{BL}_4,
\]
and by
\[
\mathrm{BL}_4 \;\Longrightarrow\; \mathrm{BL}_5
\]
whenever the first localized obstruction lies in the readout/kernel stratum.
Equivalently, a stable translation-sensitive recurrence class is not yet evidence for a second primitive carrier; it is evidence that the next reopen candidate should be searched first along the normalized matrix/phase-to-localization bridge route.
\end{proposition}

\begin{proof}
By Proposition~\ref{prop:conservative-monte-carlo-qam-probe-rule}, residual-signature words may only prioritize a reopening order inside the already normalized theorem/tool layer; they may not overrule the frozen total-channel baseline or create a new carrier hypothesis. Hence one must begin from \(\mathrm{TR}_4\), the smallest bundled total-channel anchor already fixed by the registry from Section~\ref{sec:first-explicit-primary-registry-assignments}.

Now assume the stable recurrence class \(\mathcal C\) is translation-sensitive. By definition this means that the repeated residual pattern does not become stable when reread only as a local imbalance perturbation, but does become stable once the same samples are reread through the matrix/phase common-carrier surface. Therefore the first admissible decoder target for \(\mathcal C\) cannot be merely \(M^{\mathrm{pin}}_{\mathrm{imb}}\); the smallest normalized reopening surface that can already account for the stable repeated pattern is the bridge-surface menu. In the concrete registry this is exactly the route starting at \(\mathrm{BS}_1\), the shared matrix-leakage gate, and continuing through \(\mathrm{BS}_3\), the shared phase/common-carrier surface.

Once that surface change has been made, the first honest question is whether the repeated residual pair still behaves as one masked carrier. In the normalized registry this is precisely the localization gate \(\mathrm{BL}_2\). If the recurrence class is already resolved there, no deeper reopening is justified. If it remains live, then the first further admissible step is \(\mathrm{BL}_4\), the obstruction-stratification tool. And if that stratification places the first obstruction in the readout/kernel stratum, then \(\mathrm{BL}_5\) is exactly the recorded compression step that rewrites the surviving burden as absence of one common shell-stable carrier package. Thus the displayed BS/BL priority route is exactly the smallest normalized decoding of a translation-sensitive stable recurrence class.
\end{proof}

\begin{remark}[What the first Monte-Carlo-to-route decoding now buys]
\label{rem:what-the-first-monte-carlo-to-route-decoding-buys}
The Monte-Carlo layer now has a literal output format inside the QAM vocabulary. A stable recurrence class is no longer only a vague hint that ``the bridge side may matter.'' The present line may now state, in compact normalized form: ``freeze at \(\mathrm{TR}_4\); if the recurrence is translation-sensitive, decode first by \(\mathrm{BS}_1\to\mathrm{BS}_3\to\mathrm{BL}_2\), and only then localize further by \(\mathrm{BL}_4\) or compress by \(\mathrm{BL}_5\).'' In that precise sense the residual-signature words have acquired a conservative but genuine routing semantics.
\end{remark}



\subsection*{19AZ$^{(4b)}$. First decoding of purely local residual-signature recurrences into normalized IP reopening priorities}
\label{sec:first-decoding-of-local-signature-recurrences-into-ip-priorities}
The decoder should now be made symmetric on the pinning side. The residual-signature words should not only tell us when to reopen the bridge/localization route; they should also tell us when the next honest move remains entirely inside the surviving OP~1 kernel. The immediate line therefore needs one first rule saying when a stable recurrence class should stay on the local imbalance side and decode directly into the normalized IP route.

\begin{definition}[Local-imbalance decoder for residual signatures]
\label{def:local-imbalance-decoder-residual-signatures}
Let \(\mathcal C\) be a stable recurrence class produced by the Kasiski-style scan from Definition~\ref{def:kasiski-style-recurrence-scan}. A route-priority decoder from Definition~\ref{def:route-priority-decoder-residual-signatures} is called \emph{purely local} or \emph{imbalance-stable} on \(\mathcal C\) if the repeated slot/rank pattern already stabilizes under purely local rereading on the imbalance side, so that rereading the same samples through the matrix/phase common-carrier surface adds no smaller admissible explanation. In that case the first prioritized reopening route attached to \(\mathcal C\) is the normalized pinning chain
\[
\mathrm{IP}_{10} \to \mathrm{IP}_9 \to \mathrm{IP}_8,
\]
read beneath the frozen total-channel anchor \(\mathrm{TR}_4\).
\end{definition}

\begin{proposition}[Purely local signature classes prioritize the IP route]
\label{prop:purely-local-signature-classes-prioritize-ip-route}
Assume the retained total-channel baseline remains frozen through \(\mathrm{TR}_4\), and let \(\mathcal C\) be a stable recurrence class of residual-signature words. Suppose moreover that \(\mathcal C\) is purely local in the sense of Definition~\ref{def:local-imbalance-decoder-residual-signatures}: its repeated slot/rank pattern already stabilizes under a local imbalance-side rereading and does not require an earlier matrix/phase surface change in order to become coherent. Then the first normalized reopening priority attached to \(\mathcal C\) is the pinning route
\[
\mathrm{TR}_4 \;\Longrightarrow\; \mathrm{IP}_{10} \;\Longrightarrow\; \mathrm{IP}_9 \;\Longrightarrow\; \mathrm{IP}_8,
\]
so that the sampled recurrence is read first as a compressed OP~1 burden rather than as a bridge/localization burden.
Equivalently, a stable purely local recurrence class is not evidence that one should leave the imbalance side; it is evidence that the next reopen candidate should be searched first inside the surviving normalized pinning chain.
\end{proposition}

\begin{proof}
By Proposition~\ref{prop:conservative-monte-carlo-qam-probe-rule}, the recurrence class \(\mathcal C\) may only prioritize a route already present in the normalized theorem/tool layer. Hence the total-channel baseline must remain frozen at \(\mathrm{TR}_4\), exactly as in the bridge/localization decoder.

Now assume \(\mathcal C\) is purely local. By definition, its repeated slot/rank pattern already becomes coherent under a local imbalance-side rereading, and no earlier matrix/phase reinterpretation yields a smaller admissible menu. Therefore the decoder is not allowed to jump first to \(M^{\mathrm{surf}}_{\mathrm{br}}\) or \(M^{\mathrm{loc}}_{\mathrm{br}}\). The smallest normalized reopening surface that can already account for the stable recurrence class is instead the surviving pinning menu \(M^{\mathrm{pin}}_{\mathrm{imb}}\).

Inside that menu the current live OP~1 burden has already been compressed in Proposition~\ref{prop:first-tag-first-op1-route}. The first admissible entry is therefore \(\mathrm{IP}_{10}\), the optimum-pinning comparison-rigidity bridge. Once that bridge is opened, one proceeds to \(\mathrm{IP}_9\), the retained-window-rigidity / no-second-carrier / strict-upper-gap bridge, and then to \(\mathrm{IP}_8\), the admissible-window uniqueness bridge. Thus the displayed tag route is exactly the smallest normalized decoding of a stable purely local recurrence class.
\end{proof}

\begin{remark}[What the first Monte-Carlo-to-IP decoding now buys]
\label{rem:what-the-first-monte-carlo-to-ip-decoding-buys}
The decoder is now symmetric across the two live diagnostic sides. The present line may state in compact normalized form: ``freeze at \(\mathrm{TR}_4\); if the recurrence is translation-sensitive, decode first by \(\mathrm{BS}_1\to\mathrm{BS}_3\to\mathrm{BL}_2\); if the recurrence is already imbalance-stable, decode instead by \(\mathrm{IP}_{10}\to\mathrm{IP}_9\to\mathrm{IP}_8\).'' In that precise sense the Monte-Carlo layer no longer only points vaguely toward ``some residual side''; it now separates, conservatively and in dereferenceable tags, the first bridge reopening from the first compressed OP~1 reopening.
\end{remark}


\subsection*{19AZ$^{(4c)}$. First explicit decision table for residual-signature routing}
\label{sec:first-explicit-decision-table-residual-signature-routing}
The decoder layer should now be made directly usable. The present line no longer needs only two separate propositions saying when a stable recurrence class goes to the bridge/localization side or stays on the pinning side. It now needs one explicit decision table that turns the first observed residual-signature class into a smallest admissible action. In particular, the table should also record the conservative ``hold'' outcome when no stable recurrence class appears at all.

\begin{quote}\small
\noindent\textbf{Reader cue.} Read the next table as an action selector. Its purpose is to turn an observed residual-signature class into the smallest admissible reopen instruction, not to claim that every downstream route has already been exhausted.
\end{quote}

\begin{definition}[First explicit decision table for residual-signature routing]
\label{def:first-explicit-decision-table-residual-signature-routing}
Assume the retained total-channel baseline remains frozen through \(\mathrm{TR}_4\), and let \(\mathcal C\) be the first sampled recurrence class produced by the Kasiski-style scan from Definition~\ref{def:kasiski-style-recurrence-scan}. The \emph{first explicit decision table for residual-signature routing} is the map
\[
\mathsf{RouteTab}(\mathcal C)\in \bigl\{\mathsf{Hold},\,\mathsf{IP\text{-}first},\,\mathsf{BS/BL\text{-}first},\,\mathsf{BL\text{-}continue},\,\mathsf{BL\text{-}kernel}\bigr\}
\]
with the following normal form.

\begin{center}
\renewcommand{\arraystretch}{1.2}
\begin{tabular}{>{\raggedright\arraybackslash}p{0.26\linewidth} >{\raggedright\arraybackslash}p{0.22\linewidth} >{\raggedright\arraybackslash}p{0.38\linewidth}}
\toprule
Observed residual-signature class & Decoder output & First normalized route \\
\midrule
No stable recurrence class appears & \(\mathsf{Hold}\) & no reopen; keep the theorem/tool layer frozen \\
Purely local / imbalance-stable recurrence & \(\mathsf{IP\text{-}first}\) & \(\mathrm{TR}_4 \Rightarrow \mathrm{IP}_{10} \Rightarrow \mathrm{IP}_9 \Rightarrow \mathrm{IP}_8\) \\
Translation-sensitive recurrence & \(\mathsf{BS/BL\text{-}first}\) & \(\mathrm{TR}_4 \Rightarrow \mathrm{BS}_1 \Rightarrow \mathrm{BS}_3 \Rightarrow \mathrm{BL}_2\) \\
Translation-sensitive recurrence still live after \(\mathrm{BL}_2\) & \(\mathsf{BL\text{-}continue}\) & append \(\mathrm{BL}_4\) \\
First localized obstruction lies in the readout/kernel stratum & \(\mathsf{BL\text{-}kernel}\) & append \(\mathrm{BL}_5\) \\
\bottomrule
\end{tabular}
\end{center}

The table is read top-down and always returns the smallest normalized action whose recorded carriers already account for the observed recurrence behavior.
\end{definition}

\begin{proposition}[The first explicit decision table is conservative and minimal on the current live classes]
\label{prop:first-explicit-decision-table-conservative-and-minimal}
Under the current theorem/tool layer, every sampled residual-signature class that is presently admitted by Propositions~\ref{prop:translation-sensitive-signature-classes-prioritize-bs-bl-route} and \ref{prop:purely-local-signature-classes-prioritize-ip-route} receives from Definition~\ref{def:first-explicit-decision-table-residual-signature-routing} a unique first admissible action. More precisely:
\begin{enumerate}[label=(\roman*)]
    \item if no stable recurrence class appears, the decision table returns \(\mathsf{Hold}\);
    \item if the recurrence class is purely local, it returns \(\mathsf{IP\text{-}first}\);
    \item if the recurrence class is translation-sensitive, it returns \(\mathsf{BS/BL\text{-}first}\);
    \item deeper localization steps \(\mathsf{BL\text{-}continue}\) and \(\mathsf{BL\text{-}kernel}\) may be appended only after the route has already entered the BS/BL side through \(\mathrm{BL}_2\).
\end{enumerate}
Hence the table is conservative, because it never manufactures a new carrier or bypasses the frozen total-channel anchor, and minimal, because it always picks the smallest currently recorded reopening route compatible with the observed recurrence class.
\end{proposition}

\begin{proof}
If no stable recurrence class appears at all, then Remark~\ref{rem:immediate-diagnostic-reading-signatures} already says that the theorem layer must remain untouched; this is exactly the output \(\mathsf{Hold}\). If a stable recurrence class is purely local, then Proposition~\ref{prop:purely-local-signature-classes-prioritize-ip-route} forces the first admissible reopening route to stay inside the pinning chain beneath \(\mathrm{TR}_4\), so the table must return \(\mathsf{IP\text{-}first}\) and the route \(\mathrm{TR}_4 \Rightarrow \mathrm{IP}_{10} \Rightarrow \mathrm{IP}_9 \Rightarrow \mathrm{IP}_8\). If instead the stable recurrence class is translation-sensitive, then Proposition~\ref{prop:translation-sensitive-signature-classes-prioritize-bs-bl-route} forces the first admissible reopening route to pass through \(\mathrm{BS}_1 \Rightarrow \mathrm{BS}_3 \Rightarrow \mathrm{BL}_2\), so the table must return \(\mathsf{BS/BL\text{-}first}\).

The remaining two outputs are not independent first moves. By Proposition~\ref{prop:translation-sensitive-signature-classes-prioritize-bs-bl-route}, \(\mathrm{BL}_4\) is appended only if the obstruction remains live after \(\mathrm{BL}_2\), and \(\mathrm{BL}_5\) is appended only when that localized obstruction lies in the readout/kernel stratum. Therefore the decision table never jumps directly to \(\mathrm{BL}_4\) or \(\mathrm{BL}_5\) without first exhausting the smaller admissible bridge/localization entry route. This proves both conservativity and minimality on the currently live recurrence classes.
\end{proof}

\begin{remark}[Immediate operational reading of the decision table]
\label{rem:immediate-operational-reading-of-the-decision-table}
The present line can now state the decoder in one sentence that is usable both by a reader and by the future machine layer: classify the first stable residual-signature class, then apply
\[
\mathsf{Hold}\;\vee\;\mathsf{IP\text{-}first}\;\vee\;\mathsf{BS/BL\text{-}first},
\]
with deeper localization allowed only on the BS/BL side by appending \(\mathsf{BL\text{-}continue}\) and, in the readout/kernel case, \(\mathsf{BL\text{-}kernel}\). In that sense the Monte-Carlo diagnostics, the QAM tag layer, and the normalized reopening routes now meet in one explicit working table rather than only in scattered prose.
\end{remark}


\subsection*{19AZ$^{(4d)}$. First deterministic QAM decoder syntax for residual-signature routing}
\label{sec:first-deterministic-qam-decoder-syntax-residual-signature-routing}
The decision table may now be compressed one step further. The present line no longer needs only a prose table saying which class goes to hold, IP-first, or BS/BL-first. It may now record one short deterministic decoder syntax whose words already match the normalized QAM routes. The point is not to introduce a new theory of automata; it is only to give the future machine layer and the human reader the same finite, dereferenceable instruction string.

\begin{definition}[First deterministic QAM decoder syntax]
\label{def:first-deterministic-qam-decoder-syntax}
Assume the retained total-channel anchor has been frozen through \(\mathrm{TR}_4\), and let \(\mathcal C\) be the first sampled residual-signature class from the Kasiski-style recurrence scan. The \emph{first deterministic QAM decoder syntax} is the finite rewrite alphabet
\[
\mathfrak A_{\mathrm{dec}}=\bigl\{\mathsf{classify},\,\mathsf{hold},\,\mathsf{ip},\,\mathsf{bsbl},\,\mathsf{continue},\,\mathsf{kernel}\bigr\}
\]
with the following admissible output words:
\[
\begin{aligned}
\mathsf{classify}&\to\mathsf{hold},\\
\mathsf{classify}&\to\mathsf{ip},\\
\mathsf{classify}&\to\mathsf{bsbl},\\
\mathsf{classify}&\to\mathsf{bsbl}\to\mathsf{continue},\\
\mathsf{classify}&\to\mathsf{bsbl}\to\mathsf{continue}\to\mathsf{kernel}.
\end{aligned}
\]
These words are interpreted by the normalized dereferencing map
\[
\mathsf{Dec}_{\mathrm{QAM}}:
\begin{cases}
\mathsf{hold} &\mapsto \text{keep the theorem/tool layer frozen},\\[0.3em]
\mathsf{ip} &\mapsto \mathrm{TR}_4 \Rightarrow \mathrm{IP}_{10} \Rightarrow \mathrm{IP}_9 \Rightarrow \mathrm{IP}_8,\\[0.3em]
\mathsf{bsbl} &\mapsto \mathrm{TR}_4 \Rightarrow \mathrm{BS}_1 \Rightarrow \mathrm{BS}_3 \Rightarrow \mathrm{BL}_2,\\[0.3em]
\mathsf{continue} &\mapsto \text{append }\mathrm{BL}_4,\\[0.3em]
\mathsf{kernel} &\mapsto \text{append }\mathrm{BL}_5.
\end{cases}
\]
The syntax is called \emph{deterministic} when exactly one of the following first outputs is assigned to \(\mathcal C\): \(\mathsf{hold}\), \(\mathsf{ip}\), or \(\mathsf{bsbl}\); the symbols \(\mathsf{continue}\) and \(\mathsf{kernel}\) may occur only after \(\mathsf{bsbl}\).
\end{definition}

\begin{proposition}[The deterministic decoder syntax is equivalent to the current decision table]
\label{prop:deterministic-decoder-syntax-equivalent-current-decision-table}
Under the current live theorem/tool layer, Definition~\ref{def:first-deterministic-qam-decoder-syntax} yields the same first admissible routing data as Definition~\ref{def:first-explicit-decision-table-residual-signature-routing}. More precisely:
\begin{enumerate}[label=(\roman*)]
    \item if no stable recurrence class appears, the unique admissible output word is \(\mathsf{classify}\to\mathsf{hold}\);
    \item if the recurrence class is purely local / imbalance-stable, the unique admissible output word is \(\mathsf{classify}\to\mathsf{ip}\);
    \item if the recurrence class is translation-sensitive, the unique admissible output word is \(\mathsf{classify}\to\mathsf{bsbl}\);
    \item if, after entering the BS/BL side, the obstruction remains live beyond \(\mathrm{BL}_2\), the decoder extends uniquely to \(\mathsf{classify}\to\mathsf{bsbl}\to\mathsf{continue}\);
    \item if that localized obstruction lies in the readout/kernel stratum, the decoder extends further uniquely to \(\mathsf{classify}\to\mathsf{bsbl}\to\mathsf{continue}\to\mathsf{kernel}\).
\end{enumerate}
Hence the syntax is a faithful finite normal-form compression of the decision table, with no new carriers and no loss of routing information.
\end{proposition}

\begin{proof}
The first three clauses are exactly the three primary outputs of Definition~\ref{def:first-explicit-decision-table-residual-signature-routing}: \(\mathsf{Hold}\), \(\mathsf{IP\text{-}first}\), and \(\mathsf{BS/BL\text{-}first}\). Definition~\ref{def:first-deterministic-qam-decoder-syntax} merely rewrites these three outcomes as the shorter syntax outputs \(\mathsf{hold}\), \(\mathsf{ip}\), and \(\mathsf{bsbl}\), while the dereferencing map \(\mathsf{Dec}_{\mathrm{QAM}}\) sends them back to the same normalized routes.

The remaining two clauses are likewise only compressed spellings of the already recorded BS/BL-side continuations. By Proposition~\ref{prop:first-explicit-decision-table-conservative-and-minimal}, the symbol \(\mathsf{BL\text{-}continue}\) may appear only after the bridge/localization route has already entered through \(\mathrm{BL}_2\), and \(\mathsf{BL\text{-}kernel}\) may appear only after that continuation has localized the first remaining obstruction inside the readout/kernel stratum. In the deterministic syntax these two later actions are rewritten as \(\mathsf{continue}\) and \(\mathsf{kernel}\), with the same ordering restriction. Thus every admissible table output produces exactly one admissible decoder word, and every admissible decoder word dereferences to exactly one table-consistent reopening route.

Therefore the syntax introduces no new mathematical burden. It only compresses the already normalized routing data into a finite QAM-readable word system. This proves equivalence to the current decision table and faithfulness of the compression.
\end{proof}

\begin{remark}[Immediate machine reading of the deterministic decoder syntax]
\label{rem:immediate-machine-reading-deterministic-decoder-syntax}
The present line may now be summarized by the short machine-readable sentences
\[
\mathsf{classify}\to\mathsf{hold},\qquad
\mathsf{classify}\to\mathsf{ip},\qquad
\mathsf{classify}\to\mathsf{bsbl}\to\mathsf{continue}\to\mathsf{kernel},
\]
with the understanding that later symbols may only be appended when the earlier route has actually been entered and justified. In that sense the QAM layer now has not only tags and tables, but also a first deterministic decoder syntax that a future theorem-store, proof assistant, or simulation layer could read without reopening the entire surrounding prose.
\end{remark}



\subsection*{19AZ$^{(4e)}$. First finite transition function for deterministic residual-signature routing}
\label{sec:first-finite-transition-function-deterministic-residual-signature-routing}
The deterministic decoder syntax may now be compressed once more into a tiny state transition rule. The point is again modest: not to introduce a new autonomous automaton theory, but only to expose the present routing discipline as a finite partial transition function that the future machine layer can read directly. In particular, the current line should make explicit that only three first moves are admitted from the classification state, while deeper continuation steps remain locked behind the BS/BL side.

\begin{definition}[First finite transition function for the deterministic decoder]
\label{def:first-finite-transition-function-deterministic-decoder}
Let
\[
Q_{\mathrm{dec}}=\bigl\{\mathsf{classify},\,\mathsf{hold},\,\mathsf{ip},\,\mathsf{bsbl},\,\mathsf{continue},\,\mathsf{kernel}\bigr\}
\]
be the current decoder-state set, and let
\[
\Sigma_{\mathrm{dec}}=\bigl\{\mathsf{none},\,\mathsf{local},\,\mathsf{trans},\,\mathsf{stilllive},\,\mathsf{kernelstratum}\bigr\}
\]
be the present event alphabet, where:
\begin{itemize}[leftmargin=2em]
    \item \(\mathsf{none}\) means that no stable residual-signature recurrence class appears,
    \item \(\mathsf{local}\) means that the sampled class is purely local / imbalance-stable,
    \item \(\mathsf{trans}\) means that the sampled class is translation-sensitive,
    \item \(\mathsf{stilllive}\) means that after entering the BS/BL side the obstruction remains live beyond \(\mathrm{BL}_2\),
    \item \(\mathsf{kernelstratum}\) means that the first localized remaining obstruction lies in the readout/kernel stratum.
\end{itemize}
The \emph{first finite transition function for deterministic residual-signature routing} is the partial map
\[
\delta_{\mathrm{dec}}:Q_{\mathrm{dec}}\times \Sigma_{\mathrm{dec}}\dashrightarrow Q_{\mathrm{dec}}
\]
given by the only admissible transitions
\[
\begin{aligned}
\delta_{\mathrm{dec}}(\mathsf{classify},\mathsf{none})&=\mathsf{hold},\\
\delta_{\mathrm{dec}}(\mathsf{classify},\mathsf{local})&=\mathsf{ip},\\
\delta_{\mathrm{dec}}(\mathsf{classify},\mathsf{trans})&=\mathsf{bsbl},\\
\delta_{\mathrm{dec}}(\mathsf{bsbl},\mathsf{stilllive})&=\mathsf{continue},\\
\delta_{\mathrm{dec}}(\mathsf{continue},\mathsf{kernelstratum})&=\mathsf{kernel}.
\end{aligned}
\]
All other transitions are undefined at the present stage. The state outputs are dereferenced through the same normalized map \(\mathsf{Dec}_{\mathrm{QAM}}\) from Definition~\ref{def:first-deterministic-qam-decoder-syntax}.
\end{definition}

\begin{proposition}[The finite transition function is equivalent to the deterministic decoder syntax and remains minimal on the current live events]
\label{prop:finite-transition-function-equivalent-deterministic-decoder-syntax}
Under the current theorem/tool layer, the partial transition function from Definition~\ref{def:first-finite-transition-function-deterministic-decoder} carries exactly the same routing information as the deterministic decoder syntax from Definition~\ref{def:first-deterministic-qam-decoder-syntax}. More precisely:
\begin{enumerate}[label=(\roman*)]
    \item the three first moves from \(\mathsf{classify}\) are exactly \(\mathsf{hold}\), \(\mathsf{ip}\), and \(\mathsf{bsbl}\), corresponding respectively to the events \(\mathsf{none}\), \(\mathsf{local}\), and \(\mathsf{trans}\);
    \item the only admissible continuation on the bridge/localization side is \(\mathsf{bsbl}\xrightarrow{\mathsf{stilllive}}\mathsf{continue}\);
    \item the only admissible further refinement is \(\mathsf{continue}\xrightarrow{\mathsf{kernelstratum}}\mathsf{kernel}\);
    \item every undefined transition would either skip a smaller admissible reopening route or manufacture a carrier/event type not yet recorded in the present line.
\end{enumerate}
Hence \(\delta_{\mathrm{dec}}\) is a faithful finite-state compression of the current routing discipline and is minimal on the currently live diagnostic events.
\end{proposition}

\begin{proof}
The first three transitions are exactly the three initial output words from Proposition~\ref{prop:deterministic-decoder-syntax-equivalent-current-decision-table}: \(\mathsf{classify}\to\mathsf{hold}\), \(\mathsf{classify}\to\mathsf{ip}\), and \(\mathsf{classify}\to\mathsf{bsbl}\). The transition function merely rewrites these words as event-labeled state moves. The event names \(\mathsf{none}\), \(\mathsf{local}\), and \(\mathsf{trans}\) are exactly the presently live coarse classes isolated by the decision table and the two decoder propositions.

Once the route has entered \(\mathsf{bsbl}\), the deterministic syntax allows only the optional later symbols \(\mathsf{continue}\) and \(\mathsf{kernel}\). By Definition~\ref{def:first-deterministic-qam-decoder-syntax}, \(\mathsf{continue}\) may appear only after \(\mathsf{bsbl}\), and by Proposition~\ref{prop:first-explicit-decision-table-conservative-and-minimal} it is justified precisely when the obstruction remains live beyond \(\mathrm{BL}_2\). This is the transition \(\delta_{\mathrm{dec}}(\mathsf{bsbl},\mathsf{stilllive})=\mathsf{continue}\). Likewise, \(\mathsf{kernel}\) may occur only after that continuation has localized the first remaining obstruction inside the readout/kernel stratum, giving \(\delta_{\mathrm{dec}}(\mathsf{continue},\mathsf{kernelstratum})=\mathsf{kernel}\).

No further transitions are admissible on the current line. Any direct jump from \(\mathsf{classify}\) to \(\mathsf{continue}\) or \(\mathsf{kernel}\) would bypass the smaller bridge/localization entry route, while any jump from \(\mathsf{ip}\) to a later BS/BL state would conflate two presently separated reopening disciplines. And any additional event symbol would exceed the currently normalized recurrence classes. Therefore the partial transition function is equivalent to the deterministic decoder syntax, conservative with respect to the recorded carriers, and minimal on the currently live events.
\end{proof}

\begin{remark}[Immediate automaton reading of the current decoder]
\label{rem:immediate-automaton-reading-current-decoder}
The present line can now be read as a tiny partial automaton:
\[
\begin{aligned}
\mathsf{classify} &\xrightarrow{\mathsf{none}} \mathsf{hold},\\
\mathsf{classify} &\xrightarrow{\mathsf{local}} \mathsf{ip},\\
\mathsf{classify} &\xrightarrow{\mathsf{trans}} \mathsf{bsbl}
\xrightarrow{\mathsf{stilllive}} \mathsf{continue}
\xrightarrow{\mathsf{kernelstratum}} \mathsf{kernel}.
\end{aligned}
\]
Through the dereferencing map \(\mathsf{Dec}_{\mathrm{QAM}}\), this automaton is still nothing more than the already normalized theorem/tool layer. But it now has a form that a later theorem store, proof assistant, or Monte-Carlo controller could read as a deterministic transition discipline instead of as prose plus scattered tags.
\end{remark}


\subsection*{19AZ$^{(4f)}$. First observer/event normal form for deterministic residual-signature words}
\label{sec:first-observer-event-normal-form-deterministic-residual-signature-words}
The tiny decoder automaton from the previous subsection still presupposes that the incoming event has already been named as one of the symbols \(\mathsf{none}\), \(\mathsf{local}\), \(\mathsf{trans}\), \(\mathsf{stilllive}\), or \(\mathsf{kernelstratum}\). The next conservative step is therefore to normalize not the output route once more, but the observer-side input itself. The present line does not need a large language for that purpose. It only needs a short finite event-word normal form recording how much of the bridge/localization continuation has actually become visible at readout.

\begin{definition}[First observer/event normal form for deterministic residual-signature words]
\label{def:first-observer-event-normal-form-deterministic-residual-signature-words}
Let
\[
\mathcal W^{\mathrm{nf}}_{\mathrm{dec}}
=
\bigl\{
\mathsf{none},\,
\mathsf{local},\,
\mathsf{trans},\,
\mathsf{trans}\,\mathsf{stilllive},\,
\mathsf{trans}\,\mathsf{stilllive}\,\mathsf{kernelstratum}
\bigr\}
\]
be the \emph{first observer/event normal-form language} for deterministic residual-signature routing. The associated partial extended transition map is defined recursively by
\[
\delta^{\ast}_{\mathrm{dec}}(q,\varepsilon)=q,
\qquad
\delta^{\ast}_{\mathrm{dec}}(q,aw)=
\delta^{\ast}_{\mathrm{dec}}\bigl(\delta_{\mathrm{dec}}(q,a),w\bigr)
\]
whenever the right-hand side is defined, where \(a\in\Sigma_{\mathrm{dec}}\) and \(w\) is a finite word in the same alphabet.

The \emph{observer/event normal form map}
\[
\mathsf{NF}_{\mathrm{obs}}: \text{currently normalized residual-signature cases}
\longrightarrow \mathcal W^{\mathrm{nf}}_{\mathrm{dec}}
\]
assigns the unique normal-form word as follows:
\begin{enumerate}[label=(\roman*)]
    \item no stable recurrence class is observed \(\mapsto \mathsf{none}\),
    \item a purely local / imbalance-stable class is observed \(\mapsto \mathsf{local}\),
    \item a translation-sensitive class is observed but no remaining obstruction survives beyond \(\mathrm{BL}_2\) \(\mapsto \mathsf{trans}\),
    \item a translation-sensitive class is observed and the obstruction remains live beyond \(\mathrm{BL}_2\) \(\mapsto \mathsf{trans}\,\mathsf{stilllive}\),
    \item the previous case holds and the first localized remaining obstruction lies in the readout/kernel stratum \(\mapsto \mathsf{trans}\,\mathsf{stilllive}\,\mathsf{kernelstratum}\).
\end{enumerate}
\end{definition}

\begin{proposition}[The observer/event normal forms are exactly the currently admissible decoder words and reproduce the normalized reopening routes]
\label{prop:observer-event-normal-forms-exactly-currently-admissible-decoder-words}
Under the current theorem/tool layer, the language \(\mathcal W^{\mathrm{nf}}_{\mathrm{dec}}\) from Definition~\ref{def:first-observer-event-normal-form-deterministic-residual-signature-words} is exactly the set of currently admissible observer-side event words for the deterministic decoder. More precisely:
\begin{enumerate}[label=(\roman*)]
    \item every word in \(\mathcal W^{\mathrm{nf}}_{\mathrm{dec}}\) is admitted by the partial automaton \(\delta_{\mathrm{dec}}\) when started at \(\mathsf{classify}\),
    \item the terminal states are respectively
    \[
    \delta^{\ast}_{\mathrm{dec}}(\mathsf{classify},\mathsf{none})=\mathsf{hold},\qquad
    \delta^{\ast}_{\mathrm{dec}}(\mathsf{classify},\mathsf{local})=\mathsf{ip},
    \]
    \[
    \delta^{\ast}_{\mathrm{dec}}(\mathsf{classify},\mathsf{trans})=\mathsf{bsbl},\qquad
    \delta^{\ast}_{\mathrm{dec}}(\mathsf{classify},\mathsf{trans}\,\mathsf{stilllive})=\mathsf{continue},
    \]
    \[
    \delta^{\ast}_{\mathrm{dec}}(\mathsf{classify},\mathsf{trans}\,\mathsf{stilllive}\,\mathsf{kernelstratum})=\mathsf{kernel},
    \]
    and these states dereference exactly to the normalized routes already fixed in the QAM tool layer;
    \item no other finite event word is currently admissible without either violating the transition discipline of \(\delta_{\mathrm{dec}}\) or introducing an observer-side distinction not yet normalized in the present line.
\end{enumerate}
Hence the current Monte-Carlo/QAM interface may now be read not only as a state machine, but as a finite normal-form language of observer-side diagnostic words.
\end{proposition}

\begin{proof}
The five displayed words are admitted directly by the transition function from Definition~\ref{def:first-finite-transition-function-deterministic-decoder}. The one-letter words \(\mathsf{none}\), \(\mathsf{local}\), and \(\mathsf{trans}\) are exactly the three first-step events recorded there. Once the route has entered \(\mathsf{bsbl}\), the only currently permitted continuation event is \(\mathsf{stilllive}\), and once the route has entered \(\mathsf{continue}\), the only currently permitted further refinement is \(\mathsf{kernelstratum}\). Therefore the longer words \(\mathsf{trans}\,\mathsf{stilllive}\) and \(\mathsf{trans}\,\mathsf{stilllive}\,\mathsf{kernelstratum}\) are admitted as well.

The recursive definition of \(\delta^{\ast}_{\mathrm{dec}}\) then gives the terminal states by straightforward evaluation. For \(\mathsf{none}\) and \(\mathsf{local}\), only one transition is applied, yielding \(\mathsf{hold}\) and \(\mathsf{ip}\). For \(\mathsf{trans}\), the first step yields \(\mathsf{bsbl}\). Appending \(\mathsf{stilllive}\) moves to \(\mathsf{continue}\), and appending \(\mathsf{kernelstratum}\) moves from there to \(\mathsf{kernel}\). By Definition~\ref{def:first-deterministic-qam-decoder-syntax}, each terminal state dereferences to the already normalized reopening route: \(\mathsf{hold}\) means reopen nothing, \(\mathsf{ip}\) means the chain \(\mathrm{TR}_4\Rightarrow \mathrm{IP}_{10}\Rightarrow \mathrm{IP}_9\Rightarrow \mathrm{IP}_8\), \(\mathsf{bsbl}\) means the chain \(\mathrm{TR}_4\Rightarrow \mathrm{BS}_1\Rightarrow \mathrm{BS}_3\Rightarrow \mathrm{BL}_2\), \(\mathsf{continue}\) appends \(\mathrm{BL}_4\), and \(\mathsf{kernel}\) appends \(\mathrm{BL}_5\).

It remains to exclude any additional finite event word. A word beginning with anything other than \(\mathsf{none}\), \(\mathsf{local}\), or \(\mathsf{trans}\) is impossible because \(\delta_{\mathrm{dec}}\) is undefined at \(\mathsf{classify}\) on all other events. A word beginning with \(\mathsf{none}\) or \(\mathsf{local}\) cannot be extended further without leaving the currently normalized routing discipline. And a word beginning with \(\mathsf{trans}\) may only continue with \(\mathsf{stilllive}\), then optionally with \(\mathsf{kernelstratum}\), because no further transitions are currently recorded. Thus every admissible observer-side case has exactly one normal-form word and no other finite word is currently allowed.
\end{proof}

\begin{remark}[Immediate machine reading of the observer/event normal form]
\label{rem:immediate-machine-reading-observer-event-normal-form}
The current live interface may now be summarized by the finite observer-language
\[
\mathsf{none},\qquad
\mathsf{local},\qquad
\mathsf{trans},\qquad
\mathsf{trans}\,\mathsf{stilllive},\qquad
\mathsf{trans}\,\mathsf{stilllive}\,\mathsf{kernelstratum}.
\]
This is the first point at which the tool layer, the QAM tags, the Monte-Carlo recurrence classes, the deterministic decoder syntax, and the partial transition automaton all meet in one compact object: a finite word language on the observer side whose outputs are already normalized reopening routes on the theorem side.
\end{remark}



\subsection*{19AZ$^{(4g)}$. First observer-word evaluation map into deterministic decoder outputs}
\label{sec:first-observer-word-evaluation-map-deterministic-decoder-outputs}
The observer/event normal forms now provide a finite input language for the deterministic decoder. The next conservative compression step is to package the whole passage from normalized observer word to normalized theorem-side reopening output into one explicit evaluation map. This adds no new event class and no new route. It merely composes the already stabilized normal-form language, transition automaton, and QAM dereferencing rule into one finite readout operator.

\begin{definition}[First observer-word evaluation map into deterministic decoder outputs]
\label{def:first-observer-word-evaluation-map-deterministic-decoder-outputs}
Let
\[
\mathcal W^{\mathrm{nf}}_{\mathrm{dec}}
=
\bigl\{
\mathsf{none},\,
\mathsf{local},\,
\mathsf{trans},\,
\mathsf{trans}\,\mathsf{stilllive},\,
\mathsf{trans}\,\mathsf{stilllive}\,\mathsf{kernelstratum}
\bigr\}
\]
be the observer/event normal-form language from Definition~\ref{def:first-observer-event-normal-form-deterministic-residual-signature-words}. Define the \emph{observer-word evaluation map}
\[
\mathsf{Eval}_{\mathrm{obs}}:
\mathcal W^{\mathrm{nf}}_{\mathrm{dec}}
\longrightarrow
Q_{\mathrm{dec}}\times \mathcal R_{\mathrm{tool}}^{\mathrm{out}}
\]
by
\[
\mathsf{Eval}_{\mathrm{obs}}(w)
=
\Bigl(
\delta^{\ast}_{\mathrm{dec}}(\mathsf{classify},w),
\mathsf{Dec}_{\mathrm{QAM}}\bigl(\delta^{\ast}_{\mathrm{dec}}(\mathsf{classify},w)\bigr)
\Bigr),
\]
where \(Q_{\mathrm{dec}}=\{\mathsf{hold},\mathsf{ip},\mathsf{bsbl},\mathsf{continue},\mathsf{kernel}\}\) is the current decoder-state set and \(\mathcal R_{\mathrm{tool}}^{\mathrm{out}}\) denotes the finite set of currently normalized reopening outputs:
\[
\mathcal R_{\mathrm{tool}}^{\mathrm{out}}
=
\Bigl\{
\begin{aligned}
&\varnothing,\\
&[\mathrm{TR}_4,\mathrm{IP}_{10},\mathrm{IP}_9,\mathrm{IP}_8],\\
&[\mathrm{TR}_4,\mathrm{BS}_1,\mathrm{BS}_3,\mathrm{BL}_2],\\
&[\mathrm{TR}_4,\mathrm{BS}_1,\mathrm{BS}_3,\mathrm{BL}_2,\mathrm{BL}_4],\\
&[\mathrm{TR}_4,\mathrm{BS}_1,\mathrm{BS}_3,\mathrm{BL}_2,\mathrm{BL}_4,\mathrm{BL}_5]
\end{aligned}
\Bigr\}.
\]
Equivalently,
\begin{align*}
\mathsf{Eval}_{\mathrm{obs}}(\mathsf{none})
&=(\mathsf{hold},\varnothing),\\
\mathsf{Eval}_{\mathrm{obs}}(\mathsf{local})
&=(\mathsf{ip},[\mathrm{TR}_4,\mathrm{IP}_{10},\mathrm{IP}_9,\mathrm{IP}_8]),\\
\mathsf{Eval}_{\mathrm{obs}}(\mathsf{trans})
&=(\mathsf{bsbl},[\mathrm{TR}_4,\mathrm{BS}_1,\mathrm{BS}_3,\mathrm{BL}_2]),\\
\mathsf{Eval}_{\mathrm{obs}}(\mathsf{trans}\,\mathsf{stilllive})
&=(\mathsf{continue},[\mathrm{TR}_4,\mathrm{BS}_1,\mathrm{BS}_3,\mathrm{BL}_2,\mathrm{BL}_4]),\\
\mathsf{Eval}_{\mathrm{obs}}(\mathsf{trans}\,\mathsf{stilllive}\,\mathsf{kernelstratum})
&=(\mathsf{kernel},[\mathrm{TR}_4,\mathrm{BS}_1,\mathrm{BS}_3,\mathrm{BL}_2,\mathrm{BL}_4,\mathrm{BL}_5]).
\end{align*}
\end{definition}

\begin{proposition}[The observer-word evaluation map is total on the current normal-form language and reproduces the unique normalized reopening outputs]
\label{prop:observer-word-evaluation-map-total-current-normal-form-language}
Under the current theorem/tool layer, the map \(\mathsf{Eval}_{\mathrm{obs}}\) from Definition~\ref{def:first-observer-word-evaluation-map-deterministic-decoder-outputs} is total on \(\mathcal W^{\mathrm{nf}}_{\mathrm{dec}}\). More precisely:
\begin{enumerate}[label=(\roman*)]
    \item every word in \(\mathcal W^{\mathrm{nf}}_{\mathrm{dec}}\) has exactly one decoder terminal state and exactly one currently normalized reopening output under \(\mathsf{Eval}_{\mathrm{obs}}\),
    \item the first component of \(\mathsf{Eval}_{\mathrm{obs}}(w)\) is the unique terminal state reached by the partial automaton \(\delta^{\ast}_{\mathrm{dec}}\),
    \item the second component is exactly the smallest presently admissible theorem-side reopening route compatible with that terminal state,
    \item no currently normalized observer word evaluates to two different reopening outputs.
\end{enumerate}
Hence the present live interface may now be read as one deterministic finite composite map
\[
\mathsf{NF}_{\mathrm{obs}}
\longrightarrow
\mathsf{Eval}_{\mathrm{obs}}
\longrightarrow
\text{normalized reopening output},
\]
rather than as a chain of separately cited decoder pieces.
\end{proposition}

\begin{proof}
By Proposition~\ref{prop:observer-event-normal-forms-exactly-currently-admissible-decoder-words}, every word in \(\mathcal W^{\mathrm{nf}}_{\mathrm{dec}}\) is admitted by the extended transition map \(\delta^{\ast}_{\mathrm{dec}}\) when started at \(\mathsf{classify}\), and its terminal state is uniquely determined. Thus the first component of \(\mathsf{Eval}_{\mathrm{obs}}\) is total and single-valued on the whole normal-form language.

The second component is obtained by applying the deterministic QAM dereferencing map \(\mathsf{Dec}_{\mathrm{QAM}}\) to that unique terminal state. By Definition~\ref{def:first-deterministic-qam-decoder-syntax}, each currently live terminal state carries exactly one normalized reopening output: \(\mathsf{hold}\) gives \(\varnothing\), \(\mathsf{ip}\) gives the fixed IP route \([\mathrm{TR}_4,\mathrm{IP}_{10},\mathrm{IP}_9,\mathrm{IP}_8]\), \(\mathsf{bsbl}\) gives the fixed bridge/localization entry route \([\mathrm{TR}_4,\mathrm{BS}_1,\mathrm{BS}_3,\mathrm{BL}_2]\), \(\mathsf{continue}\) appends \(\mathrm{BL}_4\), and \(\mathsf{kernel}\) appends \(\mathrm{BL}_5\). Therefore the second component is also total and single-valued.

Minimality follows from the construction of the current decoder. The state \(\mathsf{hold}\) reopens nothing, \(\mathsf{ip}\) is already the smallest admissible local route, and the bridge/localization branch enters first through \(\mathsf{bsbl}\) before optional refinement to \(\mathsf{continue}\) and \(\mathsf{kernel}\). Any alternative output would either reopen more than the current state requires or bypass an earlier normalized stage. Hence no current observer word can evaluate to two different normalized reopening outputs.
\end{proof}

\begin{remark}[Immediate operational reading of \texorpdfstring{$\mathsf{Eval}_{\mathrm{obs}}$}{Eval_obs}]
\label{rem:immediate-operational-reading-eval-obs}
The current finite live interface may now be cited in one line:
\[
\begin{aligned}
\text{observer-side recurrence class}
&\xrightarrow{\ \mathsf{NF}_{\mathrm{obs}}\ }
\text{normal-form word}\\
&\xrightarrow{\ \mathsf{Eval}_{\mathrm{obs}}\ }
(\text{decoder state},\text{normalized reopening output}).
\end{aligned}
\]
At this point, the Monte-Carlo diagnostics, the QAM tag layer, the finite decoder automaton, and the theorem-side reopening discipline have become one explicit finite evaluation pipeline. This is the first place in the present line where a later controller could read the whole diagnostic-to-route passage through a single total map on a finite language.
\end{remark}



\subsection*{19AZ$^{(4h)}$. First compact output-normalform tokens for deterministic decoder results}
\label{sec:first-compact-output-normalform-tokens-deterministic-decoder-results}
The observer-word evaluation map now packages every currently admissible observer-side normal form into one of five deterministic theorem-side outputs. The next conservative compression step is to replace those explicit route lists by a finite family of compact output tokens. This adds no new state, no new route, and no new observer word. It merely supplies a shorter QAM-side output carrier for the five already stabilized decoder results.

\begin{definition}[First compact output-normalform tokens for deterministic decoder results]
\label{def:first-compact-output-normalform-tokens-deterministic-decoder-results}
Define the finite set of \emph{output-normalform tokens}
\[
\mathcal T^{\mathrm{nf}}_{\mathrm{out}}
=
\{\mathsf{OH},\mathsf{OI},\mathsf{OB},\mathsf{OC},\mathsf{OK}\},
\]
where the tokens are read as
\begin{align*}
\mathsf{OH}&\leftrightarrow (\mathsf{hold},\varnothing),\\
\mathsf{OI}&\leftrightarrow (\mathsf{ip},[\mathrm{TR}_4,\mathrm{IP}_{10},\mathrm{IP}_9,\mathrm{IP}_8]),\\
\mathsf{OB}&\leftrightarrow (\mathsf{bsbl},[\mathrm{TR}_4,\mathrm{BS}_1,\mathrm{BS}_3,\mathrm{BL}_2]),\\
\mathsf{OC}&\leftrightarrow (\mathsf{continue},[\mathrm{TR}_4,\mathrm{BS}_1,\mathrm{BS}_3,\mathrm{BL}_2,\mathrm{BL}_4]),\\
\mathsf{OK}&\leftrightarrow (\mathsf{kernel},[\mathrm{TR}_4,\mathrm{BS}_1,\mathrm{BS}_3,\mathrm{BL}_2,\mathrm{BL}_4,\mathrm{BL}_5]).
\end{align*}
Equivalently, define the \emph{output-tokenization map}
\[
\mathsf{Tok}_{\mathrm{out}}:
Q_{\mathrm{dec}}\times \mathcal R_{\mathrm{tool}}^{\mathrm{out}}
\longrightarrow
\mathcal T^{\mathrm{nf}}_{\mathrm{out}}
\]
by sending the five currently normalized evaluation outputs from Definition~\ref{def:first-observer-word-evaluation-map-deterministic-decoder-outputs} to the corresponding tokens above. Composing with \(\mathsf{Eval}_{\mathrm{obs}}\) gives the compact output map
\[
\mathsf{Out}^{\mathrm{tok}}_{\mathrm{obs}}
:=
\mathsf{Tok}_{\mathrm{out}}\circ \mathsf{Eval}_{\mathrm{obs}}:
\mathcal W^{\mathrm{nf}}_{\mathrm{dec}}
\longrightarrow
\mathcal T^{\mathrm{nf}}_{\mathrm{out}}.
\]
Thus
\begin{align*}
\mathsf{Out}^{\mathrm{tok}}_{\mathrm{obs}}(\mathsf{none})&=\mathsf{OH},\\
\mathsf{Out}^{\mathrm{tok}}_{\mathrm{obs}}(\mathsf{local})&=\mathsf{OI},\\
\mathsf{Out}^{\mathrm{tok}}_{\mathrm{obs}}(\mathsf{trans})&=\mathsf{OB},\\
\mathsf{Out}^{\mathrm{tok}}_{\mathrm{obs}}(\mathsf{trans}\,\mathsf{stilllive})&=\mathsf{OC},\\
\mathsf{Out}^{\mathrm{tok}}_{\mathrm{obs}}(\mathsf{trans}\,\mathsf{stilllive}\,\mathsf{kernelstratum})&=\mathsf{OK}.
\end{align*}
\end{definition}

\begin{proposition}[The output-normalform tokens are faithful and complete for the current deterministic decoder results]
\label{prop:output-normalform-tokens-faithful-complete-current-deterministic-decoder-results}
Under the current live theorem/tool layer, the token map \(\mathsf{Tok}_{\mathrm{out}}\) is faithful and complete on the presently normalized decoder outputs. More precisely:
\begin{enumerate}[label=(\roman*)]
    \item every value of \(\mathsf{Eval}_{\mathrm{obs}}\) is sent to exactly one token in \(\mathcal T^{\mathrm{nf}}_{\mathrm{out}}\),
    \item different current decoder/output pairs receive different tokens,
    \item the token \(\mathsf{Out}^{\mathrm{tok}}_{\mathrm{obs}}(w)\) determines both the terminal decoder state and the unique normalized reopening route attached to that state,
    \item no information from the current finite live interface is lost when passing from explicit output pairs to output-normalform tokens.
\end{enumerate}
Hence the present diagnostic-to-route pipeline may now also be read in the compact form
\[
\text{observer class}
\longrightarrow
\text{normal-form word}
\longrightarrow
\text{output-normalform token}.
\]
\end{proposition}

\begin{proof}
By Proposition~\ref{prop:observer-word-evaluation-map-total-current-normal-form-language}, the map \(\mathsf{Eval}_{\mathrm{obs}}\) takes values in exactly five currently normalized decoder/output pairs. Definition~\ref{def:first-compact-output-normalform-tokens-deterministic-decoder-results} assigns one token to each of those five pairs and to no others, so every current evaluation value receives exactly one token.

Faithfulness is immediate from the construction: \(\mathsf{OH},\mathsf{OI},\mathsf{OB},\mathsf{OC},\mathsf{OK}\) are pairwise distinct symbols, each tied to one different decoder/output pair. Thus two different current evaluation outputs cannot collapse to the same token. Conversely, each token has exactly one prescribed decoder state and exactly one prescribed normalized reopening route, so the token determines the full current output pair.

Therefore the passage from explicit pairs to compact tokens is lossless on the present finite interface. It is merely a normalization of notation, not a quotient by any nontrivial equivalence. The compact map \(\mathsf{Out}^{\mathrm{tok}}_{\mathrm{obs}}\) preserves exactly the same current routing content as \(\mathsf{Eval}_{\mathrm{obs}}\), only in shorter QAM-readable form.
\end{proof}

\begin{remark}[Why the compact output tokens matter]
\label{rem:why-compact-output-tokens-matter}
The live interface can now be written in its shortest current form:
\[
\text{observer-side recurrence class}
\xrightarrow{\ \mathsf{NF}_{\mathrm{obs}}\ }
\text{normal-form word}
\xrightarrow{\ \mathsf{Out}^{\mathrm{tok}}_{\mathrm{obs}}\ }
\{\mathsf{OH},\mathsf{OI},\mathsf{OB},\mathsf{OC},\mathsf{OK}\}.
\]
At this point the theorem-side reopening discipline already has a compact output alphabet. A later controller or QAM overlay no longer needs to carry the full explicit route list every time; it may cite the token first and dereference the corresponding route only when the detailed theorem-side reopening path is actually needed.
\end{remark}


\subsection*{19AZ$^{(5a)}$. First signature-injectivity normal form for the live OP~1 kernel}
\label{sec:first-signature-injectivity-normal-form-live-op1-kernel}
The normalized routing layer has now done its job on the OP~1 side: it has compressed the live reopening path to
\[
\mathrm{TR}_4 \Longrightarrow \mathrm{IP}_{10} \Longrightarrow \mathrm{IP}_9 \Longrightarrow \mathrm{IP}_8.
\]
The next honest mathematical move is therefore not another routing refinement but a sharper obstruction normal form for the first live carrier in that chain. The remaining burden should now be stated as a local injectivity problem for the retained pinning signature itself: can two genuinely distinct admissible primitive photon-first-zero loops occupy the same retained post-frame corridor class while preserving the same total-channel profile, the same retained optimum location, and the same local comparison/defect readout?

\begin{quote}\small
\noindent\textbf{Reader cue.} Read the next block as an obstruction normal form for the live OP~1 kernel. Its purpose is to rewrite the remaining burden as exclusion of a second admissible primitive loop with the same retained pinning signature, not to claim that this exclusion has already been proved.
\end{quote}

\begin{definition}[Retained pinning signature on the OP~1 corridor]
\label{def:retained-pinning-signature-op1-corridor}
Assume the same-corridor reversal certificate of Definition~\ref{def:op1-same-corridor-reversal}. For every admissible primitive photon-first-zero closed-loop realization $L$ on the retained post-frame corridor, define the \emph{retained pinning signature}
\[
\mathsf{Sig}_{\mathrm{pin}}(L)
\]
to be the complete retained comparison packet carried by the present line on the intended admissible windows, namely: the retained total-channel profile, the retained optimum location, and the standing local comparison/defect data attached to $L$ in the compressed machine-shadow reading. Thus two loops have the same retained pinning signature exactly when they are indistinguishable to the current OP~1 comparison line at the level that feeds Definition~\ref{def:op1-optimum-pinning-comparison-rigidity}.
\end{definition}

\begin{definition}[Duplicate-signature witness for the live OP~1 kernel]
\label{def:duplicate-signature-witness-live-op1-kernel}
Assume Definition~\ref{def:retained-pinning-signature-op1-corridor}. A \emph{duplicate-signature witness} for the live OP~1 kernel is a pair $(L,L')$ of admissible primitive photon-first-zero closed-loop realizations on the retained post-frame corridor such that:
\begin{enumerate}[label=(\roman*)]
    \item $L$ and $L'$ are genuinely distinct up to the standing admissible corridor equivalence,
    \item both preserve the same retained total-channel profile on the intended admissible windows,
    \item \(\mathsf{Sig}_{\mathrm{pin}}(L)=\mathsf{Sig}_{\mathrm{pin}}(L')\).
\end{enumerate}
Equivalently, a duplicate-signature witness is a second primitive loop placement that the current comparison line cannot distinguish from the intended retained loop even though it remains genuinely distinct inside the retained corridor class.
\end{definition}

\begin{proposition}[The live OP~1 burden is equivalent to exclusion of duplicate-signature witnesses]
\label{prop:live-op1-burden-equivalent-exclusion-duplicate-signature-witnesses}
Assume the same-corridor reversal certificate of Definition~\ref{def:op1-same-corridor-reversal}. Then the following are equivalent on the intended admissible windows:
\begin{enumerate}[label=(\roman*)]
    \item the optimum-pinning comparison-rigidity certificate of Definition~\ref{def:op1-optimum-pinning-comparison-rigidity} holds;
    \item no duplicate-signature witness exists in the sense of Definition~\ref{def:duplicate-signature-witness-live-op1-kernel}.
\end{enumerate}
Consequently, together with Proposition~\ref{prop:op1-rigidity-gap-uniqueness-bridge} and the standing no-second-carrier / strict-upper-gap exclusions already built into the current line, exclusion of duplicate-signature witnesses implies
\[
\Delta_{\mathrm{CL}}=0.
\]
Equivalently, the present residual OP~1 burden may now be read exactly as the injectivity of the retained pinning signature on the admissible primitive loop class carried by the retained post-frame corridor.
\end{proposition}

\begin{proof}
Definition~\ref{def:op1-optimum-pinning-comparison-rigidity} says that no genuinely distinct admissible primitive loop on the retained post-frame corridor can preserve the retained total-channel profile, the retained optimum location, and the standing local comparison/defect data while remaining distinct up to the admissible corridor equivalence. But this is exactly the negation of a duplicate-signature witness from Definition~\ref{def:duplicate-signature-witness-live-op1-kernel}, since the retained pinning signature packages precisely those same preserved data. Thus \emph{(i)} and \emph{(ii)} are equivalent.

Once this injectivity statement is available, Proposition~\ref{prop:op1-optimum-pinning-implies-rigidity} yields the retained-window rigidity certificate. Proposition~\ref{prop:op1-rigidity-gap-uniqueness-bridge} then upgrades that rigidity, under the standing no-second-carrier and strict-upper-gap exclusions, to admissible-window uniqueness and hence to \(\Delta_{\mathrm{CL}}=0\).
\end{proof}

\begin{remark}[What the new normal form buys for the next OP~1 step]
\label{rem:what-new-normal-form-buys-next-op1-step}
This is a real sharpening even though it is not yet a closure theorem. Before the present step, the residual burden was phrased as ``prove optimum-pinning comparison rigidity.'' After the present step, the live target is more concrete: show that the retained comparison packet
\[
\mathsf{Sig}_{\mathrm{pin}}
\]
is injective on admissible primitive first-zero loops in the retained post-frame corridor class. In practical terms, the next OP~1 attack should therefore look for a contradiction from a hypothetical duplicate-signature witness, not for a new free continuum-limit parameter. In the normalized QAM reading, this says that the first live carrier \(\mathrm{IP}_{10}\) should now be treated as a signature-injectivity problem for the retained pinning packet.
\end{remark}




\subsection*{19AZ$^{(5b)}$. First contradiction channel split for a duplicate-signature witness}
\label{sec:first-contradiction-channel-split-duplicate-signature-witness}
The normal form of 19AZ$^{(5a)}$ says that the live OP~1 burden is equivalent to the exclusion of duplicate-signature witnesses. The next honest move is therefore to ask what such a witness would have to do inside the present line. Because the retained pinning signature already fixes the retained total-channel profile, the retained optimum location, and the local comparison/defect packet, any residual distinction between two loops carrying the same signature has only two places left to live: either as a second primitive occupancy of the same retained corridor slot, or as a second admissible completion above that same slot. This is exactly the old no-second-carrier / strict-upper-gap dichotomy, but now pulled back to the new signature normal form.

\begin{quote}\small
\noindent\textbf{Reader cue.} Read the next block as the first contradiction routing split for the live OP~1 kernel. Its purpose is to show that a hypothetical duplicate-signature witness can only survive by feeding one of the two exclusion gates already built into the current line: duplicate primitive occupancy or duplicate admissible completion.
\end{quote}

\begin{definition}[Occupancy-side and completion-side duplicate-signature channels]
\label{def:occupancy-completion-channels-duplicate-signature-op1}
Assume the same-corridor reversal certificate of Definition~\ref{def:op1-same-corridor-reversal} and let $(L,L')$ be a duplicate-signature witness in the sense of Definition~\ref{def:duplicate-signature-witness-live-op1-kernel}. We say that the witness is of \emph{occupancy-side type} if the residual distinction between $L$ and $L'$ already appears as a second primitive occupancy of the retained post-frame corridor while the retained total-channel profile is kept fixed. We say that the witness is of \emph{completion-side type} if, after the common retained corridor slot and retained comparison packet are fixed, the residual distinction can only survive as a second admissible completion route above that same retained corridor while preserving the same retained status surface.
\end{definition}

\begin{proposition}[Every duplicate-signature witness feeds a standing OP~1 exclusion gate]
\label{prop:duplicate-signature-witness-feeds-standing-op1-exclusion-gate}
Assume the same-corridor reversal certificate of Definition~\ref{def:op1-same-corridor-reversal}. Let $(L,L')$ be a duplicate-signature witness in the sense of Definition~\ref{def:duplicate-signature-witness-live-op1-kernel}. Then at least one of the following holds on the intended admissible windows:
\begin{enumerate}[label=(\roman*)]
    \item $(L,L')$ is of occupancy-side type in the sense of Definition~\ref{def:occupancy-completion-channels-duplicate-signature-op1};
    \item $(L,L')$ is of completion-side type in the sense of Definition~\ref{def:occupancy-completion-channels-duplicate-signature-op1}.
\end{enumerate}
Consequently, every hypothetical failure of signature-injectivity for the live OP~1 kernel already routes into one of the two standing exclusion channels behind Proposition~\ref{prop:op1-rigidity-gap-uniqueness-bridge}: no-second-carrier on the retained corridor, or strict-upper-gap on admissible completion above that corridor.
\end{proposition}

\begin{proof}
By Definition~\ref{def:duplicate-signature-witness-live-op1-kernel}, the loops $L$ and $L'$ are genuinely distinct up to the standing admissible corridor equivalence, yet they preserve the same retained pinning signature. In particular, they preserve the same retained total-channel profile, the same retained optimum location, and the same local comparison/defect data.

Therefore the current line has already fixed everything that belongs to the retained comparison packet. If the residual distinction between $L$ and $L'$ is visible already at the level of primitive occupancy of the retained post-frame corridor, then the witness is of occupancy-side type by definition. Otherwise the common retained corridor slot and retained comparison packet do not separate the loops, so the remaining distinction can only survive as a second admissible completion above that same retained corridor while preserving the same retained status surface; this is exactly the completion-side case.

Hence every duplicate-signature witness must feed at least one of the two channels. The final sentence is then only a translation into the standing OP~1 language: occupancy-side witnesses are precisely the kind of residual duplication targeted by the no-second-carrier discipline, while completion-side witnesses are precisely the kind of residual duplication targeted by the strict-upper-gap exclusion in Proposition~\ref{prop:op1-rigidity-gap-uniqueness-bridge}.
\end{proof}

\begin{remark}[What this contradiction split buys for the next OP~1 attack]
\label{rem:what-contradiction-split-buys-next-op1-attack}
This still does not close OP~1, but it tells us where a closure attempt must now bite. A hypothetical duplicate-signature witness no longer floats as an abstract rival loop. It is forced into one of two already named contradiction channels. In practical terms, the next OP~1 step should try to show that the retained comparison packet cannot support either occupancy-side duplication on the retained corridor or completion-side duplication above that corridor. In the normalized QAM reading, this means that the live carrier \(\mathrm{IP}_{10}\) has now been resolved into a two-channel contradiction search whose outputs are already aligned with the standing no-second-carrier / strict-upper-gap gates feeding \(\mathrm{IP}_9\).
\end{remark}



\subsection*{19AZ$^{(5c)}$. Occupancy-side duplicate-signature witnesses collapse under the retained no-second-carrier discipline}
\label{sec:occupancy-side-duplicate-signature-witnesses-collapse-no-second-carrier}
After the channel split of 19AZ$^{(5b)}$, the cleanest next question is whether both contradiction channels are still genuinely live. On the present line, the occupancy-side channel sits closest to the already retained theorematic infrastructure: if a duplicate-signature witness already survives as a second primitive occupancy of the same retained corridor while the retained total-channel profile stays fixed, then it is no longer merely a new OP~1 curiosity. It is exactly the sort of duplication that the retained no-second-carrier discipline was built to forbid on the preserved total-channel side.

\begin{quote}\small
\noindent\textbf{Reader cue.} Read the next block as the first actual collapse theorem inside the new OP~1 contradiction split. Its purpose is to show that occupancy-side duplicate-signature witnesses are already excluded by the standing retained no-second-carrier discipline, so that the live burden cannot honestly remain two-channel after this point.
\end{quote}

\begin{proposition}[Occupancy-side duplicate-signature witnesses are excluded on the retained OP~1 corridor]
\label{prop:occupancy-side-duplicate-signature-witnesses-excluded-retained-op1-corridor}
Assume the same-corridor reversal certificate of Definition~\ref{def:op1-same-corridor-reversal}. Assume moreover that the retained no-second-carrier discipline continues to hold on the intended admissible windows at the retained total-channel level, in the compressed sense already available through Theorem~\ref{thm:first-no-second-carrier-theorem-on-the-terminal-interface}. Then no duplicate-signature witness of occupancy-side type exists in the sense of Definitions~\ref{def:duplicate-signature-witness-live-op1-kernel} and~\ref{def:occupancy-completion-channels-duplicate-signature-op1}.
\end{proposition}

\begin{proof}
Suppose for contradiction that \((L,L')\) is a duplicate-signature witness of occupancy-side type. By Definition~\ref{def:duplicate-signature-witness-live-op1-kernel}, the loops \(L\) and \(L'\) are genuinely distinct up to the standing admissible corridor equivalence, yet they preserve the same retained pinning signature. In particular, they preserve the same retained total-channel profile on the intended admissible windows.

By Definition~\ref{def:occupancy-completion-channels-duplicate-signature-op1}, the residual distinction already survives as a second primitive occupancy of the retained post-frame corridor while that retained total-channel profile is kept fixed. But this is precisely the kind of extra primitive retained-carrier occupancy excluded by the standing no-second-carrier discipline at retained total-channel level. Hence the assumed witness contradicts the retained no-second-carrier theorematic layer. Therefore no occupancy-side duplicate-signature witness can exist.
\end{proof}

\begin{corollary}[The live OP~1 duplicate-signature burden reduces to the completion-side channel]
\label{cor:live-op1-duplicate-signature-burden-reduces-completion-side-channel}
Under the hypotheses of Proposition~\ref{prop:occupancy-side-duplicate-signature-witnesses-excluded-retained-op1-corridor}, the following are equivalent on the intended admissible windows:
\begin{enumerate}[label=(\roman*)]
    \item a duplicate-signature witness for the live OP~1 kernel exists;
    \item a duplicate-signature witness of completion-side type exists.
\end{enumerate}
Consequently, after the present step the live OP~1 kernel is no longer honestly read as a two-channel contradiction search. The occupancy-side branch is already discharged by the retained no-second-carrier discipline, so the surviving burden is the exclusion of a second admissible completion above the retained corridor that preserves the same retained pinning signature.
\end{corollary}

\begin{proof}
If a duplicate-signature witness exists, Proposition~\ref{prop:duplicate-signature-witness-feeds-standing-op1-exclusion-gate} shows that it must be either of occupancy-side type or of completion-side type. Proposition~\ref{prop:occupancy-side-duplicate-signature-witnesses-excluded-retained-op1-corridor} eliminates the occupancy-side case. Hence every surviving duplicate-signature witness must be completion-side. The converse is immediate, since every completion-side witness is in particular a duplicate-signature witness.
\end{proof}

\begin{remark}[What this first collapse buys for the next OP~1 step]
\label{rem:what-first-collapse-buys-next-op1-step}
This is the first point in the new post-release OP~1 line where one of the two contradiction channels is not merely named but actually discharged. The live burden is therefore smaller than it looked in 19AZ$^{(5b)}$: the retained comparison packet is already rigid enough against duplicate \emph{occupancy} once the inherited no-second-carrier discipline is invoked. What remains is subtler and now sharply localized. One must show that the same retained pinning signature cannot support a second \emph{admissible completion} above the already fixed retained corridor. In the normalized QAM reading, the live carrier \(\mathrm{IP}_{10}\) has therefore collapsed from a two-channel contradiction search to a completion-side exclusion problem feeding directly into the strict-upper-gap gate behind \(\mathrm{IP}_9\).
\end{remark}




\noindent\textbf{Reader cue.} Read the next block as the second actual collapse theorem in the new OP~1 contradiction split. Its purpose is to show that completion-side duplicate-signature witnesses are already excluded by the standing retained strict-upper-gap discipline, so that no duplicate-signature witness survives the present OP~1 kernel once both standing exclusion gates are read in the new signature normal form.

\subsection*{19AZ$^{(5d)}$. Completion-side duplicate-signature witnesses collapse under the retained strict-upper-gap discipline}
\label{sec:completion-side-duplicate-signature-witnesses-collapse-strict-upper-gap}
After 19AZ$^{(5c)}$, the live duplicate-signature burden has only one honest place left to survive: as a second admissible completion above the already fixed retained corridor while the retained pinning signature stays unchanged. But this is exactly the sort of residual duplication that the retained strict-upper-gap language was already designed to exclude. Once the retained corridor slot and retained comparison packet are fixed, a genuinely second admissible completion preserving the same retained status surface is no longer a new geometric freedom; it is a forbidden duplicate completion above the same retained corridor.

\begin{proposition}[Completion-side duplicate-signature witnesses are excluded on the retained OP~1 corridor]
\label{prop:completion-side-duplicate-signature-witnesses-excluded-retained-op1-corridor}
Assume the same-corridor reversal certificate of Definition~\ref{def:op1-same-corridor-reversal}. Assume moreover that the retained strict-upper-gap discipline continues to hold on the intended admissible windows in the compressed sense already used by Proposition~\ref{prop:op1-rigidity-gap-uniqueness-bridge}. Then no duplicate-signature witness of completion-side type exists in the sense of Definitions~\ref{def:duplicate-signature-witness-live-op1-kernel} and~\ref{def:occupancy-completion-channels-duplicate-signature-op1}.
\end{proposition}

\begin{proof}
Suppose for contradiction that \((L,L')\) is a duplicate-signature witness of completion-side type. By Definition~\ref{def:duplicate-signature-witness-live-op1-kernel}, the loops \(L\) and \(L'\) are genuinely distinct up to the standing admissible corridor equivalence, yet they preserve the same retained pinning signature. In particular, they preserve the same retained total-channel profile, the same retained optimum location, and the same local comparison/defect packet on the intended admissible windows.

By Definition~\ref{def:occupancy-completion-channels-duplicate-signature-op1}, once the common retained corridor slot and retained comparison packet are fixed, the residual distinction between \(L\) and \(L'\) can only survive as a second admissible completion above that same retained corridor while preserving the same retained status surface. But this is precisely the kind of duplicate admissible completion excluded by the standing strict-upper-gap discipline in the retained OP~1 reading. Hence the assumed witness contradicts the retained strict-upper-gap layer. Therefore no completion-side duplicate-signature witness can exist.
\end{proof}

\begin{corollary}[No duplicate-signature witness survives the retained OP~1 kernel]
\label{cor:no-duplicate-signature-witness-survives-retained-op1-kernel}
Assume the hypotheses of Proposition~\ref{prop:occupancy-side-duplicate-signature-witnesses-excluded-retained-op1-corridor} and Proposition~\ref{prop:completion-side-duplicate-signature-witnesses-excluded-retained-op1-corridor}. Then no duplicate-signature witness for the live OP~1 kernel exists on the intended admissible windows.
\end{corollary}

\begin{proof}
If a duplicate-signature witness existed, Corollary~\ref{cor:live-op1-duplicate-signature-burden-reduces-completion-side-channel} would force it to be of completion-side type. Proposition~\ref{prop:completion-side-duplicate-signature-witnesses-excluded-retained-op1-corridor} eliminates that case. Hence no duplicate-signature witness survives.
\end{proof}

\begin{theorem}[Closure of the live OP~1 kernel under the standing retained exclusion disciplines]
\label{thm:closure-live-op1-kernel-standing-retained-exclusion-disciplines}
Assume the same-corridor reversal certificate of Definition~\ref{def:op1-same-corridor-reversal}. Assume moreover that the retained no-second-carrier discipline and the retained strict-upper-gap discipline both hold on the intended admissible windows in the senses used above. Then the optimum-pinning comparison-rigidity certificate of Definition~\ref{def:op1-optimum-pinning-comparison-rigidity} holds on the live OP~1 kernel. Consequently,
\[
\Delta_{\mathrm{CL}}=0.
\]
Equivalently, under the standing retained exclusion disciplines the normalized OP~1 route
\[
\mathrm{TR}_4 \Longrightarrow \mathrm{IP}_{10} \Longrightarrow \mathrm{IP}_9 \Longrightarrow \mathrm{IP}_8
\]
closes on the intended admissible windows.
\end{theorem}

\begin{proof}
By Corollary~\ref{cor:no-duplicate-signature-witness-survives-retained-op1-kernel}, no duplicate-signature witness exists. Proposition~\ref{prop:live-op1-burden-equivalent-exclusion-duplicate-signature-witnesses} therefore yields the optimum-pinning comparison-rigidity certificate. The final sentence of that same proposition then gives \(\Delta_{\mathrm{CL}}=0\) via the already established retained-window-rigidity / admissible-window-uniqueness bridge.
\end{proof}

\begin{remark}[What the completion-side collapse means for the role of the Monte-Carlo probes]
\label{rem:completion-side-collapse-role-monte-carlo-probes}
After the present step, the Monte-Carlo probes with the fusion/decoder reading should not be read as replacements for the OP~1 theorematic closure. Their real leverage is earlier in the search. They compress the search space by classifying residual behavior into a small number of decoder-relevant recurrence classes and by telling us whether a surviving defect should first be treated as \texttt{ip}-local or as \texttt{bsbl}-translation-sensitive. That is why the work now feels faster: the probes do not prove the retained no-second-carrier or strict-upper-gap exclusions, but they stop us from reopening the wrong parts of the document and from carrying broad rival-loop pictures once the residual signature is already locally pinned.

In the present OP~1 line, this means the probes mainly accelerated the route to the right normal form. They helped isolate the live burden as an injectivity problem, then as a two-channel contradiction search, and finally as the completion-side residue that can be handed directly to the strict-upper-gap gate. So the speed-up is real, but it is a \emph{search-space reduction and routing advantage}, not a substitute for the structural closure theorems themselves.
\end{remark}



\noindent\textbf{Reader cue.} Read the next block as the first honest post-OP~1 normalization of the live OP~5 Step~(B) burden. Its purpose is not to declare the routed photon-frame balance defect solved, but to identify the smallest currently licensed structural kernel on which that burden now lives: the production or exclusion of a carrier-upgraded routed-balance witness at the routed photon-frame cutoff.

\subsection*{19AZ$^{(6a)}$. First carrier-upgraded witness normal form for the live OP~5 Step~(B) kernel}
\label{sec:first-carrier-upgraded-witness-normal-form-live-op5b-kernel}
After the closure of the live OP~1 kernel, the honest next burden is OP~5 Step~(B), i.e. the routed photon-frame balance defect \(\Delta_{\Theta R}(T)\). The current line should not reopen that burden through every packet-wave, breathing-envelope, frame-flicker, or helix picture at once. The earlier carrier-first OP~5 architecture already teaches that visible mismatches come in two grades: either they remain screening-only and therefore subordinate to one still-available common closure-bearing return-side carrier package, or they upgrade and certify loss of that common carrier package itself. The cleanest first OP~5(B) normalization is therefore to ask not for a new global formula immediately, but for the smallest carrier-grade witness through which a still-live routed balance defect could honestly persist.

\begin{definition}[Routed-balance witness for the live OP~5 Step~(B) kernel]
\label{def:routed-balance-witness-live-op5b-kernel}
Fix the routed photon-frame cutoff \(T=\gamma_L^{-3/2}\) and the routed balance defect \(\Delta_{\Theta R}(T)\) from Definition~\ref{def:delta-theta-r-residual}. Under the standing carrier-first OP~5 architecture, an admitted visible witness \(w\) is called a \emph{routed-balance witness for the live OP~5 Step~(B) kernel} if \(w\) is cited to support the claim that the routed photon-frame balance layer remains structurally live at the cutoff \(T\). Such a witness is called
\begin{enumerate}[leftmargin=2em]
  \item \emph{screening-only} if it remains screening-only in the sense of Lemma~\ref{lem:carrier-witness-separation};
  \item \emph{carrier-upgraded} if it is carrier-upgraded in the same sense.
\end{enumerate}
Thus a routed-balance witness is not a new witness species; it is the OP~5 Step~(B) reading of the already licensed witness-grade dichotomy at the routed photon-frame cutoff.
\end{definition}

\begin{proposition}[The live OP~5 Step~(B) kernel is a carrier-upgraded witness problem]
\label{prop:live-op5b-kernel-carrier-upgraded-witness-problem}
Assume Proposition~\ref{prop:carrier-first-op5-interface}, Proposition~\ref{prop:packet-phase-epsilon-screening}, Lemma~\ref{lem:first-diagnostic-witness}, Lemma~\ref{lem:screening-witness-compression}, Lemma~\ref{lem:carrier-witness-separation}, and Corollary~\ref{cor:carrier-escalation-threshold}. Then the honest live OP~5 Step~(B) kernel on the present line is the following alternative:
\begin{enumerate}[leftmargin=2em,label=(\roman*)]
  \item either every routed-balance witness at the routed photon-frame cutoff remains screening-only, in which case no carrier/package-level OP~5 Step~(B) obstruction has yet been exhibited;
  \item or there exists a carrier-upgraded routed-balance witness, in which case the live OP~5 Step~(B) burden has genuinely reached the carrier/package layer.
\end{enumerate}
In particular, the first honest structural reopening of \(\Delta_{\Theta R}(T)\) is not an arbitrary visible mismatch, but the production or exclusion of a carrier-upgraded routed-balance witness.
\end{proposition}

\begin{proof}
By Lemma~\ref{lem:carrier-witness-separation}, every admitted visible witness at the carrier-first OP~5 layer lies in exactly one of two grades: screening-only or carrier-upgraded. Definition~\ref{def:routed-balance-witness-live-op5b-kernel} does nothing more than read that existing dichotomy at the routed photon-frame cutoff \(T\) and relative to the residual \(\Delta_{\Theta R}(T)\). If a cited witness remains screening-only, then by Corollary~\ref{cor:carrier-escalation-threshold} its upstream content is still representable through some common closure-bearing return-side carrier package; hence no carrier/package-level OP~5 Step~(B) obstruction has yet been exhibited. If, on the other hand, a cited witness upgrades to carrier grade, then the same corollary says that no such common carrier package remains available, so the routed balance burden has genuinely reached the carrier/package layer. Since no third currently licensed witness grade exists, the displayed alternative is exhaustive. Therefore the first honest structural OP~5 Step~(B) kernel is exactly the production or exclusion of a carrier-upgraded routed-balance witness.
\end{proof}

\begin{corollary}[First normalized reopening route for the live OP~5 Step~(B) kernel]
\label{cor:first-normalized-reopening-route-live-op5b-kernel}
Under the hypotheses of Proposition~\ref{prop:live-op5b-kernel-carrier-upgraded-witness-problem}, the smallest currently licensed normalized reopening route for a genuinely live OP~5 Step~(B) witness is
\[
\mathrm{TR}_4 \Longrightarrow \mathrm{BS}_1 \Longrightarrow \mathrm{BS}_3 \Longrightarrow \mathrm{BL}_2 \Longrightarrow \mathrm{BL}_4 \Longrightarrow \mathrm{BL}_5,
\]
with the last two steps used only when the routed-balance witness remains live after the masked-single-carrier test and localizes at the readout/kernel stratum.
\end{corollary}

\begin{proof}
The total-channel baseline remains frozen by \(\mathrm{TR}_4\). Because the routed photon-frame balance defect is not a pure local optimum-pinning problem but a return-side witness-grade problem, the first honest move is bridge-side rather than imbalance-side: the route enters through the matrix/phase surface-change gates \(\mathrm{BS}_1\) and \(\mathrm{BS}_3\). Once that surface change has been made, the first honest diagnostic question is whether the putative routed-balance witness is still compatible with one masked single carrier, which is precisely the gate \(\mathrm{BL}_2\). If it collapses there, no deeper reopening is licensed. If it remains live, then \(\mathrm{BL}_4\) is exactly the normalized stratification step that localizes the first remaining obstruction. And when that localized obstruction sits in the readout/kernel stratum, \(\mathrm{BL}_5\) is the recorded compression to absence of one common shell-stable carrier package. Therefore the displayed BS/BL route is exactly the smallest currently licensed reopening route for a genuinely live OP~5 Step~(B) witness.
\end{proof}

\begin{remark}[What this first OP~5 Step~(B) normalization buys]
\label{rem:what-first-op5b-normalization-buys}
This is not yet an OP~5 Step~(B) closure theorem. The routed photon-frame balance defect \(\Delta_{\Theta R}(T)\) is still honestly live. But the burden is now sharper. It is no longer carried as an undifferentiated packet/helix/flicker mismatch surface. It is carried as a witness-grade question: does any routed-balance witness actually upgrade to carrier/package level, or do all currently visible routed witnesses remain screening-only beneath one still-available common return-side carrier package? In the present QAM language, OP~1 was compressed onto the \(\mathrm{IP}\)-route, whereas the live OP~5 Step~(B) burden now announces itself as a \(\mathrm{BS}/\mathrm{BL}\)-route burden from the start.
\end{remark}


\noindent\textbf{Reader cue.} Read the next block as the first negative-side sharpening of the live OP~5 Step~(B) kernel. Its purpose is not yet to close the routed photon-frame balance defect, but to identify the first retained situation in which a routed-balance witness cannot honestly escalate to carrier/package grade at all.

\subsection*{19AZ$^{(6b)}$. Carrier-stable routed returns prevent witness-grade escalation on the live OP~5 Step~(B) kernel}
\label{sec:carrier-stable-routed-returns-prevent-witness-grade-escalation-live-op5b-kernel}
The previous normalization reduced the live OP~5 Step~(B) burden to the production or exclusion of a carrier-upgraded routed-balance witness. The next honest question is whether every routed-balance witness must be tested all the way down to the carrier/package layer, or whether one already has a retained positive branch on which such escalation is structurally forbidden. The earlier routed-return package answers this. Once the routed realization lies on the carrier-stable routed return normal form, every admitted visible routed mismatch is already subordinate to one common closure-bearing return-side carrier. In that situation the witness-grade question collapses on the positive side: a routed-balance witness may still exist as screening-level visible data, but it can no longer count as a carrier-grade OP~5 Step~(B) witness.

\begin{proposition}[Carrier-stable routed returns prevent routed-balance witness escalation]
\label{prop:carrier-stable-routed-returns-prevent-routed-balance-witness-escalation}
Assume Proposition~\ref{prop:live-op5b-kernel-carrier-upgraded-witness-problem}, Corollary~\ref{cor:first-normalized-reopening-route-live-op5b-kernel}, Corollary~\ref{cor:carrier-escalation-threshold}, Theorem~\ref{thm:carrier-stable-diagnostic-collapse}, and Theorem~\ref{thm:first-routed-return-normal-form}. Let $(b,r,s;w)$ be a routed branch-shell realization at the cutoff \(T=\gamma_L^{-3/2}\) together with a routed-balance witness \(w\) in the sense of Definition~\ref{def:routed-balance-witness-live-op5b-kernel}. If $(b,r,s;w)$ lies on the \emph{carrier-stable routed return normal form} of Theorem~\ref{thm:first-routed-return-normal-form}, then \(w\) remains screening-only and cannot be carrier-upgraded. In particular, no carrier-upgraded routed-balance witness exists on the carrier-stable routed-return branch.
\end{proposition}

\begin{proof}
By Theorem~\ref{thm:first-routed-return-normal-form}, the carrier-stable routed return normal form means precisely that there exists one common closure-bearing return-side carrier through which the routed realization is read and that every admitted visible routed mismatch collapses beneath that carrier to the carrier-stable visible normal form. Apply this to the routed-balance witness \(w\). Theorem~\ref{thm:carrier-stable-diagnostic-collapse} then says that, over that same stable carrier package, \(w\) has only one current structural role: it witnesses subordinate screening-loss normal form and does not define an independent carrier-relevant visible species. But Corollary~\ref{cor:carrier-escalation-threshold} identifies carrier upgrade exactly with the first failure of every common closure-bearing return-side carrier-supported realization. Since the carrier-stable routed return normal form explicitly retains such a common carrier package, the threshold for carrier upgrade is not met. Hence \(w\) remains screening-only and cannot be carrier-upgraded. Because \(w\) was arbitrary, no carrier-upgraded routed-balance witness exists on the carrier-stable routed-return branch.
\end{proof}

\begin{corollary}[The live OP~5 Step~(B) burden survives only on the obstructed routed-return branch]
\label{cor:live-op5b-burden-survives-only-on-obstructed-routed-return-branch}
Under the hypotheses of Proposition~\ref{prop:carrier-stable-routed-returns-prevent-routed-balance-witness-escalation}, a genuinely live OP~5 Step~(B) burden can survive only on the \emph{obstructed routed return normal form} of Theorem~\ref{thm:first-routed-return-normal-form}. Equivalently, after the gates \(\mathrm{TR}_4\Rightarrow\mathrm{BS}_1\Rightarrow\mathrm{BS}_3\Rightarrow\mathrm{BL}_2\), any routed-balance witness that still honestly threatens carrier/package loss must already belong to the negative routed-return branch; and when that negative branch localizes at the readout/kernel stratum, \(\mathrm{BL}_5\) compresses the remaining burden to absence of one common shell-stable carrier package.
\end{corollary}

\begin{proof}
By Proposition~\ref{prop:live-op5b-kernel-carrier-upgraded-witness-problem}, the live OP~5 Step~(B) kernel is exactly the production or exclusion of a carrier-upgraded routed-balance witness. Proposition~\ref{prop:carrier-stable-routed-returns-prevent-routed-balance-witness-escalation} excludes such a witness on the carrier-stable routed-return branch. Therefore any genuinely live routed-balance burden must sit on the only remaining routed-return branch, namely the obstructed routed return normal form of Theorem~\ref{thm:first-routed-return-normal-form}. Corollary~\ref{cor:first-normalized-reopening-route-live-op5b-kernel} already records that, after the masked-single-carrier gate \(\mathrm{BL}_2\), the next honest localization step is \(\mathrm{BL}_4\), and that a readout/kernel localization is then compressed by \(\mathrm{BL}_5\) to absence of one common shell-stable carrier package. This is exactly the stated reformulation.
\end{proof}

\begin{remark}[What this second OP~5 Step~(B) sharpening buys]
\label{rem:what-second-op5b-sharpening-buys}
The live OP~5 Step~(B) burden is now smaller than in r05. It is no longer ``any routed-balance witness that might still upgrade.'' It is ``a routed-balance witness on the obstructed routed-return branch that survives the masked-single-carrier test and, if needed, localizes further.'' This is also the point at which the Monte-Carlo/fusion probes materially accelerate the search: once the probes classify a run as carrier-stable on the routed-return side, one no longer needs to pursue carrier-grade OP~5 escalation there. The probes still do not prove the theorematic branch separation by themselves, but they do cut away large amounts of blind BS/BL reopening by steering attention directly to the small class of routed returns that still look obstruction-like after the carrier-stable branch has been excluded.
\end{remark}



\noindent\textbf{Reader cue.} Read the next block as the first localization split of the remaining OP~5 Step~(B) burden on the negative routed-return side. Its purpose is not yet to declare the obstructed routed return harmless, but to separate the part that is already inherited from the upstream OP3/QAM obstruction package from the part that would remain as a genuinely downstream carrier-grade OP~5 obstruction.

\subsection*{19AZ$^{(6c)}$. Obstructed routed returns split into upstream-screened versus downstream carrier-grade OP~5 Step~(B) witnesses}
\label{sec:obstructed-routed-returns-split-upstream-screened-vs-downstream-carrier-grade-op5b}
After 19AZ$^{(6b)}$, the live OP~5 Step~(B) burden survives only on the obstructed routed-return branch. But even this negative branch is still too coarse to be the final witness kernel. The already extracted global obstruction package says that, once the admissible run is inside the bundled OP1$\to$branch$\to$OP3/QAM$\to$OP5 channel, the first licensed failure is already localized either upstream in the OP3/QAM obstruction package or downstream at the carrier-upgraded OP5 obstruction layer. Therefore the obstructed routed-return branch must itself be split: either the routed obstruction is only the return-side shadow of an upstream OP3/QAM obstruction already fixed before carrier-grade OP5 escalation, or it is a genuinely downstream carrier-grade OP5 obstruction that survives that upstream screening.

\begin{definition}[Upstream-screened versus downstream carrier-grade obstructed routed-return witness]
\label{def:upstream-screened-vs-downstream-carrier-grade-obstructed-routed-return-witness}
Assume the hypotheses of Corollary~\ref{cor:live-op5b-burden-survives-only-on-obstructed-routed-return-branch}. Let $(b,r,s;w)$ be a routed branch-shell realization on the \emph{obstructed routed return normal form} together with a routed-balance witness \(w\).
\begin{enumerate}[leftmargin=2em,label=(\roman*)]
\item We say that $(b,r,s;w)$ is an \emph{upstream-screened obstructed routed-return witness} if the currently licensed negative routed-return reading of $(b,r,s;w)$ is already forced by the upstream OP3/QAM obstruction package of Proposition~\ref{prop:global-obstruction-tool}, so that the witness does not yet certify an independent downstream carrier-grade OP~5 obstruction.
\item We say that $(b,r,s;w)$ is a \emph{downstream carrier-grade obstructed routed-return witness} if the witness remains honestly live after the upstream OP3/QAM obstruction screening and therefore certifies a genuinely downstream carrier-upgraded OP~5 obstruction layer in the sense of Proposition~\ref{prop:global-obstruction-tool}.
\end{enumerate}
\end{definition}

\begin{proposition}[The remaining obstructed routed-return burden is an upstream/downstream split]
\label{prop:remaining-obstructed-routed-return-burden-is-upstream-downstream-split}
Assume Proposition~\ref{prop:live-op5b-kernel-carrier-upgraded-witness-problem}, Corollary~\ref{cor:live-op5b-burden-survives-only-on-obstructed-routed-return-branch}, Proposition~\ref{prop:global-obstruction-tool}, and Proposition~\ref{prop:global-transfer-tool}. Let $(b,r,s;w)$ be a genuinely live routed-balance witness on the obstructed routed return normal form. Then exactly one of the following two cases is currently licensed.
\begin{enumerate}[leftmargin=2em,label=(\roman*)]
\item \textbf{Upstream-screened case.} The negative routed-return reading of $(b,r,s;w)$ is already accounted for by the upstream OP3/QAM obstruction package, and the witness is therefore not yet an independent downstream carrier-grade OP~5 obstruction.
\item \textbf{Downstream carrier-grade case.} The witness survives the upstream OP3/QAM obstruction screening and therefore remains as a genuinely downstream carrier-grade OP~5 obstruction witness.
\end{enumerate}
No third autonomous obstructed routed-return witness class is currently licensed.
\end{proposition}

\begin{proof}
By Corollary~\ref{cor:live-op5b-burden-survives-only-on-obstructed-routed-return-branch}, every genuinely live OP~5 Step~(B) witness must already lie on the obstructed routed return normal form. Once the run is inside the bundled OP1$\to$branch$\to$OP3/QAM$\to$OP5 channel, Proposition~\ref{prop:global-obstruction-tool} says that the first licensed failure is already localized either upstream in the OP3/QAM obstruction package or downstream at the carrier-upgraded OP5 obstruction layer. Proposition~\ref{prop:global-transfer-tool} adds that no third autonomous terminal visible class is licensed outside that bundled obstruction reading. Applying this dichotomy to the obstructed routed-return witness $(b,r,s;w)$ yields exactly the two stated cases. The first is the upstream-screened case, where the negative routed-return picture is merely the return-side shadow of an obstruction already fixed upstream. The second is the downstream carrier-grade case, where the witness remains live after that upstream screening and therefore truly belongs to the OP~5 obstruction layer. Exclusivity and exhaustiveness are precisely the content of the bundled global obstruction dichotomy.
\end{proof}

\begin{corollary}[The genuinely new OP~5 Step~(B) burden survives only on the downstream branch]
\label{cor:genuinely-new-op5b-burden-survives-only-on-downstream-branch}
Under the hypotheses of Proposition~\ref{prop:remaining-obstructed-routed-return-burden-is-upstream-downstream-split}, any genuinely \emph{new} OP~5 Step~(B) burden survives only in the downstream carrier-grade case of Definition~\ref{def:upstream-screened-vs-downstream-carrier-grade-obstructed-routed-return-witness}. Equivalently, once the obstructed routed return has been shown not to be merely the return-side shadow of the upstream OP3/QAM obstruction package, the remaining burden is already localized to the downstream carrier-upgraded OP~5 obstruction layer.
\end{corollary}

\begin{proof}
In the upstream-screened case, the negative routed-return witness is already accounted for by the upstream OP3/QAM obstruction package and therefore does not constitute a genuinely new downstream OP~5 burden. Hence only the downstream carrier-grade case can carry a genuinely new OP~5 Step~(B) burden after the upstream screening has been discharged. This is exactly the stated reformulation.
\end{proof}

\begin{remark}[What this third OP~5 Step~(B) sharpening buys]
\label{rem:what-third-op5b-sharpening-buys}
The live OP~5 Step~(B) burden is now smaller again. It is no longer merely ``a witness on the obstructed routed-return branch.'' It is ``a witness on that negative branch that is not already explained by the upstream OP3/QAM obstruction package.'' This is also where the Monte-Carlo/fusion probes become strategically useful on the negative side: they do not prove the upstream/downstream split, but they help tell much earlier whether a run looks like an inherited OP3/QAM obstruction shadow or like a genuinely downstream carrier-grade obstruction. In practical search terms, this means that many apparently live negative routed returns can now be discarded upstream before one spends time reopening the deepest OP~5 carrier-grade layer.
\end{remark}


\noindent\textbf{Reader cue.} Read the next block as the discharge of the upstream-screened side of the remaining OP~5 Step~(B) burden. Its purpose is not yet to close the downstream carrier-grade obstruction kernel, but to remove every negative routed-return witness whose content is already absorbed by the upstream OP3/QAM obstruction package.

\subsection*{19AZ$^{(6d)}$. Upstream-screened obstructed routed returns are discharged before the downstream OP~5 Step~(B) carrier-grade kernel begins}
\label{sec:upstream-screened-obstructed-routed-returns-discharged-before-downstream-op5b-kernel}
After 19AZ$^{(6c)}$, the live negative routed-return burden has only two currently licensed readings: upstream-screened or downstream carrier-grade. But the upstream-screened case is not merely weaker; it is not an independent OP~5 Step~(B) burden at all. By definition it is already the routed return shadow of an obstruction whose first licensed failure sits upstream in the OP3/QAM package. Therefore the honest live OP~5 Step~(B) kernel starts only after that upstream explanation has been discharged.

\begin{proposition}[Upstream-screened obstructed routed-return witnesses do not define an independent OP~5 Step~(B) burden]
\label{prop:upstream-screened-obstructed-routed-return-witnesses-do-not-define-independent-op5b-burden}
Assume Proposition~\ref{prop:remaining-obstructed-routed-return-burden-is-upstream-downstream-split}, Corollary~\ref{cor:genuinely-new-op5b-burden-survives-only-on-downstream-branch}, Proposition~\ref{prop:global-obstruction-tool}, and Proposition~\ref{prop:global-transfer-tool}. Let $(b,r,s;w)$ be an upstream-screened obstructed routed-return witness in the sense of Definition~\ref{def:upstream-screened-vs-downstream-carrier-grade-obstructed-routed-return-witness}. Then $(b,r,s;w)$ does not define an independent live OP~5 Step~(B) burden; its entire negative routed-return content is already discharged by the upstream OP3/QAM obstruction package.
\end{proposition}

\begin{proof}
By Definition~\ref{def:upstream-screened-vs-downstream-carrier-grade-obstructed-routed-return-witness}, an upstream-screened obstructed routed-return witness is precisely one whose currently licensed negative routed-return reading is already forced by the upstream OP3/QAM obstruction package. Proposition~\ref{prop:global-obstruction-tool} says that once the admissible run is inside the bundled OP1$\to$branch$\to$OP3/QAM$\to$OP5 channel, the first licensed failure is already localized either upstream in that OP3/QAM obstruction package or downstream at the carrier-upgraded OP5 obstruction layer. Since the witness is, by hypothesis, upstream-screened, the first licensed failure is already on the upstream side of that dichotomy. Proposition~\ref{prop:global-transfer-tool} excludes any third autonomous terminal visible class outside that bundled obstruction reading. Therefore no additional independent OP~5 Step~(B) burden remains attached to $(b,r,s;w)$ once the upstream package has been acknowledged. In other words, the routed-return shadow carried by $(b,r,s;w)$ is fully discharged upstream and does not itself define a genuinely new OP~5 Step~(B) burden.
\end{proof}

\begin{corollary}[The live OP~5 Step~(B) kernel reduces to the downstream carrier-grade obstructed routed-return branch]
\label{cor:live-op5b-kernel-reduces-to-downstream-carrier-grade-obstructed-routed-return-branch}
Under the hypotheses of Proposition~\ref{prop:upstream-screened-obstructed-routed-return-witnesses-do-not-define-independent-op5b-burden}, the honest live OP~5 Step~(B) kernel is exactly the production or exclusion of a downstream carrier-grade obstructed routed-return witness. Equivalently, once the upstream-screened negative routed-return cases have been discharged, the only remaining live OP~5 Step~(B) burden sits on the downstream carrier-grade branch isolated in Definition~\ref{def:upstream-screened-vs-downstream-carrier-grade-obstructed-routed-return-witness}.
\end{corollary}

\begin{proof}
By Proposition~\ref{prop:remaining-obstructed-routed-return-burden-is-upstream-downstream-split}, every genuinely live obstructed routed-return witness lies in exactly one of two currently licensed classes: upstream-screened or downstream carrier-grade. Proposition~\ref{prop:upstream-screened-obstructed-routed-return-witnesses-do-not-define-independent-op5b-burden} removes the first class as an independent OP~5 Step~(B) burden. Therefore only the downstream carrier-grade class remains available as the honest live OP~5 Step~(B) kernel. This is exactly the stated reformulation.
\end{proof}

\begin{remark}[What this fourth OP~5 Step~(B) sharpening buys]
\label{rem:what-fourth-op5b-sharpening-buys}
The live OP~5 Step~(B) burden is now concentrated on one branch only. It is no longer ``an obstructed routed-return witness that is not yet known to be harmless.'' It is ``a downstream carrier-grade obstructed routed-return witness after the upstream OP3/QAM shadow cases have been discharged.'' This is the first point on the OP~5 side where the negative branch has been reduced to a single honest carrier-grade kernel. Strategically, this is also where the Monte-Carlo/fusion probes become most valuable as a search accelerator: once a negative routed return looks upstream-screened, it can be discarded without reopening the downstream carrier-grade layer, so the expensive search is reserved only for the small class of runs that still look honestly downstream after the upstream package has been factored out.
\end{remark}



\noindent\textbf{Reader cue.} Read the next block as the compression of the remaining downstream OP~5 Step~(B) kernel to one shell-package deficit question on the retained branch-shell corridor. Its purpose is not yet to declare OP~5 Step~(B) globally closed, but to show that once the upstream shadow cases are gone, the honest remaining burden is exactly the absence of one common shell-stable carrier package.

\subsection*{19AZ$^{(6e)}$. Downstream carrier-grade obstructed routed returns compress to one common shell-stable carrier-package deficit on the retained branch-shell corridor}
\label{sec:downstream-carrier-grade-obstructed-routed-returns-compress-shell-stable-carrier-package-deficit}
After 19AZ$^{(6d)}$, the live OP~5 Step~(B) burden is concentrated on the downstream carrier-grade obstructed routed-return branch. The smallest honest next question is therefore no longer whether some vague routed mismatch survives, but whether the retained branch-shell corridor still admits one common shell-stable carrier package once that downstream witness has survived every earlier screening. On the current line this is exactly the readout/kernel compression question behind \(\mathrm{BL}_5\): if the branch and shell strata themselves remain intact on the retained routed corridor, then every remaining downstream carrier-grade witness is nothing but failure of one common shell-stable carrier package.

\begin{definition}[Downstream shell-package deficit witness for the live OP~5 Step~(B) kernel]
\label{def:downstream-shell-package-deficit-witness-live-op5b-kernel}
Fix the live OP~5 Step~(B) setup of Definition~\ref{def:routed-balance-witness-live-op5b-kernel}. A quadruple $(b,r,s;w)$ is called a \emph{downstream shell-package deficit witness for the live OP~5 Step~(B) kernel} if
\begin{enumerate}[leftmargin=2em,label=(\roman*)]
  \item $(b,r,s;w)$ is a downstream carrier-grade obstructed routed-return witness in the sense of Definition~\ref{def:upstream-screened-vs-downstream-carrier-grade-obstructed-routed-return-witness}, and
  \item on the retained routed branch-shell corridor determined after the gates \(\mathrm{TR}_4\Rightarrow\mathrm{BS}_1\Rightarrow\mathrm{BS}_3\Rightarrow\mathrm{BL}_2\), the branch and shell strata remain intact, so the remaining localized burden is entirely in the readout/kernel stratum and therefore equals failure of one common shell-stable carrier package in the sense of Lemma~\ref{lem:qam-readout-kernel-obstruction-compression}.
\end{enumerate}
\end{definition}

\begin{proposition}[Downstream carrier-grade OP~5 witnesses compress to one common shell-stable carrier-package deficit on the retained branch-shell corridor]
\label{prop:downstream-carrier-grade-op5-witnesses-compress-to-one-common-shell-stable-carrier-package-deficit}
Assume Corollary~\ref{cor:live-op5b-kernel-reduces-to-downstream-carrier-grade-obstructed-routed-return-branch}, Corollary~\ref{cor:first-normalized-reopening-route-live-op5b-kernel}, Theorem~\ref{thm:first-routed-return-normal-form}, Lemma~\ref{lem:qam-shell-obstruction-stratification}, and Lemma~\ref{lem:qam-readout-kernel-obstruction-compression}. Let $(b,r,s;w)$ be a genuinely live downstream carrier-grade obstructed routed-return witness for the OP~5 Step~(B) kernel. If, on the retained routed branch-shell corridor carried through \(\mathrm{TR}_4\Rightarrow\mathrm{BS}_1\Rightarrow\mathrm{BS}_3\Rightarrow\mathrm{BL}_2\), the branch and shell strata remain intact, then $(b,r,s;w)$ is exactly a downstream shell-package deficit witness in the sense of Definition~\ref{def:downstream-shell-package-deficit-witness-live-op5b-kernel}. Equivalently, under the retained branch-shell intactness assumptions, the whole remaining downstream OP~5 Step~(B) burden compresses to absence of one common shell-stable carrier package.
\end{proposition}

\begin{proof}
By Corollary~\ref{cor:live-op5b-kernel-reduces-to-downstream-carrier-grade-obstructed-routed-return-branch}, every genuinely live OP~5 Step~(B) burden is already a downstream carrier-grade obstructed routed-return witness. Corollary~\ref{cor:first-normalized-reopening-route-live-op5b-kernel} records that after the retained gates \(\mathrm{TR}_4\Rightarrow\mathrm{BS}_1\Rightarrow\mathrm{BS}_3\Rightarrow\mathrm{BL}_2\), the first further honest localization step is \(\mathrm{BL}_4\), and that whenever the first remaining localized obstruction lies in the readout/kernel stratum, the compression gate \(\mathrm{BL}_5\) rewrites the burden as absence of one common shell-stable carrier package. Now assume that on the retained routed branch-shell corridor the branch and shell strata themselves remain intact. Then Lemma~\ref{lem:qam-shell-obstruction-stratification} forces every remaining first local obstruction into the readout/kernel stratum, and Lemma~\ref{lem:qam-readout-kernel-obstruction-compression} identifies that stratum exactly with failure of one common shell-stable carrier package. Therefore the surviving downstream carrier-grade witness $(b,r,s;w)$ is not a broader residual species: on the retained branch-shell corridor it is precisely a downstream shell-package deficit witness. This is the stated compression.
\end{proof}

\begin{corollary}[The live OP~5 Step~(B) kernel reduces to one common shell-stable carrier-package deficit problem on the retained branch-shell corridor]
\label{cor:live-op5b-kernel-reduces-to-one-common-shell-stable-carrier-package-deficit-problem}
Under the hypotheses of Proposition~\ref{prop:downstream-carrier-grade-op5-witnesses-compress-to-one-common-shell-stable-carrier-package-deficit}, the honest live OP~5 Step~(B) kernel on the retained branch-shell corridor is exactly the production or exclusion of a downstream shell-package deficit witness. Equivalently, once the upstream-screened cases have been discharged and the branch/shell strata remain intact, OP~5 Step~(B) is reduced to the single question whether one common shell-stable carrier package is absent.
\end{corollary}

\begin{proof}
Corollary~\ref{cor:live-op5b-kernel-reduces-to-downstream-carrier-grade-obstructed-routed-return-branch} says that the live OP~5 Step~(B) burden is already concentrated on downstream carrier-grade obstructed routed-return witnesses. Proposition~\ref{prop:downstream-carrier-grade-op5-witnesses-compress-to-one-common-shell-stable-carrier-package-deficit} then shows that, under retained branch-shell intactness, every such witness is exactly a downstream shell-package deficit witness. Hence the remaining kernel is neither a diffuse routed-return problem nor a broader carrier-grade witness search; it is exactly the stated one-package deficit problem.
\end{proof}

\begin{remark}[What this fifth OP~5 Step~(B) sharpening buys]
\label{rem:what-fifth-op5b-sharpening-buys}
This is the first point on the OP~5 side where the surviving kernel is no longer merely a branch classification problem. It has become a concrete deficit question. The live burden is now: not ``some downstream obstruction somewhere on the negative routed-return branch,'' but ``the absence of one common shell-stable carrier package on the retained branch-shell corridor.'' This is also why the Monte-Carlo/fusion probes feel so effective here. Once the probes indicate that the run is both downstream and branch/shell-intact, the search no longer needs to explore broad alternative witness pictures; it can go straight to the small shell-package deficit kernel that \(\mathrm{BL}_5\) was designed to expose.
\end{remark}




oindent\textbf{Reader cue.} Read the next block as the final conservative test of the live OP~5 Step~(B) kernel on the retained branch-shell corridor. Its purpose is not to claim a new global closure principle from scratch, but to show that once the surviving downstream kernel is already compressed to one shell-package deficit question, the inherited OP~3/QAM shell-closure package removes exactly that final deficit under the retained shell-transfer discipline.

\subsection*{19AZ$^{(6f)}$. The inherited shell-closure package discharges the downstream shell-package deficit on the retained branch-shell corridor}
\label{sec:inherited-shell-closure-package-discharges-downstream-shell-package-deficit}
After 19AZ$^{(6e)}$, the live OP~5 Step~(B) kernel has been compressed as far as possible without reusing the earlier OP~3/QAM shell package. The honest remaining question is therefore very narrow: can the retained branch-shell corridor both inherit the shell-closure package and yet still fail to admit one common shell-stable carrier package? On the present line the answer is no. Once the retained branch-shell corridor is genuinely read through the already established shell certificate / shell-stable kernel-readout package, the missing common shell-stable carrier of 19AZ$^{(6e)}$ is exactly the kind of deficit that the inherited package forbids.

\begin{definition}[Retained shell-transfer discipline for the live OP~5 Step~(B) kernel]
\label{def:retained-shell-transfer-discipline-live-op5b-kernel}
Fix the live OP~5 Step~(B) setup of Definition~\ref{def:routed-balance-witness-live-op5b-kernel}. We say that the retained routed branch-shell corridor satisfies the \emph{retained shell-transfer discipline} if, after the retained gates \(\mathrm{TR}_4\Rightarrow\mathrm{BS}_1\Rightarrow\mathrm{BS}_3\Rightarrow\mathrm{BL}_2\), the routed branch-shell realization inherits the shell admissibility certificate and the equivalent shell-stable kernel/readout package from Theorem~\ref{thm:qam-op3-shell-closure}. Equivalently, on that retained corridor there exists one common shell-stable kernel/readout carrier through which all closure-relevant routed readout factors.
\end{definition}

\begin{proposition}[The inherited shell-closure package excludes downstream shell-package deficit witnesses on the retained branch-shell corridor]
\label{prop:inherited-shell-closure-package-excludes-downstream-shell-package-deficit-witnesses}
Assume Corollary~\ref{cor:live-op5b-kernel-reduces-to-one-common-shell-stable-carrier-package-deficit-problem}, Definition~\ref{def:retained-shell-transfer-discipline-live-op5b-kernel}, Theorem~\ref{thm:qam-op3-shell-closure}, and Theorem~\ref{thm:qam-shell-iff-package}. Then no downstream shell-package deficit witness in the sense of Definition~\ref{def:downstream-shell-package-deficit-witness-live-op5b-kernel} can exist on the retained branch-shell corridor. Equivalently, once the retained branch-shell corridor inherits the shell certificate / shell-stable kernel-readout package, the remaining OP~5 Step~(B) kernel from 19AZ$^{(6e)}$ is discharged.
\end{proposition}

\begin{proof}
Assume toward contradiction that a downstream shell-package deficit witness exists on the retained branch-shell corridor. By Definition~\ref{def:downstream-shell-package-deficit-witness-live-op5b-kernel}, this means that the surviving downstream carrier-grade routed-return witness is already compressed to absence of one common shell-stable carrier package on that corridor. But by Definition~\ref{def:retained-shell-transfer-discipline-live-op5b-kernel}, the same retained corridor inherits the shell admissibility certificate and hence, by Theorem~\ref{thm:qam-op3-shell-closure} together with Theorem~\ref{thm:qam-shell-iff-package}, admits one common shell-stable kernel/readout carrier through which the closure-relevant routed readout factors. This is exactly the negation of the shell-package deficit asserted by the witness. The contradiction shows that no downstream shell-package deficit witness can exist once the retained shell-transfer discipline holds. Therefore the remaining OP~5 Step~(B) kernel from Corollary~\ref{cor:live-op5b-kernel-reduces-to-one-common-shell-stable-carrier-package-deficit-problem} is discharged on that corridor.
\end{proof}

\begin{corollary}[Closure of the live OP~5 Step~(B) kernel under the standing retained branch-shell and shell-transfer disciplines]
\label{cor:closure-live-op5b-kernel-under-standing-retained-branch-shell-and-shell-transfer-disciplines}
Assume the standing retained branch-shell intactness hypotheses from Proposition~\ref{prop:downstream-carrier-grade-op5-witnesses-compress-to-one-common-shell-stable-carrier-package-deficit} together with the retained shell-transfer discipline of Definition~\ref{def:retained-shell-transfer-discipline-live-op5b-kernel}. Then the live OP~5 Step~(B) kernel closes:
\[
\Delta_{\Theta R}=0.
\]
Equivalently, under the current retained no-upstream-shadow, branch/shell-intact, and inherited shell-package disciplines, no honest live OP~5 Step~(B) witness survives.
\end{corollary}

\begin{proof}
By Corollary~\ref{cor:live-op5b-kernel-reduces-to-one-common-shell-stable-carrier-package-deficit-problem}, the live OP~5 Step~(B) kernel under the retained branch-shell hypotheses is exactly the production or exclusion of a downstream shell-package deficit witness. Proposition~\ref{prop:inherited-shell-closure-package-excludes-downstream-shell-package-deficit-witnesses} excludes precisely that witness class once the retained shell-transfer discipline holds. Hence no live OP~5 Step~(B) witness remains, so the kernel closes and \(\Delta_{\Theta R}=0\).
\end{proof}

\begin{remark}[What this sixth OP~5 Step~(B) sharpening buys]
\label{rem:what-sixth-op5b-sharpening-buys}
This is the OP~5 analogue of the OP~1 local closure step. The result is intentionally conservative. It does not say that every historical reading of OP~5 is now globally absorbed into one theorematic sentence. It says something narrower and structurally cleaner: once the live OP~5 Step~(B) burden has been reduced to one downstream shell-package deficit on the retained branch-shell corridor, the inherited OP~3/QAM shell-closure package forbids exactly that final deficit. Strategically, this is also the point where the Monte-Carlo/fusion probes show their full value on the negative routed-return side. They do not prove the inherited shell package, but they steer the search quickly onto the tiny downstream/branch-shell-intact corridor where the shell-transfer discipline can be applied and the final witness class disappears.
\end{remark}
\begin{remark}[Historical transition note for the archived v2.27.x release-control fragment]
Read the next local post-release fragment as archival proof-memory from an earlier public hardening cycle. Its purpose is to preserve release genealogy, checkpoint recall, and local baseline memory, not to function as a current roadmap, freeze checklist, or active control surface for the present v\docversion\ line.
\end{remark}

\subsection*{Archived v2.27.3 local post-release note (historical baseline memory)}



% --- Integrated v3.31.0 groundwork-complete pre-OP frontier synthesis ---

